\documentclass[11pt]{article}
\usepackage[a4paper,left=27mm,right=27mm,top=25mm,bottom=25mm]{geometry}
\usepackage[T1]{fontenc}
\usepackage{lmodern,amsmath,amssymb,amsthm,mathtools,mathrsfs}
\usepackage{booktabs,tabularx,array,enumitem,microtype,xcolor,hyperref,needspace,longtable}
\setlist[enumerate]{itemsep=3pt,topsep=5pt,parsep=0pt}
\definecolor{inkblue}{RGB}{25,55,85}
\hypersetup{hidelinks,
 pdftitle={Vertex Operator Algebra Realizations of Twisted Haagerup-Izumi Doubles},
 pdfauthor={David E. Evans and Andrew Riesen},pdfsubject={Regular nonunitary VOA realizations of twisted Haagerup-Izumi centers}}
\pagestyle{plain}
\setlength{\parindent}{1.2em}\setlength{\parskip}{0pt}
\setlength{\emergencystretch}{2em}
\setcounter{tocdepth}{1}
\setcounter{secnumdepth}{2}
\newtheorem{theorem}{Theorem}[section]
\newtheorem{lemma}[theorem]{Lemma}
\newtheorem{proposition}[theorem]{Proposition}
\newtheorem{corollary}[theorem]{Corollary}
\newcommand{\C}{\mathbb C}\newcommand{\Q}{\mathbb Q}\newcommand{\Z}{\mathbb Z}
\newcommand{\one}{\mathbf1}\newcommand{\D}{\mathcal D}\newcommand{\Dcat}{\mathcal D}
\newcommand{\Ncat}{\mathcal N}\newcommand{\Pcat}{\mathcal P}\newcommand{\Fcat}{\mathcal F}
\newcommand{\Scat}{\mathcal S}\newcommand{\U}{\mathcal U}\newcommand{\bo}{\boxtimes}
\newcommand{\FP}{\operatorname{FPdim}}\newcommand{\Irr}{\operatorname{Irr}}
\newcommand{\rank}{\operatorname{rank}}\newcommand{\End}{\operatorname{End}}
\newcommand{\Hom}{\operatorname{Hom}}\newcommand{\id}{\operatorname{id}}
\newcommand{\Vir}{\operatorname{Vir}}\newcommand{\Rep}{\operatorname{Rep}}
\newcommand{\Ext}{\operatorname{Ext}}\newcommand{\wt}{\operatorname{wt}}
\newcommand{\Spec}{\operatorname{Spec}}\newcommand{\soc}{\operatorname{soc}}\newcommand{\ch}{\operatorname{ch}}

\newcommand{\N}{\mathcal N}\newcommand{\F}{\mathcal F}
\newcommand{\zetaM}[1]{\zeta_{#1}}
\numberwithin{equation}{section}
\newcommand{\R}{\mathbb R}
\begin{document}
\title{Vertex Operator Algebra Realizations of\\Twisted Haagerup--Izumi Doubles}
\author{David E. Evans \and Andrew Riesen}
\date{22 October 2026}
\maketitle
\begin{abstract}
We construct vertex operator algebras realizing an infinite family of twisted, nonunitary Haagerup--Izumi doubles. For every odd integer $N\geq3$, we construct a spherical fusion category $\Fcat_N$ with the Haagerup--Izumi fusion rules for $\Z/N\Z$ and pointed associator of exact order $N$, and realize its Drinfeld center as the ribbon category of modules of an explicit vertex operator algebra $\mathscr{V}_N$. The construction begins with a uniform family of strongly rational Virasoro extensions, whose irreducible modules and theta-function characters we determine. A simple-current extension involving a conformally shifted $E_7$ lattice vertex operator algebra removes the additional semion factor. The resulting $\mathscr{V}_N$ is simple, regular, self-contragredient, and $C_2$-cofinite, with central charge $-6N^2-10$, effective central charge eight, and $N^2+3$ irreducible modules with explicit characters. Its weight-zero subspace has dimension $56$, so it is not of CFT type; we establish modularity and the ribbon equivalence directly. The uniform rationality argument combines singular-vector recurrences, a $2$-adic finiteness obstruction, and hypergeometric control of the Zhu algebra.
\end{abstract}

\section{Introduction}
The realization of Haagerup--Izumi doubles by vertex operator algebras is part of the broader problem of realizing modular tensor categories by rational conformal field theories. The first author and Gannon~\cite{B:EG11} introduced families of twisted Haagerup--Izumi modular data and studied their possible VOA realizations. In this paper, we construct vertex operator algebras realizing an infinite family of these twisted nonunitary doubles, with a ribbon equivalence identifying the representation categories.

Recent work has also advanced the realization problem for untwisted doubles. Gannon, Schopieray and Yadav~\cite{B:GSY} construct Haagerup--Izumi fusion categories for every odd cyclic group and determine the modular data of their centers. Lattice constructions for the untwisted Haagerup double at central charge eight appear in work of the second author~\cite{I:Riesen3} and T.-C. Huang~\cite{I:Huang3}. Huang~\cite{I:Huang5} also constructs a unitary strongly rational VOA realizing the positive order-five Haagerup--Izumi double inside $V_{A_4\oplus A_4}$.

Our starting point is the diagonal Virasoro construction of the first author~\cite{A:Evans}. It was motivated by the S-fold theories of Gang and Kim~\cite{I:GK} and the boundary characters proposed by Gang, Kim and Lee~\cite{I:GKL}; its theta-function characters are also related to Flohr's earlier work on two-generator $\mathcal W$-algebras~\cite{I:Flohr}. For odd $n\geq3$, the underlying Virasoro module is
\begin{equation}\label{A:construction}
V=V_{n+2}=\bigoplus_{j\geq0}L(1-6n,nj(j+1)),
\end{equation}
where $L(c,h)$ denotes the simple Virasoro module of central charge $c$ and lowest weight $h$. The multiplication is the canonical diagonal multiplication of \cite[Theorem~3.2]{A:Evans}, recalled in Section~\ref{sec:diagonal}. The first author established strong rationality and identified the twisted center factor at $n=9$, and formulated the uniform extension of these results in \cite[Conjecture~11.1]{A:Evans}.

We establish the structure of this family for every odd $n\geq3$. In particular, we prove strong rationality, classify the irreducible modules and determine their characters. If $J$ denotes the order-two simple current, write
\[
\Scat_n=\langle J\rangle,\qquad
\D_n=\operatorname{Cent}_{\Rep(V_{n+2})}(J),\qquad m=n+4.
\]
There is a ribbon factorization
\begin{equation}\label{E:factorintro}
\Rep(V_{n+2})\simeq\Scat_n\bo\D_n,
\end{equation}
where $\Scat_n$ is a semion category with $d(J)=-1$. A lattice extension of $V_{n+2}$ determines a connected \'etale algebra in $\D_n$, whose module category has near-group fusion rules. For $n=N^2$, a further condensation constructs a spherical fusion category $\Fcat_N$ with Haagerup--Izumi fusion rules and gives a ribbon equivalence $\D_{N^2}\simeq Z(\Fcat_N)$. More generally, $\D_n$ is a Drinfeld center exactly when $n$ is a square.

The remaining step is to realize $\D_{N^2}$ itself by a VOA. We use a simple-current extension to cancel the semion factor in \eqref{E:factorintro}, following the reconstruction method of Gannon and the second author~\cite[Theorem~2.3]{E:GR}. The spherical structure is essential: the auxiliary semion must have the braiding and twist compatible with the negative dimension of $J$. A minuscule conformal shift of the $E_7$ lattice VOA provides the required inverse ribbon factor and integral extension weights. This gives the following main result.

\begin{theorem}[Twisted-center realization]\label{E:main}
Let $N\geq3$ be odd and put $n=N^2$. Let $K=E_7$, let $u\in K^*\setminus K$ be a minuscule weight with $u^2=3/2$, and set
\[
\mathcal W=(V_K,\omega_K+u(-2)\one),\qquad
M=(V_{K+u},\omega_K+u(-2)\one).
\]
In the second expression the conformal vector acts on the lattice module. Then
\begin{equation}\label{E:def}
\mathscr{V}_N=(V_{N^2+2}\otimes \mathcal W)\oplus(J\otimes M)
\end{equation}
admits a simple, regular, $C_2$-cofinite VOA extension structure, unique up to an isomorphism fixing $V_{N^2+2}\otimes \mathcal W$. It is self-contragredient and nonnegatively integer graded, and
\begin{equation}\label{E:maininvariants}
\begin{gathered}
\Rep(\mathscr{V}_N)\simeq Z(\Fcat_N),\qquad |\Irr(\mathscr{V}_N)|=N^2+3,\\
c(\mathscr{V}_N)=-6N^2-10,\qquad \dim(\mathscr{V}_N)_0=56,\\
h_{\min}(\mathscr{V}_N)=-\frac{N^2+3}{4},\qquad c_{\mathrm{eff}}(\mathscr{V}_N)=8.
\end{gathered}
\end{equation}
The equivalence is ribbon. The category $\Fcat_N$ is the category constructed in Section~\ref{A:sec:categorical}, with fusion rules
\[
\alpha_a\alpha_b=\alpha_{a+b},\quad
\rho\alpha_a=\alpha_{-a}\rho,\quad
\rho^2=\one\oplus\bigoplus_{a\in\Z/N}\alpha_a\rho.
\]
Its noninvertible spherical dimension and pointed associator are
\begin{equation}\label{E:HIinvariants}
\begin{gathered}
d(\rho)=\frac{N-\sqrt{N^2+4}}2,\\
\Omega_N(a,b,c)=\exp\!\left(\frac{2\pi i c}{N}
\left\lfloor\frac{a+b}{N}\right\rfloor\right),\qquad 0\leq a,b,c<N.
\end{gathered}
\end{equation}
The class $[\Omega_N]$ has exact order $N$, also for composite $N$.
\end{theorem}

The proof is completed in Section~\ref{E:proof-main} (p.~\pageref{E:proof-main}). The construction and uniqueness of the VOA are proved in Proposition~\ref{E:existence}; self-contragredience and modularity in Propositions~\ref{E:selfdual}--\ref{E:nonCFTmodular}; and the ribbon equivalence and complete module list in Proposition~\ref{E:category}. The category $\Fcat_N$ and its associator are constructed in Sections~\ref{A:cat:condensation}--\ref{A:cat:associator}, and the numerical invariants are calculated in Section~\ref{E:invariants}.

Here \emph{twisted} refers to the associator on the pointed subcategory of $\Fcat_N$. The category $\Fcat_N$ is constructed by condensation, so the equivalence in \eqref{E:maininvariants} identifies the ribbon categories, including their associativity data. Its comparison with the notation $D^1\mathrm{Hg}_N$ of~\cite{B:EG11} is stated at the level of modular data. The VOAs $\mathscr{V}_N$ are nonnegatively graded and regular, with a $56$-dimensional weight-zero space. Sections~\ref{E:auxiliary}--\ref{E:characters} establish their modularity, give an explicit conformal embedding into a deformed rank-eight lattice VOA and determine all their characters.

\subsection*{The diagonal Virasoro family}
The uniform analysis of $V_{n+2}$ also gives results at nonsquare parameters. Its rank-one lattice extension has the form $\mathcal U=V\oplus X$, where $X$ is a second simple module of lowest weight zero. The corresponding algebra $A=\one\oplus X$ lies in $\D_n$. Passing to $A$-modules relates the representation theory of $V$ to near-group categories and gives the following reconstruction theorem.

\begin{theorem}\label{A:cat:main}
For every odd $n\geq3$, the following statements hold.
\begin{enumerate}[label=(\roman*),leftmargin=8mm]
\item The simple $V$-modules are exactly those in Theorem~\ref{B:characterthm}. They are self-contragredient and have one-dimensional lowest spaces. In particular,
\[
 \rank\Rep(V_{n+2})=2n+6,\qquad \rank\D_n=n+3.
\]
Their modular matrices are given by \eqref{B:Sfull}, \eqref{B:Tfull} and \eqref{B:VOAmatrices}. The data of $\D_n$ are Galois conjugate to the cyclic factor $D_n$ of~\cite{A:EG}.
\item Every simple VOA extension structure on $V\oplus X$ is isomorphic to the lattice extension by an isomorphism restricting to the identity on $V$.
\item The algebra $A$ is connected \'etale. The spherical fusion category $\Ncat_n=(\D_n)_A$ has simple objects $g\in\Z/n$ and $R$, with
\begin{equation}\label{A:cat:NG}
 gR=Rg=R,\qquad R^*=R,\qquad
 R^2=\bigoplus_{g\in\Z/n}g\oplus nR.
\end{equation}
There is a braided equivalence
\begin{equation}\label{A:cat:NGcenter}
 Z(\Ncat_n)\simeq\D_n\bo\Pcat_n^{\rm rev},\qquad
 \Pcat_n=\mathcal C(\Z/n,q(x)=e^{2\pi ix^2/n}).
\end{equation}
\item If $n=N^2$, there is a spherical fusion category $\Fcat_N$ with simple objects $\alpha_a,\alpha_a\rho$, $a\in\Z/N$, and fusion rules
\begin{equation}\label{A:cat:HI}
 \alpha_a\alpha_b=\alpha_{a+b},\qquad
 \rho\alpha_a=\alpha_{-a}\rho,\qquad
 \rho^2=\one\oplus\bigoplus_{a\in\Z/N}\alpha_a\rho.
\end{equation}
Its noninvertible spherical dimension is $d(\rho)=(N-\sqrt{N^2+4})/2$. With the tensor convention of Section~\ref{sec:tensorconvention}, its pointed subcategory is $\mathrm{Vec}_{\Z/N}^{\omega_1}$, where
\begin{equation}\label{A:cat:omega}
 \omega_1(a,b,c)=\exp\!\left(\frac{2\pi ic}{N}
                       \left\lfloor\frac{a+b}{N}\right\rfloor\right),
 \qquad 0\leq a,b,c<N.
\end{equation}
The class $[\omega_1]$ has order $N$, and
\begin{equation}\label{A:cat:realization}
 \D_{N^2}\simeq Z(\Fcat_N),\qquad
 \Rep(V_{N^2+2})\simeq\Scat_{N^2}\bo Z(\Fcat_N)
\end{equation}
as ribbon categories. The center data are Galois conjugate to $D^1\mathrm{Hg}_N$ by the automorphism in \eqref{B:CRTHI}. More generally, $\D_n$ is a Drinfeld center if and only if $n$ is a square.
\end{enumerate}
\end{theorem}

The proof of all four parts is assembled in Section~\ref{A:cat:proof-main} (p.~\pageref{A:cat:proof-main}). Part~\textup{(i)} follows from Sections~\ref{sec:spectrum}--\ref{B:sec:characters}; part~\textup{(ii)} is proved in Section~\ref{sec:extension}; and parts~\textup{(iii)}--\textup{(iv)} are proved in Section~\ref{A:sec:categorical}, with the Galois comparison in Section~\ref{B:sec:HIcomparison}. All these arguments use the foundational properties proved in Theorem~\ref{E:foundations}.

\subsection*{Outline of the proof}
The foundational properties of the diagonal family are collected in Theorem~\ref{E:foundations} and proved in Sections~\ref{A:app:c2}--\ref{A:app:rat}. The level-nine singular vector of the primary generator yields recurrences in the $C_2$-algebra. A uniform $2$-adic obstruction proves that these recurrences force finite dimensionality. The deformed lattice construction supplies the extension $V\oplus X$ and the simple current $J$. Relations from the level-four singular vector of $J$ restrict the lowest weights and control extensions between modules; projectivity of the vacuum then gives rationality. Appendices~\ref{H:uniform}--\ref{H:gauge} provide an alternative rationality proof through an explicit hypergeometric description of the Zhu equation.

To complete the module classification, we use the near-group category of $A$-modules. Each simple module not already supplied by lattice branching induces to a multiple $b_iR$ of its unique noninvertible simple object. A bound from the Zhu bimodule of $X$ forces $b_i=1$, determining the rank. Galois symmetry and a theta-character identity recover the remaining modular block, and modular forms of weight $1/2$ determine the characters. At square parameters, condensation then gives the Haagerup--Izumi fusion rules and the pointed associator of exact order $N$.

The construction of $\mathscr{V}_N$ pairs $J$ with the nontrivial module of the shifted $E_7$ lattice VOA. The diagonal simple-current extension removes the semion while preserving the required ribbon structure. The conformal shift contributes $-11$ to the central charge, and the lattice construction determines the remaining invariants and the character formulas.

\subsection*{Organization}
Section~\ref{sec:inputs} recalls the diagonal construction and fixes notation. Sections~\ref{A:app:c2}--\ref{A:app:rat} prove $C_2$-cofiniteness, establish the lattice realization and prove rationality. Section~\ref{A:sec:categorical} constructs the near-group and twisted Haagerup--Izumi categories. Sections~\ref{sec:spectrum}--\ref{B:sec:characters} classify the simple modules and determine the modular data and characters, completing Theorem~\ref{A:cat:main} in Section~\ref{A:cat:proof-main}. Sections~\ref{E:auxiliary}--\ref{E:characters} construct $\mathscr{V}_N$, its lattice parent and its characters, completing Theorem~\ref{E:main} in Section~\ref{E:proof-main}. Section~\ref{E:scope} discusses CFT type and the relation to the untwisted unitary problem. Appendix~\ref{app:null} records the singular-vector coefficients; Appendices~\ref{H:uniform}--\ref{H:gauge} give the alternative hypergeometric proof and its gauge coefficients.

\subsection*{AI disclosure}
Both authors used ChatGPT for proof exploration, symbolic calculations, organization and drafting. The text includes AI-assisted mathematical arguments. Responsibility for the mathematical content and conclusions remains with the authors.

\clearpage
\tableofcontents
\clearpage
\section{The diagonal Virasoro family}\label{sec:inputs}
\subsection{The diagonal construction}\label{sec:diagonal}
Recall the first author's construction \cite[Theorem~3.2]{A:Evans} so that \eqref{A:construction} specifies an algebra rather than only a character. Choose $t>1$ with $t+t^{-1}=n+2$. This number is irrational. At $c=13-6(t+t^{-1})=1-6n$, set
\[
h_{r,s}=\frac{(rt-s)^2-(t-1)^2}{4t}.
\]
The category of $C_1$-cofinite Virasoro modules at this generic central charge is semisimple and rigid, with simple modules $L(c,h_{r,s})$ for $r,s>0$ \cite{A:CJORY}. Its two subcategories generated by the odd $r$ labels $L(c,h_{r,1})$ and the odd $s$ labels $L(c,h_{1,s})$ centralize each other. Indeed,
\[
h_{r,s}-h_{r,1}-h_{1,s}=-\tfrac12(r-1)(s-1)\in2\Z
\]
for odd $r,s$, and $L(c,h_{r,1})\bo L(c,h_{1,s})\cong L(c,h_{r,s})$.

The first subcategory is the generic integer-spin category $SO(3)_q$ with $q^2=e^{2\pi it}$; the second has parameter $q'^2=e^{2\pi i/t}=q^{-2}$. Thus the second is braid-reversed equivalent to the first, taking label $r$ to label $r$. The generic type-$A_1$ braiding and fusion rules give this equivalence, including the twists; on a diagonal label, $h_{r,1}+h_{1,r}\in\Z$. The canonical diagonal algebra in the product of a rigid category with its braid reverse has underlying object
\[
\bigoplus_{r>0\;\mathrm{odd}}L(c,h_{r,1})^*\bo L(c,h_{1,r})
\cong\bigoplus_{j\geq0}L(c,nj(j+1)).
\]
Its multiplication is obtained by pairing dual bases of the fusion spaces; it is nonzero in every allowed channel. The gluing construction and its direct-limit extension theorem identify it with a VOA extension \cite{A:CKM,A:Ind}. It is $\Z_{\geq0}$-graded with finite-dimensional graded pieces and weight-zero space $\C\one$, so the VOA-extension correspondence applies.

This VOA is simple. An ideal is a Virasoro submodule, hence a sum of diagonal simple summands. Multiplying a nonzero summand by its dual reaches the vacuum because the corresponding pairing is nonzero. Thus every nonzero ideal contains the unit. Let $W$ span the top of the $j=1$ summand, of weight $2n$. Fusion with this summand reaches $j+1$ from $j$, so $T,W$ generate the VOA. The singular $WW$ OPE has weights at most $4n-1$, below the next diagonal primary of weight $6n$. It therefore consists of descendants of the vacuum and $W$; closure of these generator OPEs proves strong generation. Finally $V_1=0$, so the invariant-form criterion for a simple CFT-type VOA \cite{A:Li} makes $V$ self-contragredient.

\subsection{Notation}
Write $M(c,h)$ for a Virasoro Verma module and $L(c,h)$ for its simple quotient. Vertex modes use $Y(a,z)=\sum_{m\in\Z}a_mz^{-m-1}$, $L_m=T_{m+1}$, $D=L_{-1}$, and $o(a)=a_{\wt(a)-1}$. In particular
\[
C_2(V)=\operatorname{span}\{a_{-2}b:a,b\in V\},\qquad
C_1(M)=\operatorname{span}\{a_{-1}m:\wt(a)>0\}.
\]
An ordinary module has diagonalizable $L_0$, finite-dimensional weight spaces, and weights bounded below. Generalized modules allow finite Jordan blocks. CFT type means $V_j=0$ for $j<0$ and $V_0=\C\one$. Regularity means that every weak module is a direct sum of ordinary irreducible modules. Strong rationality means simplicity, CFT type, self-contragredience, $C_2$-cofiniteness, and rationality. A ribbon equivalence is a braided tensor equivalence preserving the twist. For a spherical fusion category $\mathcal F$, $Z(\mathcal F)$ denotes its Drinfeld center with the induced ribbon structure.
We reserve $\mathscr{V}_N$ for the final realizing VOA, write $V=V_{n+2}$ for the diagonal Virasoro extension, and use $V_L$ for the lattice VOA associated to an even positive-definite lattice $L$.
The symbols $\FP$ and $d$ denote Frobenius--Perron and spherical dimensions; the latter may be negative. The rank of a fusion category is its number of simple isomorphism classes. The centralizer of $J$ consists of objects whose double braiding with $J$ is the identity. The symbol $\bo$ denotes a VOA fusion product or a Deligne product as its arguments dictate. All categories are over $\C$. The letter $W$ denotes the primary generator of $V$; the shifted $E_7$ lattice VOA used in Sections~\ref{E:auxiliary}--\ref{E:scope} is denoted by $\mathcal W$.

For a commutative algebra $A$ in a braided category, an $A$-module is \emph{local} if its action is unchanged by double braiding with $A$. Write $\D_A$ for all $A$-modules and $\D_A^{\mathrm{loc}}$ for the local ones. An algebra is \emph{connected \'etale} if $\Hom(\one,A)=\C$ and its multiplication has an $A$-bimodule section. Passing to $A$-modules is the condensation used below.

\subsection{Tensor convention}\label{sec:tensorconvention}
For a commutative algebra $A$ in a braided category $\D$, $\D_A$ means right modules with left action on $M$ obtained from $c^{-1}_{M,A}:A\otimes M\to M\otimes A$. This is the convention of \cite[Section 3.3]{A:DMNO}. For an algebra $B$ in a center $Z(\Ncat)$, the quotient $\operatorname{Rep}_{\Ncat}(B)$ uses its specified central half-braiding $B\otimes M\to M\otimes B$ to define the left action. These two conventions agree under the condensation-in-stages identification used below. Using the opposite convention changes the displayed pointed cocycle to its inverse.

\section{\texorpdfstring{Uniform $C_2$-cofiniteness}{Uniform C2-cofiniteness}}\label{A:app:c2}
The generator $W$ has weight $2n$ and belongs to the Virasoro module $L(1-6n,h_{3,3})$. Its level-nine null vector gives a differential equation for the image of $Y(W,z)W$ at a point of the $C_2$-algebra. A nonzero point would produce a polynomial solution; a $2$-adic recurrence rules out such a solution.

\begin{theorem}\label{A:c2:thm:main}
For every positive odd $n$, the diagonal algebra $V_{n+2}$ is $C_2$-cofinite. If $n\geq5$, then its $C_2$-algebra $R_n$, generated by $t=[T]$ and $w=[W]$, satisfies
\[
 t^{2n+4}=t^{n+4}w=0,\qquad \dim R_n\leq3n+8.
\]
Its Zhu algebra is finite-dimensional and commutative, so every irreducible module has a one-dimensional lowest space.
\end{theorem}

The Zhu assertion follows from strong generation by $T,W$, centrality of $[T]$, and the surjection $R_n\twoheadrightarrow\operatorname{gr}A(V)$. The coefficients of the normalized level-nine singular vector and its differential recurrence are recorded in Appendix~\ref{app:null}.

\subsection{The differential equation at a $C_2$-point}
Strong generation gives a surjection $\C[t,w]\twoheadrightarrow R_n$. Since the weight-$4n$ space lies below the next Virasoro highest weight, reduction of $W_{(-1)}W$ gives
\begin{equation}\label{A:c2:eq:quad}
 w^2=A t^{2n}+B t^n w
\end{equation}
for scalars $A,B$. At any point of $\Spec R_n$ with $t=0$, equation \eqref{A:c2:eq:quad} gives $w=0$. Thus a nonzero point must have $t\ne0$. By the grading action it can be normalized to $t=1$.

Suppose such a point exists, and let $\ell:V_{n+2}\to R_n\to\C$ be its evaluation functional. Then
\[
 \ell L_{-2}=\ell,\qquad \ell L_{-1}=0,\qquad
 \ell L_{-r}=0\quad(r\geq3).
\]
Set $F(z)=\ell(Y(W,z)W)$. Every positive power of $z$ has coefficient in $C_2(V)$, and odd-weight classes vanish because $R_n$ is generated in even weights. Hence
\begin{equation}\label{A:c2:eq:finite}
 F(z)\in\C[z^{-2}],\qquad
 F(z)=\nu z^{-4n}+\cdots+w^2,\quad \nu\ne0.
\end{equation}
The nonzero leading coefficient is the nonzero vacuum pairing of the $W$ summand.

For a descendant of the inserted field, the Virasoro Ward identities and the displayed properties of $\ell$ give
\begin{align*}
 D_1&=\partial_z,\\
 D_r&=\delta_{r,2}+(-1)^{r-2}
       \big((r-1)h z^{-r}-z^{1-r}\partial_z\big),\qquad r\geq2.
\end{align*}
These formulas also hold when the inserted vector is already a descendant. They follow by normal ordering: only $L_{-2}$ survives on the left, while $L_0$ and $L_{-1}$ are the only nonzero nonnegative-side modes on the right highest-weight vector.

Let $s_{3,3}=\sum_{\lambda\vdash9}c_\lambda(n)L_{-\lambda}W$ be the level-nine singular vector, normalized by $c_{(1^9)}=1$. Define $P_j(e,n)$ by
\begin{equation}\label{A:c2:eq:defP}
 \sum_{\lambda\vdash9}c_\lambda(n)
 D_{\lambda_1}\cdots D_{\lambda_\ell}z^e
   =\sum_{j=0}^4P_j(e,n)z^{e-9+2j}.
\end{equation}
The coefficients are listed in Appendix~\ref{app:null}. Commuting each $L_m$, $1\leq m\leq9$, to the right gives $L_m s_{3,3}=0$ as a polynomial identity in $n$. Since this singular vector vanishes in $L(c,2n)$,
\begin{equation}\label{A:c2:eq:L}
 \mathcal L_n F=0,\qquad
 \mathcal L_n=\sum_{j=0}^4z^{-9+2j}P_j(z\partial_z,n).
\end{equation}
In particular,
\[
 P_4(e,n)=256en(n+4).
\]
With $x=z^{-2}$ and $\Theta=x\partial_x$, multiplication of \eqref{A:c2:eq:L} by $z$ gives the polynomial-coefficient operator
\begin{equation}\label{A:c2:eq:A}
 \mathcal A_n=\sum_{j=0}^4x^jR_j(\Theta,n),\qquad
 R_j(k,n)=P_{4-j}(-2k,n).
\end{equation}
Thus \eqref{A:c2:eq:finite} would be a nonzero polynomial in the kernel of $\mathcal A_n$.

\subsection{The $2$-adic obstruction}
Write $D=n(n+4)$. The diagonal coefficient is
\[
 R_0(k,n)=-512Dk.
\]
Since $D\ne0$, a nonzero formal power-series solution at $x=0$ must have nonzero constant term and is uniquely determined by that term. Normalize it as
\[
 C(x)=\sum_{a\geq0}C_a x^a,\qquad C_0=1.
\]
The coefficients are rational and satisfy
\begin{equation}\label{A:c2:eq:rec}
 512aD\, C_a=\sum_{j=1}^{\min(4,a)}R_j(a-j,n)C_{a-j}.
\end{equation}

\begin{lemma}[Integer divisibilities]\label{A:c2:lem:div}
For odd $n$ and integral $k$,
\[
 512\mid R_1(k,n),\quad256\mid R_2(k,n),\quad128\mid R_3(k,n).
\]
If $4\mid k$, then $512\mid R_2(k,n)$. Furthermore
\begin{align*}
 R_4(k,n)={}&-32(k+n)(n-k)(2n-k)A(k,n)B(k,n)C(k,n),\\
 A(k,n)={}&4k^2+4k-4n^2-16n+1,\\
 B(k,n)={}&k^2-kn+4k-2n^2-11n+4,\\
 C(k,n)={}&4k^2-12kn+4k+8n^2-10n+1.
\end{align*}
Each of $A,B,C$ is odd for odd $n$ and integral $k$.
\end{lemma}
\begin{proof}
The recurrence polynomials obtained from \eqref{A:c2:eq:defP} have contents $512,256,128,64,32$ for $R_0,\ldots,R_4$, respectively. For odd $n$, their primitive parts give
\[
 \frac{R_1}{256}\equiv k(k+1),\qquad
 \frac{R_2}{128}\equiv k(k+1),\qquad
 \frac{R_3}{64}\equiv k^4+k\pmod2.
\]
All three expressions vanish at every integral $k$. Direct substitution $k=4v$ into $R_2$ gives the additional divisibility by $512$. The factorization of $R_4$ is an exact polynomial identity. The parity assertions follow immediately, including $B\equiv k^2-k-1\equiv1\pmod2$.

These divisibilities can be verified by expansion after $n=2u+1$ and $k=2v+r$ (or $k=4v$): coefficient divisibility in $\Z[u,v]$ gives the assertions on every residue class. This is a finite polynomial identity check, independent of any specialization of $n$.
\end{proof}

Let $K=\Q_2(\pi)$, $\pi^2=2$, and let $\mathcal O=\Z_2[\pi]$ be its valuation ring. Its residue field is $\mathbb F_2$. Define
\[
 b_a=\pi^a a!C_a.
\]
Equation \eqref{A:c2:eq:rec} becomes
\begin{equation}\label{A:c2:eq:scaled}
 b_a=\sum_{j=1}^{\min(4,a)}q_j(a,n)b_{a-j},\qquad
 q_j=\frac{\pi^j(a-1)\cdots(a-j+1)R_j(a-j,n)}{512D}.
\end{equation}
The empty product for $j=1$ is $1$. Since $n$ is odd, $D$ is a unit.

\Needspace{8\baselineskip}
\begin{lemma}\label{A:c2:lem:units}
Every coefficient $q_j(a,n)$ in \eqref{A:c2:eq:scaled} lies in $\mathcal O$. Moreover:
\begin{enumerate}
\item $q_1,q_3\in\pi\mathcal O$ for every applicable $a$;
\item $q_2\in\pi\mathcal O$ unless $4\mid a$;
\item $q_4\equiv1\pmod\pi$ if $4\mid a$, and $q_4\equiv0\pmod\pi$ otherwise.
\end{enumerate}
\end{lemma}
\begin{proof}
The $q_1$ assertion follows from $512\mid R_1$. For $q_3$, use $128\mid R_3$ and the even factor in $(a-1)(a-2)$: the resulting valuation is at least that of $\pi$.

For $q_2=(a-1)R_2(a-2,n)/(256D)$, integrality follows from $256\mid R_2$. If $a$ is odd, $a-1$ is even. If $a\equiv2\pmod4$, then $a-2\equiv0\pmod4$, giving the extra factor of $2$ in Lemma~\ref{A:c2:lem:div}.

For $q_4$, the exact factorization gives the following lower bounds, with equality in the first row. Here $v_2(2)=1$:
\[
\begin{array}{c|c|c|c}
 a\bmod4 & v_2(R_4(a-4,n)) &
 v_2((a-1)(a-2)(a-3)) & v_2(q_4)\\\hline
 0&6&1&0\\
 2&\geq7&\geq2&\geq2\\
 1,3&\geq8&\geq3&\geq4
\end{array}
\]
For instance, in the first row $a-4$ is divisible by $4$, the factors $k+n$ and $n-k$ are odd, and $v_2(2n-k)=1$. A unit has residue $1$ in $\mathbb F_2$. Zero factors only increase the lower bounds.
\end{proof}

\begin{proposition}\label{A:c2:prop:no-poly}
For every positive odd integer $n$, the only polynomial solution of $\mathcal A_nF=0$ over $\C$ is zero.
\end{proposition}
\begin{proof}
By Lemma~\ref{A:c2:lem:units} and induction, all $b_a$ lie in $\mathcal O$. The same induction modulo $\pi$ gives
\begin{equation}\label{A:c2:eq:key}
 \displaystyle b_a\equiv
 \begin{cases}1&4\mid a,\\0&4\nmid a\end{cases}\pmod\pi.
\end{equation}
Indeed, if $a\not\equiv0\pmod4$, all applicable coefficients vanish modulo $\pi$. If $4\mid a$, the only additional possible term is $q_2b_{a-2}$, whose second factor is already zero modulo $\pi$; the fourth-step term propagates $b_{a-4}\equiv1$.

Thus $C_{4m}\ne0$ for every $m\geq0$, so the normalized formal solution is not a polynomial. Any nonzero polynomial solution must have nonzero constant term, because $R_0(d,n)\ne0$ for every positive integer $d$; after normalization it would equal $C(x)$ coefficientwise. This proves the claim, also for complex coefficients.
\end{proof}

\subsection{Consequences for the $C_2$-algebra}
Proposition~\ref{A:c2:prop:no-poly} contradicts the polynomial forced by any nonzero point in the preceding Ward-equation argument. Thus $\Spec R_n$ has only its origin. Since $R_n$ is a finitely generated commutative $\C$-algebra, it is zero-dimensional and finite-dimensional. This proves $C_2$-cofiniteness for every positive odd $n$.

The same operator gives a linear bound on $\dim R_n$.

\begin{proposition}\label{A:c2:prop:bound}
For odd $n>4$, one has $t^{2n+4}=t^{n+4}w=0$ in $R_n$.
\end{proposition}
\begin{proof}
On each of the vacuum and $W$ Virasoro channels, evaluate PBW monomials by sending pure $L_{-2}$ powers of the highest-weight vector to $1$ and every other negative-mode monomial to $0$. This is well defined: the negative-mode character kills $L_{-1}\mathbf1$, and kills the odd-level singular vector $s_{3,3}$ and all its descendants. The associated channel series are
\[
 F_0(z)=z^{-4n}\sum_{a\geq0}\alpha_a z^{2a},\qquad
 F_1(z)=z^{-2n}\sum_{a\geq0}\beta_a z^{2a},\qquad
 \alpha_0=\beta_0=1.
\]
Both satisfy $\mathcal L_nF_i=0$. Their distinct leading powers make them linearly independent.

The $4\times2$ matrix
\[
 M_n=\begin{pmatrix}
 \alpha_{2n+1}&\beta_{n+1}\\
 \alpha_{2n+2}&\beta_{n+2}\\
 \alpha_{2n+3}&\beta_{n+3}\\
 \alpha_{2n+4}&\beta_{n+4}
 \end{pmatrix}
\]
has rank two. Otherwise a nonzero combination $G$ of $F_0,F_1$ has zero coefficients of $z^2,z^4,z^6,z^8$. Its truncation $G_{\leq0}$ is a nonzero polynomial in $z^{-2}$, while $G-G_{\leq0}=O(z^{10})$. Since $\mathcal L_n$ lowers powers by at most nine, $\mathcal L_n(G-G_{\leq0})$ has only positive powers. Conversely $\mathcal L_nG_{\leq0}$ has only negative powers. The equation $\mathcal L_nG=0$ forces both to vanish, contradicting Proposition~\ref{A:c2:prop:no-poly}.

For $j=1,\ldots,4$, the negative product $W_{(-2j-1)}W$ belongs to $C_2(V)$ and has weight $4n+2j<6n$. Its reduction therefore gives
\[
 \nu\alpha_{2n+j}t^{2n+j}
   +\mu\beta_{n+j}t^{n+j}w=0,
 \qquad \nu\mu\ne0,
\]
where $\nu,\mu$ are the two channel normalizations. Multiply the $j$th relation by $t^{4-j}$. Rank two now gives both claimed nilpotence relations.
\end{proof}
Together with \eqref{A:c2:eq:quad}, these relations span $R_n$ by
\[
 1,t,\ldots,t^{2n+3},\qquad w,tw,\ldots,t^{n+3}w.
\]
There are $3n+8$ listed monomials, completing Theorem~\ref{A:c2:thm:main}.


The rationality argument in Section~\ref{A:app:rat} addresses the further semisimplicity assertion.

\Needspace{10\baselineskip}
\section{The lattice realization}\label{A:app:lat}
We construct the conformal lattice extension and the simple current used in the rationality argument. A charge filtration and the ordinary Virasoro block decomposition determine their irreducible summands and characters. The local-module calculation then identifies the pointed category that enters the categorical reconstruction.
Write
\[
 c=1-6n,\qquad V=V_{n+2}=\bigoplus_{j\geq0}L(c,nj(j+1)),\qquad n\geq1.
\]
Let $L=\Z\beta$, $(\beta,\beta)=2n$, $H=\beta_{-1}\one$, $\varepsilon=e^\beta$, and $\lambda\ne0$. On the ordinary lattice vertex algebra put
\begin{equation}\label{A:lat:stress}
 T_\lambda=\frac{1}{4n}H_{-1}^{2}\one+\frac12H_{-2}\one+\lambda D^2\varepsilon,
 \qquad \U=(V_L,T_\lambda).
\end{equation}
The symbols $\mathcal M_a$ denote the $2n$ lattice modules $V_{L+a\beta/(2n)}$.

\begin{theorem}\label{A:lat:main}
For every positive integer $n$, the following hold.
\begin{enumerate}[label=(\roman*),leftmargin=7mm]
\item There is a conformal embedding $V\hookrightarrow\U$. Its image is exactly the sum of the $C_1$-cofinite simple Virasoro submodules of $\U$.
\item As $V$-modules, $\U=V\oplus X$ and $\mathcal M_n=J\oplus Y$, where $X,J,Y$ are simple and
\[
 J\big|_{\Vir}=\bigoplus_{j\geq1}L\left(c,n\left(j^2-\tfrac14\right)\right).
\]
The lowest weights of $X,J,Y$ are $0,3n/4,-n/4$, respectively.
\item For $a\not\equiv0,n\pmod{2n}$, $\mathcal M_a$ is simple upon restriction to $V$, and $\mathcal M_a\cong \mathcal M_{2n-a}$ as $V$-modules. The restrictions with $1\leq a<n$, together with $V,X,J,Y$, give $n+3$ distinct simple $V$-modules.
\end{enumerate}
For positive odd $n$, using the preceding uniform $C_2$ result, $J$ is an order-two simple current in the full category $\mathcal C_n$ of grading-restricted generalized $V$-modules. Moreover,
\begin{equation}\label{A:lat:factor}
 \mathcal C_n\simeq\Scat_n\boxtimes\Dcat_n,
 \qquad \Dcat_n=\operatorname{Cent}_{\mathcal C_n}(J),
\end{equation}
as braided monoidal categories with twist. The semion factor has
\[
 J^2=\one,\qquad F(J,J,J)=-1,\qquad c_{J,J}=i^n,\qquad
 \theta_J=i^{-n},\qquad \dim J=-1.
\]
Semisimplicity is addressed in Section~\ref{A:app:rat}.
\end{theorem}

In addition, the extension algebra $A=V\oplus X$ belongs to $\Dcat_n$, and
\[
 J\boxtimes X\cong Y,\qquad
 (\Dcat_n)^{\mathrm{loc}}_A\simeq
 \mathcal C\left(\Z/n,\ q(a)=e^{2\pi i a^2/n}\right).
\]
For $n=N^2$ with $N$ odd, an explicit further extension in $\Dcat_{N^2}$ has local module category $\mathrm{Vec}$; see Section~\ref{A:lat:condensation}. 

\subsection{The deformed lattice and the Virasoro blocks}
The lattice facts needed in the proof are recalled explicitly. With $T_0$ the first two terms of \eqref{A:lat:stress}, $\varepsilon$ is primary of weight zero and has regular self-OPE. Thus
\[
 L_m^\lambda=L_m^0+\lambda m(m+1)\varepsilon_{m-1}.
\]
The Virasoro relations follow by substitution, with unchanged central charge $c$. Both $L_{-1}$ and $L_0$ are unchanged. A charge-$k$ state with oscillator degree $r$ has weight $nk(k-1)+r$. In particular,
\[
 \U_0=\C\one\oplus\C\varepsilon,\quad
 \U_1=\C H\oplus\C D\varepsilon,\quad
 L_1^\lambda H=-2n(\one+2\lambda\varepsilon).
\]
Since $L_1^\lambda D\varepsilon=0$, the invariant-form functional can be normalized as
\begin{equation}\label{A:lat:ell}
 \ell(\one)=1,\qquad \ell(\varepsilon)=-\frac1{2\lambda}.
\end{equation}
The corresponding invariant form on $\U$ is nondegenerate: its radical is an ideal of the simple lattice vertex algebra. The positive-definite lattice vertex algebra is regular; its weak module decomposition is unchanged, and each lattice coset remains ordinary for the shifted grading. This also gives regularity of $\U$.

For $t+t^{-1}=n+2$, $t$ is irrational. Generic Virasoro Verma modules at a positive Kac label have one proper simple submodule; the $C_1$-cofinite category is semisimple, with positive-label simple objects \cite{A:CJORY}. The Virasoro block decomposition is functorial \cite{A:BNW}. The standard Feigin--Fuchs determinant for the map from a Verma module to a Fock module gives the orientations below \cite{A:FF,A:Fr}.

Set
\[
 h_j=nj(j+1),\quad h_j^+=h_j+(2j+1)^2,
 \quad L_j=L(c,h_j),\quad K_j=L(c,h_j^+),\quad M_j=M(c,h_j).
\]
The two-element block is $\{L_j,K_j\}$, with
\begin{equation}\label{A:lat:Verma}
 0\longrightarrow K_j\longrightarrow M_j\longrightarrow L_j\longrightarrow0.
\end{equation}
Write $\mathsf F_k$ for the charge-$k$ Fock space with the $T_0$ action. Then
\begin{equation}\label{A:lat:Fock}
 \mathsf F_{j+1}\cong M_j,\qquad \mathsf F_{-j}\cong M_j^\vee.
\end{equation}
To check the convention, take $\alpha_+=\sqrt{2t}$ and $\alpha_-=-\sqrt{2/t}$, so $\beta=\alpha_++\alpha_-$. The determinant zeros for a Fock momentum are
$\alpha_{r,s}=((1-r)\alpha_++(1-s)\alpha_-)/2$, with $r,s>0$.
The equality $k\beta=\alpha_{r,s}$ forces $r=s=1-2k$. Hence the Verma-to-Fock map is invertible for $k=j+1>0$. Fock contragredience sends $k$ to $1-k$, proving the other orientation in \eqref{A:lat:Fock}. In particular, $K_j$ is an irreducible Verma module and is not $C_1$-cofinite.

Distinct $j$ give disjoint blocks. The only potential cross-equality would imply
\[
 n(2j+1)^2=(n+4)(2k+1)^2,
\]
which would make $n(n+4)=(n+2)^2-4$ a square. It is not a square for $n\geq1$.


\begin{lemma}[Block projectors on weight-finite unions]\label{A:blocks:union}
Suppose an ordinary lower-bounded Virasoro module $M$ has finite-dimensional weight spaces, and each of its finitely generated submodules has the functorial exact block decomposition in category $\mathcal O$. The block projectors extend compatibly to $M$, preserve its submodules, and remain exact on short exact sequences of such modules.
\end{lemma}
\begin{proof}
A finitely generated submodule is in category $\mathcal O$: the positive modes lower weight, so their orbit on finitely many generators lies in a finite-dimensional sum of weight spaces. If $M_1\subset M_2$ are two such submodules, functoriality makes their block projections agree on $M_1$. Taking their directed union therefore defines a projector on $M$. Every vector has only finitely many block components because it belongs to one finitely generated submodule. A submodule is preserved by applying the same construction to its finitely generated submodules.

For exactness, lift any finitely many vectors from a quotient to the middle module and generate a finitely generated Virasoro submodule there. Exactness of its block decomposition gives the required lifts of block components. This proves exactness of the extended projections. In particular, one may apply a block projector to the charge filtration and its quotients before taking the weightwise finite associated graded. This argument is restricted to ordinary modules; it is not a statement about generalized $L_0$-blocks.
\end{proof}

\begin{lemma}\label{A:lat:blocklemma}
Under the deformed Virasoro action,
\[
 \U=\bigoplus_{j\geq0}B_j,\qquad
 0\longrightarrow M_j\longrightarrow B_j\longrightarrow M_j^\vee\longrightarrow0.
\]
Each $B_j$ has precisely one simple $C_1$-cofinite submodule, isomorphic to $L_j$. Its highest vector $W_j$ has lowest charge $-j$, with nonzero coefficient of $e^{-j\beta}$.
\end{lemma}
\begin{proof}
The filtration $F^r\U=\bigoplus_{k\geq r}\mathsf F_k$ is Virasoro-stable: the perturbation raises charge by one. It is finite and exhaustive on every conformal-weight space. Its graded quotients have the undeformed Fock action.

Apply Lemma~\ref{A:blocks:union} to the charge filtration. In the $j$th block, only the two quotients \eqref{A:lat:Fock} remain. Weightwise finiteness of the filtration therefore gives the displayed exact sequence.

The bottom space of $B_j$ is two-dimensional. Apply the Verma singular operator at level $(2j+1)^2$ to this space. It is nonzero on the bottom of the embedded $M_j$. Its image in $M_j^\vee$ is zero. Its image in $M_j$ consists of singular vectors, a one-dimensional space. Thus the map has rank exactly one. Its one-dimensional kernel generates $L_j$; no other $C_1$-cofinite simple can occur in this block. A kernel vector cannot lie in the embedded Verma module, which proves the assertion about its lowest charge.
\end{proof}

\subsection{The conformal embedding}
Let $S=\bigoplus_{j\geq0}L_j\subset\U$ be the subspace in Lemma~\ref{A:lat:blocklemma}. An intertwining operator between two $C_1$-cofinite modules has $C_1$-cofinite image when the target is lower-bounded \cite[Theorem 4.1.7]{A:CJORY}. Applying this to the lattice products, and using generic semisimplicity, shows that $S$ is closed under all vertex products. It contains $\one,T_\lambda$ and is a CFT-type sub-VOA.

Normalize $W=W_1$ so that its $e^{-\beta}$ coefficient is one. The lowest-charge term of
\[
 W_{-2nj-1}W_j
\]
is a nonzero multiple of $e^{-(j+1)\beta}$. Since its weight is $h_{j+1}$, it has nonzero projection onto $L_{j+1}$. Thus $T_\lambda,W$ generate $S$. The singular $WW$ OPE lies in $L_0\oplus L_1$, because $h_2=6n>4n-1$. This also proves strong generation.

\paragraph{The spin-one recoupling coefficients.}
For generic $q$, put
\[
[a]_q=\frac{q^a-q^{-a}}{q-q^{-1}},\qquad [a]_q!=\prod_{b=1}^a[b]_q,\qquad
\Delta_j=\frac{[2-j]_q![j]_q!^2}{[j+3]_q!}.
\]
For $0\leq j,k\leq2$, the rational Racah matrix is
\[
F_{jk}=[2k+1]_q\Delta_j\Delta_k
\sum_{z=\max(2+j,2+k)}^{\min(4,2+j+k)}
\frac{(-1)^z[z+1]_q!}{[z-2-j]_q!^2[z-2-k]_q!^2[4-z]_q![2+j+k-z]_q!^2}.
\]
Every factorial argument is nonnegative. Rational simplification gives $F^2=I$ and
$F_{0r}=(-1)^r[2r+1]_q/[3]_q\ne0$. In particular a diagonal matrix commuting with $F$ is scalar. These are the exact recoupling facts used next.

\begin{lemma}\label{A:lat:coupling}
The vacuum coefficient and the $W$ self-coupling in $WW$ are both nonzero for $\lambda\ne0$.
\end{lemma}
\begin{proof}
The charge-$-1$ term of $W$ gives a coefficient one of $\one$ in $o(W)\varepsilon$. Since $W$ is primary of even weight, its zero mode is self-adjoint. Consequently \eqref{A:lat:ell} and $o(W)\one=0$ give
\begin{equation}\label{A:lat:zero}
 o(W)\varepsilon=\one+2\lambda\varepsilon.
\end{equation}
The Zhu product $[W]*[W]$ is a sum of descendants in $L_0\oplus L_1$: even its highest contributing weight is $4n<6n$. If $\nu$ and $\mu$ denote the vacuum coefficient and self-coupling, it has the form
\[
 y^2=\nu f_n(x)+\mu g_n(x)y,\qquad f_n(0)=0.
\]
Here $x=[T_\lambda]$, $y=[W]$; $f_n(0)=0$ follows because every contributing vacuum descendant has positive weight. Applying this relation to $\U_0$ gives $\mu g_n(0)=2\lambda$, so $\mu\ne0$.

To determine the vacuum coefficient, use the integer-spin factorization in Section~\ref{sec:diagonal}. Compare multiplication with the canonical diagonal algebra in $\mathcal E\boxtimes\mathcal E^{\mathrm{rev}}$, where $\mathcal E=SO(3)_q$ and $q$ is not a root of unity. Normalize every channel relative to that canonical multiplication, and write its relative scalar as $a_{ij}^k$. Commutativity gives $a_{ij}^k=a_{ji}^k$.

For three spin-one inputs and spin-one output, the two intermediate labels in the product category must coincide. In bases in the second factor dual to the bases in the first, the coefficient matrix is diagonal:
\[
 D=\operatorname{diag}\big(a_{11}^0,(a_{11}^1)^2,a_{11}^2a_{21}^1\big).
\]
The canonical diagonal multiplication is the identity pairing between the two spaces of three-point compositions. Therefore changing their parenthesization acts on these two spaces by $F$ and its inverse transpose. Associativity is exactly $FDF^{-1}=D$. 

In a rational Racah gauge,
\[
 F_{0r}=(-1)^r\frac{[2r+1]_q}{[3]_q},\qquad 0\leq r\leq2.
\]
All three entries are nonzero. Comparing entries $(0,r)$ in $FD=DF$ makes $D$ scalar. Since $a_{11}^1\ne0$, it follows that $a_{11}^0\ne0$, and hence $\nu\ne0$. The displayed Racah coefficients give the required nonzero matrix entries explicitly.
\end{proof}

The restricted invariant form on $S$ has radical equal to its maximal proper ideal. In its simple quotient, $W$ survives by Lemma~\ref{A:lat:coupling}. This quotient is a simple commutative diagonal algebra in the same product category. The support theorem for such algebras \cite[Proposition 4.4]{A:CKM} says that the surviving labels form a tensor subcategory. Spin one generates all integer spins. Hence every $L_j$ survives, the radical is zero, and $S$ is simple.

Finally, the same recoupling calculation gives
\[
 \frac{\nu}{\nu_{\rm can}}=
 \left(\frac{\mu}{\mu_{\rm can}}\right)^2.
\]
After rescaling the canonical generator, the singular generator OPEs of $S$ and $V$ agree: the relevant Virasoro intertwining spaces are one-dimensional. Both are simple, CFT-type, and strongly generated by these generators. The common universal vacuum vertex algebra with these generator relations has a unique maximal proper graded ideal, so their simple quotients agree. This proves $S\cong V$ and the embedding in Theorem~\ref{A:lat:main}.

\subsection{Simplicity of the complement}
The nondegenerate form on $V$ gives a $V$-module orthogonal decomposition
\[
 \U=V\oplus X.
\]
The form on $X$ is nondegenerate. Distinct Virasoro blocks are orthogonal, and
$X=\bigoplus X_j$ with composition factors $L_j,K_j,K_j$. There is no $L_j$ submodule of $X_j$, by the uniqueness in Lemma~\ref{A:lat:blocklemma}.

\begin{lemma}\label{A:lat:essential}
The socle of $X_j$ is one copy of $K_j$, and every nonzero Virasoro submodule of $X_j$ contains it. This copy is the singular submodule of the embedded $M_j$, and has lowest charge $j+1$.
\end{lemma}
\begin{proof}
Only $K_j$ can occur in the socle. If two copies occurred there, self-duality would put two copies in the head as well. The kernel of the quotient onto that head would be a copy of $L_j$, contradicting the absence of an $L_j$ submodule. The socle is therefore a single $K_j$. Every nonzero submodule of a finite-length module contains a simple submodule. The singular copy in $M_j$ already provides this $K_j$ and fixes its lowest charge.
\end{proof}

Write $\phi_j$ for a highest vector of this socle. For $j\geq1$, the lowest-charge part of $Y(W,z)\phi_j$ is
\begin{equation}\label{A:lat:lowering}
 Y(e^{-\beta},z)\phi_j^{(0)}\in\mathsf F_j\{z\},
\end{equation}
where $\phi_j^{(0)}$ is the nonzero Feigin--Fuchs singular vector in $\mathsf F_{j+1}$. This series is nonzero: the exponential lattice field is a product of invertible creation and annihilation exponentials and a charge shift.

Its image in the simple quotient $L_{j-1}$ of $\mathsf F_j\cong M_{j-1}$ is zero. Indeed, a nonzero intertwiner of type $\binom{L_{j-1}}{L_1\ K_j}$ would, by contragredient symmetry, give a nonzero intertwiner of type $\binom{K_j}{L_1\ L_{j-1}}$. The latter is impossible by the $C_1$ image theorem, since $K_j$ is not $C_1$-cofinite. Thus \eqref{A:lat:lowering} has nonzero image in $K_{j-1}\subset\mathsf F_j$.

Exact block projection now shows that some coefficient of $Y(W,z)\phi_j$ has a nonzero component in $X_{j-1}$. By Lemma~\ref{A:lat:essential}, any $V$-submodule containing $\phi_j$ consequently contains $\phi_{j-1}$. Repeating this reaches
\[
 \phi_0=D\varepsilon.
\]
Equation \eqref{A:lat:zero} and translation covariance give
\begin{equation}\label{A:lat:x0}
 W_{2n}D\varepsilon=2n(\one+2\lambda\varepsilon)=2n x_0,
 \qquad x_0:=\one+2\lambda\varepsilon\in X_0.
\end{equation}
The weight-zero vector $x_0$ satisfies
\[
 \langle x_0,x_0\rangle
 =\ell(\one+4\lambda\varepsilon)=-1,
\]
because $\varepsilon_{-1}\varepsilon=0$.

Every nonzero $V$-submodule of $X$ contains some $\phi_j$, by block projection and Lemma~\ref{A:lat:essential}; hence it contains $x_0$. The submodule generated by $x_0$ is therefore simple. Its invariant form is nonzero and consequently nondegenerate. Its orthogonal complement in $X$ would be another submodule and, if nonzero, would also contain $x_0$, an impossibility. Thus $X$ is simple. This argument proves simplicity for all $n$, without assuming rationality.

\subsection{The half-lattice decomposition and the simple current}
Consider $\mathcal M_n=\bigoplus_{k\in\Z}\mathsf F_{k+1/2}$. Put
\[
 u_j=n(j^2-\tfrac14),\quad v_j=u_j+4j^2=(n+4)j^2-\tfrac n4
 \quad(j\geq1).
\]
The argument of Lemma~\ref{A:lat:blocklemma} applies with even Kac label $2j$. The corresponding two Fock quotients have charges $1/2-j$ and $1/2+j$. There is exactly one $C_1$ simple submodule $L(c,u_j)$ in their block. The charge-$1/2$ Fock quotient is a separate one-element block, the irreducible Verma module $L(c,-n/4)$. Thus the $C_1$ socle is precisely
\[
 J=\bigoplus_{j\geq1}L(c,u_j).
\]
It is stable under $V$, again by the $C_1$ image theorem. Its $j$th highest vector has lowest-charge term $e^{(1/2-j)\beta}$, and the lowest-charge lattice product with $W$ gives a nonzero transition from the $j$th to the $(j+1)$st summand.

The lattice coset $\mathcal M_n$ is self-contragredient and has a nondegenerate invariant pairing: the lattice intertwiner $\mathcal M_n\times \mathcal M_n\to\U$, paired with the invariant form on $\U$, identifies $\mathcal M_n$ with its contragredient. Equivalently, this follows from the order-two lattice fusion rule. The next argument shows that the restriction of this pairing to $J$ is nondegenerate.

The lattice intertwiner sends $J\times J$ into $V$, by $C_1$ closure. Let $R$ be the radical of the restricted pairing. It is a $V$-submodule. Invariance and skew-symmetry imply that $Y(R,z)J$ pairs to zero with all of $V$, hence is zero. But take the lowest vector $j_1$ of $J$, whose lowest-charge term is $e^{-\beta/2}$. For any nonzero $r\in R$, the lowest-charge part of $Y(j_1,z)r$ is a nonzero exponential field acting on the lowest-charge part of $r$. It cannot vanish. This contradiction proves nondegeneracy.

Each Virasoro summand of $J$ now has nondegenerate pairing. Any $V$-submodule is a sum of these distinct simple summands and therefore has a $V$-stable orthogonal complement. A nonzero submodule contains all sufficiently large labels by the upward transition above. Two nonzero complementary submodules cannot both do this. Thus $J$ is simple.

Consequently $\mathcal M_n=J\oplus Y$ orthogonally. In each two-element Virasoro block, $Y$ has one $L(c,u_j)$ and two $L(c,v_j)$ factors and has essential socle $L(c,v_j)$, by Lemma~\ref{A:lat:essential}. The same lowest-charge lowering argument sends this socle to the one with label $j-1$. At $j=1$ it reaches the isolated irreducible block $L(c,-n/4)$. Every nonzero submodule of $Y$ therefore contains that block. Its bottom space is the one-dimensional bottom space of $\mathcal M_n$ and has nonzero pairing. The argument used for $X$ proves that $Y$ is simple.

The standard lattice abelian intertwining algebra $\U\oplus \mathcal M_n$ contains the closed subalgebra
\[
 E=V\oplus J.
\]
Closure follows from the $C_1$ image theorem for each pair of summands. Both summands are simple $V$-modules, and the pairing on $J$ makes the vacuum channel in $JJ$ nonzero. Hence $E$ is simple. The involution acting as $+1$ on $V$ and $-1$ on $J$ has fixed-point algebra $V$.

For odd $n$, the preceding $C_2$ theorem supplies the full braided module category of $V$. McRae's abelian-intertwining-algebra orbifold theorem \cite[Corollary 4.8]{A:Orb} applies to $E$ and this involution, without assuming rigidity of that full category. It gives $J\boxtimes_VJ\cong V$. The inherited lattice associator and braiding give $F(J,J,J)=-1$ and $c_{J,J}=i^n$. Also $h_J=3n/4$, so $\theta_J=i^{-n}$, and the ribbon dimension on this rigid rank-two subcategory is $-1$.

\subsection{The remaining lattice modules}
Let $a\not\equiv0,n\pmod{2n}$ and $r=a/(2n)$. The Fock quotients of $\mathcal M_a$ have highest weights
\[
 h(k)=n(k+r)(k+r-1),\qquad k\in\Z.
\]
Each is an irreducible Verma module: a positive Kac label would force $2(k+r)$ to be an integer. They lie in distinct one-element Virasoro blocks. The charge filtration therefore gives a multiplicity-free direct sum of these irreducibles for the deformed Virasoro action. Denote the corresponding submodules by $Q_k$.

The lowest-charge term of $W$ gives a nonzero transition $Q_k\to Q_{k-1}$, just as in \eqref{A:lat:lowering}. The lattice contragredient of $\mathcal M_a$ is $\mathcal M_{2n-a}$; its block paired with $Q_k$ has charge $1-r-k$. The downward transition on that dual module, together with primary-field adjointness, gives a nonzero upward transition $Q_k\to Q_{k+1}$. Every nonzero $V$-submodule consequently contains every $Q_k$. Thus the restriction is simple.

The paired modules have the same lowest conformal weight and the same $W$ zero-mode eigenvalue, since $W$ has even weight. Their one-dimensional top spaces are therefore isomorphic as Zhu-algebra modules: the Zhu algebra is generated by $[T]$ and $[W]$. Uniqueness of the simple module with a specified simple top proves $\mathcal M_a\cong \mathcal M_{2n-a}$ upon restriction.

Put $\Lambda_{b,m}=\theta_{b,m}/\eta$, with
$\theta_{b,m}=\sum_{k\in\Z}q^{(b+2mk)^2/(4m)}$. The modules constructed above have characters
\begin{center}
\begin{tabular}{ccl}
\toprule
Module & Lowest weight & Character\\
\midrule
$V$ & $0$ & $\frac12(\Lambda_{n,n}-\Lambda_{n+4,n+4})$\\[3pt]
$X$ & $0$ & $\frac12(\Lambda_{n,n}+\Lambda_{n+4,n+4})$\\[3pt]
$J$ & $3n/4$ & $\frac12(\Lambda_{0,n}-\Lambda_{0,n+4})$\\[3pt]
$Y$ & $-n/4$ & $\frac12(\Lambda_{0,n}+\Lambda_{0,n+4})$\\[3pt]
$\mathcal M_a$, $1\leq a<n$ & $\frac{(n-a)^2}{4n}-\frac n4$ & $\Lambda_{n-a,n}$\\
\bottomrule
\end{tabular}
\end{center}
These identities follow from the decompositions stated above and the Fock characters. The $n+3$ modules are distinct: all listed lowest weights except the two zeros are distinct, and \eqref{A:lat:zero} separates $V$ and $X$ by the eigenvalues $0$ and $2\lambda$.

\subsection{A uniform polynomial zero-mode identity}
The charge-zero part of the normalized $W$ is a degree-$2n$ Heisenberg polynomial. Let $p_n(s)$ be its zero-mode eigenvalue on a Fock highest vector with $H_0$ eigenvalue $s$. Explicitly, if this part is $\sum_{|p|=2n}c_pH_{-p}\one$, then
\[
 p_n(s)=\sum_{|p|=2n}c_p(-1)^{2n-\ell(p)}s^{\ell(p)}.
\]
Thus $\deg p_n\leq2n$, $p_n(0)=0$, and \eqref{A:lat:zero} gives $p_n(2n)=2\lambda$.

\begin{proposition}\label{A:lat:reflection}
For all $n\geq1$,
\begin{equation}\label{A:lat:polyreflection}
 p_n(s)-p_n(2n-s)=
 \frac{2\lambda}{(2n-1)!}\prod_{a=1}^{2n-1}(s-a).
\end{equation}
\end{proposition}
\begin{proof}
For every integer $1\leq a\leq2n-1$, the bottom of $\mathcal M_a$ is one-dimensional, represented by $e^{a\beta/(2n)}$. The nonzero-charge terms of $W$ cannot preserve its bottom weight, so its zero-mode eigenvalue is $p_n(a)$. Contragredient pairing with $\mathcal M_{2n-a}$ gives $p_n(a)=p_n(2n-a)$; this argument also applies at $a=n$, without any simplicity assumption on its restriction. The left side of \eqref{A:lat:polyreflection} has degree at most $2n-1$, and hence is a scalar multiple of the displayed product. Evaluation at $s=0$ fixes the scalar.
\end{proof}

Equivalently, there is a polynomial $Q_n$ of degree at most $n$, with $Q_n(0)=\lambda$, such that
\[
 p_n(s)=Q_n\left(\frac{s(s-2n)}{4n}\right)
 +\frac{\lambda}{(2n-1)!}\prod_{a=1}^{2n-1}(s-a).
\]
This is a polynomial identity for every $n$. It constrains the first-radius branch of the Zhu spectrum; it is not, by itself, a relation proving the whole Zhu algebra reduced.

\subsection{The semion factor}
Here is the elementary categorical point needed for \eqref{A:lat:factor}. In any abelian braided monoidal category containing an invertible $J$ with $J^2\cong\one$ and $c_{J,J}^2=-\mathrm{id}$, transport the monodromy of $J$ with an object $M$ across the equivalence $J\otimes-$. It defines a natural endomorphism $\kappa_M$ of $M$. The hexagon and $J^2\cong\one$ give $\kappa_M^2=\mathrm{id}$. Thus
\[
 M=M_+\oplus M_-,\qquad \kappa|_{M_\pm}=\pm\mathrm{id}.
\]
The plus subcategory is the centralizer of $J$. Tensoring with $J$ exchanges the signs. Therefore every object is a direct sum of an object of $\Dcat_n$ and $J\otimes D$ with $D\in\Dcat_n$. Morphisms between the two signs vanish, while tensoring with $J$ is fully faithful on either sign. The tensor-product functor
\[
 \Scat_n\boxtimes\Dcat_n\longrightarrow\mathcal C_n
\]
is consequently fully faithful and essentially surjective. Its usual tensor structure is braided because the two factors centralize each other; balancing gives compatibility with twists. This proves \eqref{A:lat:factor}. Full rigidity has not been used.

\subsection{Local modules and the square case}\label{A:lat:condensation}
Regard $A=\U=V\oplus X$ as a commutative algebra object of $\mathcal C_n$. Its unit is simple as an $A$-module. The induction functor $F_A$ is strong monoidal into the category of all $A$-modules. Hence $F_A(J)$ is invertible and in particular simple. The inclusion $J\hookrightarrow \mathcal M_n$ induces a nonzero $A$-module map
\[
 F_A(J)\longrightarrow \mathcal M_n.
\]
Both objects are simple $A$-modules, so it is an isomorphism. Comparing underlying $V$-modules gives
\[
 A\boxtimes J\cong J\oplus Y,\qquad X\boxtimes J\cong Y.
\]
Moreover, $F_A(J)$ is local. Locality of a free module implies that its inducing object centralizes $A$: insert the algebra unit into the local-action identity to recover the monodromy of $A$ with that object. Thus $J$ centralizes $A$, and $A$ lies in $\Dcat_n$.

The equivalence \eqref{A:lat:factor} identifies local $A$-modules in $\mathcal C_n$ with $\Scat_n\boxtimes(\Dcat_n)_A^{\mathrm{loc}}$. The former category is the lattice module category of $\U$. Its objects are $\mathcal M_a$, with fusion addition modulo $2n$, braiding inherited from the lattice, and twist
\[
 \theta_a=\exp\left(2\pi i\left[\frac{a^2}{4n}-\frac a2\right]\right).
\]
The image of $J$ is $\mathcal M_n$. Its centralizer consists exactly of the even labels $a=2x$. Therefore
\begin{equation}\label{A:lat:pointed}
 (\Dcat_n)_A^{\mathrm{loc}}\simeq\mathcal P_n
 :=\mathcal C\left(\Z/n,\ q(x)=e^{2\pi i x^2/n}\right).
\end{equation}
In a convenient gauge $\mathcal P_n$ has trivial associator, braiding $c_{x,y}=e^{2\pi ixy/n}$, and dimensions one.

Now let $n=N^2$ with $N$ odd and put $\alpha=\beta/N$. The subgroup $N\Z/N^2$ is Lagrangian in $\mathcal P_{N^2}$. Its regular algebra is realized by the concrete overlattice extension
\[
 B_N=(V_{\Z\alpha},T_\lambda),\qquad (\alpha,\alpha)=2.
\]
As a $\U$-module it is $\bigoplus_{r=0}^{N-1}\mathcal M_{2Nr}$. As a $V$-module, it is the explicitly semisimple object
\begin{equation}\label{A:lat:BN}
 B_N\cong V\oplus X\oplus
 2\bigoplus_{r=1}^{(N-1)/2}Z_{Nr},
 \qquad Z_{Nr}:=\mathcal M_{2Nr}\big|_V,\quad h_{Z_{Nr}}=r(r-N).
\end{equation}
Here pairing $r$ with $N-r$ uses the proved isomorphism of generic lattice restrictions. All these summands belong to $\Dcat_{N^2}$: the regular algebra of the Lagrangian subgroup is an object of $(\Dcat_{N^2})_A^{\mathrm{loc}}$ by \eqref{A:lat:pointed}; forgetting its $A$-module structure gives $B_N$ in $\Dcat_{N^2}$, which is closed under direct summands.

The category of local $B_N$-modules in $\Dcat_{N^2}$ is $\mathrm{Vec}$, since further local induction corresponds to local modules of that Lagrangian algebra in $\mathcal P_{N^2}$. This also follows directly from the two lattice modules of $V_{\Z\alpha}$: the surviving nontrivial one is the image of $J$, so the centralizer has only the unit. The grading of $B_N$ is bounded below but can have negative weights; no CFT-type claim is made for this overlattice extension.

Thus \eqref{A:lat:BN} supplies the concrete algebra for the center reconstruction in Section~\ref{A:sec:categorical}. After rationality is established in Section~\ref{A:app:rat}, Lemma~\ref{A:cat:etale} and condensation in stages prove that its full module category is the spherical fusion category $\Fcat_N$ and that $Z(\Fcat_N)\simeq\D_{N^2}$.


\Needspace{10\baselineskip}
\section{Uniform rationality}\label{A:app:rat}
We now prove rationality using the $C_2$-cofiniteness theorem of Section~\ref{A:app:c2} and the lattice and simple-current constructions of Section~\ref{A:app:lat}. The argument first controls the ordered pair of lowest weights under the simple current, then excludes extensions of the vacuum. A Wronskian zero count removes the final weight-one obstruction. The section concludes by collecting these results in Theorem~\ref{E:foundations}. 
Let $n\geq3$ be odd and put
\[
 V=V_{n+2},\qquad c=1-6n,\qquad h_W=2n,\qquad h_J=\frac{3n}{4}.
\]
Sections~\ref{A:app:c2} and \ref{A:app:lat} establish the following properties: $V$ is simple, self-contragredient, of CFT type and $C_2$-cofinite; it is strongly generated by $T$ and the even-weight primary $W$; and its full grading-restricted generalized module category $\mathcal C$ contains the simple current
\begin{equation}\label{A:rat:Jdecomp}
 J\bo J\cong V,\qquad
 J\big|_{\Vir}=\bigoplus_{r\geq1}L(c,u_r),\qquad
 u_r=n\left(r^2-\frac14\right).
\end{equation}
Write $j_r$ for a highest vector of the $r$th summand and $j=j_1$. The lattice construction gives a nonzero upward component
\begin{equation}\label{A:rat:up}
 W_{-n(2r-1)-1}j_r=\gamma_rj_{r+1}
       +\text{descendants in summands with label at most }r,
 \qquad \gamma_r\ne0.
\end{equation}
It also supplies a simple $Y$ of lowest weight $-n/4$ and a simple $X$ of weight zero with $o(W)|_{X_0}\ne0$, with $J\bo X=Y$. All these properties hold for every odd $n\geq3$.

\begin{theorem}\label{A:rat:main}
For every odd $n\geq3$, the following assertions hold.
\begin{enumerate}[label=(\roman*),leftmargin=7mm]
\item If $S$ is simple and $S^J=J\bo S$, with lowest weights $h,h^J$, then for some $a\in\Z$ either
\begin{equation}\label{A:rat:pair1}
 (h,h^J)=\left(\frac{a^2}{4n}-\frac n4,
                    \frac{(a+n)^2}{4n}-\frac n4\right),
\end{equation}
or
\begin{equation}\label{A:rat:pair2}
 (h,h^J)=\left(\frac{a^2}{4(n+4)}-\frac n4,
                    \frac{(a+n+4)^2}{4(n+4)}-\frac n4\right).
\end{equation}
The label $a$ is odd in the centralizer of $J$ and even in the other monodromy sector. There is no cutoff on $a$ in this theorem.
\item All simple lowest weights satisfy $h\geq-n/4$. Since $Y$ attains that bound, $c_{\mathrm{eff}}=c-24h_{\min}=1$.
\item The ordered pair $(h,h^J)$ determines $S$ up to isomorphism. In particular, at most two simple modules have any prescribed lowest weight.
\item For simple modules $S,T$ with $h_S=h_T$,
\begin{equation}\label{A:rat:extvanish}
 \Ext^1_{\mathcal C}(S,T)=0.
\end{equation}
In particular, every simple module has no non-split self-extension, including logarithmic self-extensions. The same vanishing holds when $h_{S^J}=h_{T^J}$.
\end{enumerate}
\end{theorem}

\subsection{The simple-current Zhu bimodule is cyclic}
Let $A(V)$ be Zhu's algebra and $A(J)=J/O_V(J)$ the Frenkel--Zhu bimodule. Write $x_l,x_r$ for left and right multiplication by $[T]$, respectively. Since $A(V)$ is generated by $x=[T]$ and $y=[W]$, and $x$ is central, $A(V)$ is commutative. Its simple modules, and therefore the lowest spaces of simple $V$-modules, are one-dimensional.

The standard bimodule identities give, for homogeneous $v\in J$,
\begin{equation}\label{A:rat:VirZhu}
 [L_{-m}v]=(-1)^m\big(mx_r-x_l+\wt v\big)[v],\qquad m\geq1.
\end{equation}
For example, $[L_{-1}v]=(x_l-x_r-\wt v)[v]$ and
$[L_{-2}v]=(2x_r-x_l+\wt v)[v]$. Thus the image of each Virasoro summand is generated by $[j_r]$ over $\C[x_l,x_r]$.

\begin{lemma}\label{A:rat:cyclic}
\[
 A(J)=\C[x_l,x_r][j].
\]
Moreover, for polynomials $F_n,G_n$,
\begin{equation}\label{A:rat:WG}
 [W]*[j]=F_n(x_l,x_r)[j],\qquad
 [j]*[W]=G_n(x_l,x_r)[j].
\end{equation}
\end{lemma}
\begin{proof}
The higher circle relations
\[
 \operatorname{Res}_z\frac{(1+z)^{2n}}{z^k}Y(W,z)v\in O_V(J),\qquad k\geq2,
\]
reduce $[W_{-k}v]$, $k\geq2$, to a linear combination of classes $[W_mv]$ with $m\geq-1$. This follows by induction on $k$, since the leading term of the displayed relation is $W_{-k}v$.

Apply this to \eqref{A:rat:up}. The reduced terms have weights at most $u_r+2n$, whereas
\[
 u_{r+1}-u_r=n(2r+1)>2n.
\]
Hence all reduced terms lie in Virasoro summands with label at most $r$. The other terms in \eqref{A:rat:up} have classes polynomial in $x_l,x_r$ times $[j_s]$, $s\leq r$. Since $\gamma_r\ne0$, induction expresses every $[j_{r+1}]$ in terms of $[j]$. This proves cyclicity.

Every term in $W*j$ or $j*W$ has weight at most $u_1+2n<u_2$. It therefore belongs to the first Virasoro summand. Applying \eqref{A:rat:VirZhu} gives \eqref{A:rat:WG}.
\end{proof}

For simple $S,S^J$, the top-space map associated to the nonzero intertwiner of type $\binom{S^J}{J\ S}$ is nonzero by the Frenkel--Zhu--Li injectivity theorem; only injectivity, not the generally false unrestricted surjectivity statement, is used \cite[Theorem 2.12]{A:AA}. Lemma~\ref{A:rat:cyclic} then implies that its value on $[j]$ is nonzero. Thus
\begin{equation}\label{A:rat:nonzero}
 \mathcal Y(j,z)s=z^{h^J-h-h_J}(t+\text{higher powers of }z),
 \qquad 0\ne t\in S^J_{h^J},\quad 0\ne s\in S_h.
\end{equation}


\subsection{The level-four polynomial and the spectrum}
The $(2,2)$ null vector in the first summand of $J$ is
\begin{equation}\label{A:rat:null}
 \bigl(L_{-1}^4-2(n+2)L_{-2}L_{-1}^2+n(n+4)L_{-2}^2
       -2(n-1)L_{-3}L_{-1}-3nL_{-4}\bigr)j=0.
\end{equation}
Applying \eqref{A:rat:VirZhu}, or the three-point Ward identities, gives the exact bimodule relation
\begin{equation}\label{A:rat:bimodP}
 P_n(x_r,x_l-x_r)[j]=0,
\end{equation}
where
\begin{equation}\label{A:rat:P}
 P_n(h,d)=
 \left[\left(d-\frac n4\right)^2-n\left(h+\frac n4\right)\right]
 \left[\left(d-\frac{n+4}{4}\right)^2-(n+4)\left(h+\frac n4\right)\right].
\end{equation}
Both identities are obtained by applying the positive Virasoro modes to the displayed level-four vector and then substituting \eqref{A:rat:VirZhu}.

For $d=h^J-h$, \eqref{A:rat:nonzero} and \eqref{A:rat:bimodP} imply $P_n(h,d)=0$. Finite order of $J$ implies that its monodromy with a simple module is $\pm1$. Balancing therefore gives
\begin{equation}\label{A:rat:half}
 e^{2\pi i(d-3n/4)}=\pm1,\qquad
 d-\frac{3n}{4}\in\frac12\Z.
\end{equation}
This uses the categorical finite-order monodromy fact, which does not require rigidity of the ambient category \cite[Lemma 2.13]{A:CKL}.

If the first factor of \eqref{A:rat:P} vanishes, put $a=2d-n/2$. Then $a\in\Z$ by \eqref{A:rat:half}, and direct substitution gives \eqref{A:rat:pair1}. For the second factor put $a=2d-(n+4)/2$; this gives \eqref{A:rat:pair2}. In the plus monodromy sector, $d-3n/4\in\Z$, so both possible integer labels are odd. In the minus sector they are even. The lower bound and effective central charge follow immediately.

To prove uniqueness for an ordered pair, evaluate the second relation in \eqref{A:rat:WG} on the nonzero map \eqref{A:rat:nonzero}. The $W$ eigenvalue of $S_h$ is
\begin{equation}\label{A:rat:weightW}
 w_S=G_n(h^J,h).
\end{equation}
Thus the ordered pair determines its one-dimensional $A(V)$-module, and hence its simple $V$-module. For fixed $h\ne-n/4$, only one quadratic lattice can occur: simultaneous occurrence with nonzero integer labels would make $n(n+4)$ a rational square. That number lies strictly between $(n+1)^2$ and $(n+2)^2$. The remaining branch gives at most two values of $h^J$. At $h=-n/4$, the two possible values are $0$ and $1$. This proves the two-module bound.


\begin{lemma}[Lowest-space residues for a log-free intertwiner]\label{A:lowest:residue}
Let $J$ be an ordinary module, and let $M,N$ be grading-restricted generalized modules with lowest generalized weight spaces $M_{[h]},N_{[h']}$. Suppose a log-free intertwiner $\mathcal Y$ of type $\binom N{J\ M}$ commutes with the nilpotent parts of $L_0$. For homogeneous $v\in J$, let $o_{\mathcal Y}(v)$ be the coefficient of $z^{h'-h-\mathrm{wt}(v)}$ on these spaces. Then
\[
 [v]\longmapsto o_{\mathcal Y}(v)
\]
annihilates $O_V(J)$ and respects the left and right Zhu actions, with $[\omega]$ acting by the actual operators $L_0|_{N_{[h']}}$ and $L_0|_{M_{[h]}}$, respectively. No injectivity assertion is made for nonsimple $M,N$.
\end{lemma}
\begin{proof}
The Zhu top-space residue identities use the Jacobi identity, the derivative identity on $J$, and the fact that every positive-degree lowering mode annihilates a lowest space. These statements hold for generalized lowest spaces too. Expand the factors $(1+x)^{\mathrm{wt}(a)}$ for homogeneous $a\in V$ in the usual finite binomial sums. Taking the specified lowest-to-lowest coefficient in the Jacobi identity gives
\[
 o_{\mathcal Y}(a*v)=o(a)\,o_{\mathcal Y}(v),\qquad
 o_{\mathcal Y}(v*a)=o_{\mathcal Y}(v)\,o(a),\qquad
 o_{\mathcal Y}(a\circ v)=0.
\]
The same residue calculation gives the higher-circle identities. Log-freeness ensures that coefficient extraction introduces no derivatives of $\log z$. Compatibility with nilpotent parts ensures that $L_0$ is retained as an operator in these equations, rather than replaced prematurely by its eigenvalue. This proves the stated descent. In the application below, invertibility of the particular top map comes from naturality with the simple submodule and quotient; it does not come from extending the injectivity theorem for simple modules.
\end{proof}

\subsection{Vanishing of self-extensions}
The polynomial relation can now be applied to generalized lowest spaces.

Suppose
\begin{equation}\label{A:rat:E}
 0\longrightarrow S\longrightarrow E\longrightarrow S\longrightarrow0.
\end{equation}
It is enough to consider indecomposable $E$. Tensoring by $J$ is an equivalence and hence exact, so $E^J=J\bo E$ is a self-extension of $S^J$. Their generalized lowest spaces have dimension two. Write $N_E=L_0-(L_0)_{\rm ss}$ and similarly $N_{E^J}$. Both are $V$-module endomorphisms and square to zero.

Naturality shows that the monodromy on $J\bo E$ induces the same scalar on its simple submodule and quotient. Its difference from that scalar therefore squares to zero. Since $J$ has finite order, Lemma~2.13 of~\cite{A:CKL} makes this nilpotent part vanish, so the monodromy is scalar. Taking nilpotent logarithms in the balancing identity gives
\begin{equation}\label{A:rat:Ntransport}
 N_{E^J}=\id_J\bo N_E.
\end{equation}
Naturality of the universal intertwiner then says
$N_{E^J}\mathcal Y_E(j,z)=\mathcal Y_E(j,z)N_E$.
The usual $L_0$-bracket relation shows that the coefficients of positive powers of $\log z$ vanish: on each fixed pair of generalized weight spaces the remaining differential equation is the ordinary homogeneous Euler equation. Thus $\mathcal Y_E$ is log-free.

Let $h,h^J$ be the simple lowest weights and $d=h^J-h$. The coefficient of $z^{d-h_J}$ gives a map
\[
 T_E:E_{[h]}\longrightarrow E^J_{[h^J]}.
\]
Naturality with respect to the submodule and quotient in \eqref{A:rat:E} gives a morphism between the two short exact sequences of lowest spaces. Its two outer maps are the nonzero one-dimensional maps \eqref{A:rat:nonzero}. Consequently $T_E$ is an isomorphism. In particular,
\[
 N_{E^J}T_E=T_EN_E.
\]

Apply Lemma~\ref{A:lowest:residue} to this log-free coefficient map on the generalized lowest spaces. Relation \eqref{A:rat:bimodP}, with the common nilpotent part cancelled in $x_l-x_r$, gives
\begin{equation}\label{A:rat:nilpoly}
 T_E P_n(h+N_E,d)=0.
\end{equation}
Both spaces are lowest spaces of self-extensions, so the coefficient map has no contribution from lower weights.

The crucial derivative is nonzero. Indeed,
\[
 \operatorname{disc}_{h}P_n(h,d)
       =\frac{(16d^2-n(n+4))^2}{16}.
\]
Since $d$ is rational and $n(n+4)$ is not a square, the two roots in $h$ are distinct. More explicitly, at the roots supplied by the first and second factors,
\begin{equation}\label{A:rat:deriv}
 \partial_hP_n=\phantom{-}\frac{16d^2-n(n+4)}4,
 \qquad
 \partial_hP_n=-\frac{16d^2-n(n+4)}4,
\end{equation}
respectively. Expanding \eqref{A:rat:nilpoly} using $N_E^2=0$, $P_n(h,d)=0$ and invertibility of $T_E$ gives $N_E=0$ on $E_{[h]}$.

The module $E$ is generated by that lowest space: it contains the lowest space of the submodule $S$ and maps onto the lowest space of the quotient $S$. Hence $N_E=0$ on all of $E$. The same holds for $E^J$.

Now evaluate \eqref{A:rat:WG} on $T_E$. Since $x_l=h^J$ and $x_r=h$ are scalar,
\[
 T_E\,o(W)=G_n(h^J,h)T_E=w_ST_E.
\]
Thus $o(W)$ is also scalar on $E_h$. All of $A(V)$ acts there as the character of $S_h$.

Choose a lift $s\in E_h$ of the quotient's lowest vector and let $F=V\cdot s$. Its lowest space is $A(V)s=\C s$. Therefore $F$ cannot contain the submodule $S$, whose lowest vector is linearly independent of $s$. Since $S$ is simple, $F\cap S=0$. The image of $F$ in the quotient is all of $S$, so $F\cong S$ and \eqref{A:rat:E} splits. This proves $\Ext^1(S,S)=0$.

If $S,T$ are distinct simples with the same lowest weight, their one-dimensional Zhu characters are distinct. A finite-dimensional commutative algebra has no extensions between distinct characters. Hence the two-dimensional lowest space of an extension of $S$ by $T$ splits as an $A(V)$-module. A lift in the character of $S$ generates a complement by the same intersection argument. This proves \eqref{A:rat:extvanish}. Applying the exact equivalence $J\bo-$ proves the assertion about equal partner weights.

\subsection{Circle-product relations}
The cyclicity lemma allows the growing level-nine calculation for $W$ to be supplemented by a fourth-order one. Put $q=o(W)|_{J_{h_J}}$. Then $q\ne0$: otherwise the projection of $Y(W,z)j$ to the first Virasoro summand would vanish, so all terms of weight at most $h_J+2n$ would vanish. This would give $[j]*[W]=0$, forcing $w_S=0$ for every simple $S$ by \eqref{A:rat:nonzero}. The already constructed $X$ has a nonzero $W$ eigenvalue, a contradiction.

Normalize the projection by $q$ and write
\begin{equation}\label{A:rat:Fseries}
 q^{-1}[Y(W,z)j]=z^{-2n}\sum_{\ell\geq0}f_\ell(x_l,x_r)z^\ell[j],
 \qquad f_0=1.
\end{equation}
Before $\ell=3n$ only the first Virasoro summand occurs. Thus these initial coefficients are determined without any higher-channel normalization.

Here is a compact specification of the complete fourth-order differential operator. Set $\theta=z\partial_z$. When moving $L_{-r}$ from the second input, whose existing descendant level is $L$, use
\begin{equation}\label{A:rat:Dr}
 \mathscr D_r^{(L)}=
 (-1)^r\big(rx_r-x_l+h_J+2n+L+\theta\big)
 +z^{-r}\big((r-1)2n-\theta\big).
\end{equation}
Apply these operators in the order dictated by each monomial of \eqref{A:rat:null}, starting with its rightmost factor, and sum with the null-vector coefficients. Multiply by $z^4$ to obtain
$\sum_{s=0}^{4}z^sQ_s(\theta)$. Its indicial denominator on \eqref{A:rat:Fseries} is exactly
\begin{equation}\label{A:rat:indicial}
 D_\ell=\ell(\ell-3n)
      \bigl(\ell^2-(3n+2)\ell-6n+1\bigr).
\end{equation}
The quadratic discriminant is $9n(n+4)$, so it has no rational roots. For odd $n\geq3$, all denominators through $\ell=2n+2$ are therefore nonzero. This gives an exact four-step recurrence over $\Q[x_l,x_r]$.

The two circle products give genuine bimodule relations
\begin{equation}\label{A:rat:circles}
 R_{n,k}(x_l,x_r)[j]=0,\qquad
 R_{n,k}=\sum_{i=0}^{2n}\binom{2n}{i}f_{2n+k-i},\qquad k=1,2.
\end{equation}
Also
\[
 q^{-1}G_n(x_l,x_r)=\sum_{i=0}^{2n-1}\binom{2n-1}{i}f_{2n-i}.
\]
In particular, $w_S/q\in\Q$ for every simple $S$, since its ordered weight pair is rational. This is a statement about the semisimple quotient of the Zhu algebra; it does not exclude its radical.

\paragraph{The recurrence.}
The following finite rule spells out every coefficient of the operator in \eqref{A:rat:Dr}. For a PBW term of \eqref{A:rat:null} with coefficient $b$, start with one pair $(s,b)=(0,b)$ and descendant level $L=0$. For each negative mode $L_{-r}$, starting on the right, replace each pair $(s,b)$ by
\[
\begin{split}
&\bigl(s,\ b(-1)^r(rx_r-x_l+3n/4+2n+L+e+s)\bigr),\\
&\bigl(s-r,\ b((r-1)2n-e-s)\bigr).
\end{split}
\]
Combine equal shifts, replace $L$ by $L+r$, and continue. Sum the five PBW terms and add four to every shift. The coefficient at shift $s$ is $Q_s(e)$. With $f_0=1$, the exact recurrence is
\[
D_\ell f_\ell=-\sum_{s=1}^{\min(4,\ell)}Q_s(-2n+\ell-s)f_{\ell-s},
\]
whenever $D_\ell\ne0$. Together with \eqref{A:rat:circles}, this specifies the circle polynomials explicitly.

\subsection{Extensions of the vacuum}
Theorem~\ref{A:rat:main} removes all same-weight extension pairs. Generalized $L_0$ eigenspaces modulo $\Z$, and the two monodromy sectors of $J$, are direct-sum decompositions of the category. Consequently a remaining nonzero $\Ext^1(S,T)$ must satisfy
\[
 h_S-h_T\in\Z\setminus\{0\},\qquad
 h_{S^J}-h_{T^J}\in\Z\setminus\{0\},
 \qquad \kappa_S=\kappa_T.
\]
In particular, both $\Ext^1(V,X)$ and $\Ext^1(J,Y)$ vanish.

There is a useful additional reduction for extensions of the vacuum. If a simple $T$ has $h_T>1$, an extension $0\to T\to E\to V\to0$ splits: lift the vacuum to $v\in E_0$. Then $L_{-1}v\in T_1=0$. A translation-invariant module vector is vacuum-like, so $a\mapsto a_{-1}v$ gives a module map $V\to E$ splitting the quotient. Nonintegral $h_T$ also gives a split extension by the $L_0$ residue decomposition, and $h_T=0$ is covered by Theorem~\ref{A:rat:main}.

At this stage a possible obstruction to projectivity of $V$ must have lowest weight in the explicit finite arithmetic set
\begin{equation}\label{A:rat:badset}
 \begin{split}
 \mathcal B_n=\{k\in\Z:\;&-\lfloor n/4\rfloor\leq k\leq1,\quad k\ne0,\\
 &n(n+4k)\text{ or }(n+4)(n+4k)\text{ is a nonnegative integer square}\}.
 \end{split}
\end{equation}
There are at most two simple candidates at each such weight. The next argument removes all negative weights from this obstruction set, including at the odd-square parameters relevant to the Haagerup--Izumi family.

\subsection{Extensions with negative lowest weight}
Fix $t>1$ with $t+t^{-1}=n+2$, and use momentum notation
\[
 h(p)=\frac{p^2-(t-1)^2}{4t}.
\]
For rational labels $r,s$, the weight $h(rt-s)$ is rational precisely when $r^2=s^2$. This follows by reducing $t^{-1}=n+2-t$ and using irrationality of $t$.

The level-nine null vector of $W=L(c,h_{3,3})$ gives the usual nine-root fusion restriction
\begin{equation}\label{A:rat:Wroots}
 p\longmapsto p+2it-2j,\qquad i,j\in\{-1,0,1\}.
\end{equation}
To see the restriction explicitly, apply the Virasoro bimodule identity \eqref{A:rat:VirZhu} to the level-nine singular vector in Appendix~\ref{app:null}, with input weight $h$ and output weight $\widetilde h$. The resulting monic polynomial is
\begin{equation}\label{A:rat:Wpoly}
\begin{split}
 P_{3,3}(h,\widetilde h)={}&(\widetilde h-h)
 \prod_{a\in\{1,t,t-1,t+1\}}\left[
 \left(\widetilde h-h-\frac{a^2}{t}\right)^2
 -\frac{a^2\bigl(4th+(t-1)^2\bigr)}{t^2}\right].
\end{split}
\end{equation}
For $h=h(p)$, each quadratic factor has roots $h(p+2a)$ and $h(p-2a)$. Together with the root $h(p)$ these are exactly the nine weights in \eqref{A:rat:Wroots}. Every nonzero three-point coupling annihilates this polynomial, which gives the necessary fusion restriction used below.

\begin{lemma}\label{A:rat:genericorbit}
Suppose a simple ordinary $V$-module has lowest weight $h(k(t-1))$ or $h(k(t+1))$, where $k\in\Q\setminus\Z$. Its Virasoro constituents lie, respectively, in the sets
\[
 \{L(c,h((k+2m)(t-1))):m\in\Z\},\qquad
 \{L(c,h((k+2m)(t+1))):m\in\Z\}.
\]
All these constituents are irreducible Verma modules in singleton Virasoro blocks. In particular, none lies in a block containing a constituent of $V$.
\end{lemma}
\begin{proof}
For the first family, a shift in \eqref{A:rat:Wroots} changes the rational labels to $r=k+2i$, $s=k+2j$. Rational output weight requires
\[
 (k+2i)^2-(k+2j)^2=4(i-j)(k+i+j)=0.
\]
Since $k\notin\Z$, this forces $i=j$, giving the stated step-two orbit. For the second family the corresponding equation is
\[
 (k+2i)^2-(k-2j)^2=4(i+j)(k+i-j)=0,
\]
so $i=-j$. These nonintegral labels cannot be integer Kac labels, even after a momentum sign change. Feigin--Fuchs linkage therefore makes each orbit weight a singleton block \cite[Theorem 3.1]{A:BNW}.

To propagate the restriction through the module, restrict the ordinary module to Virasoro and apply Lemma~\ref{A:blocks:union} to extend the category-$\mathcal O$ block projections. The hypotheses here include ordinary $L_0$-action; generalized modules are not being projected by this argument. A fixed singleton or two-element block has finite length, since its possible composition-factor lowest spaces are finite-dimensional. If the projection of $Y(W,z)$ from a simple constituent to a block is nonzero, first take the Virasoro submodule spanned by its coefficients and then a simple quotient of that image. The block has finite length by the preceding weight-space bound, so such a quotient exists and the composed intertwiner is nonzero. Its output weight must satisfy \eqref{A:rat:Wroots}. Any permitted output belongs to a singleton block in the displayed orbit; hence a block outside that orbit cannot receive a nonzero image. Thus no block outside the displayed orbit is reached. Singleton ordinary blocks are semisimple. Iterating from the lowest vector and using strong generation by $T,W$ proves the claim. The blocks of $V$ have positive integral diagonal Kac labels and their linked negative-label partners, so they are disjoint from these orbits.
\end{proof}

\begin{proposition}\label{A:rat:negativeext}
For every simple $T$ with negative lowest weight,
\[
 \Ext^1(V,T)=0.
\]
Consequently the only remaining possible obstruction to projectivity of $V$ is a simple module of lowest weight one.
\end{proposition}
\begin{proof}
Nonintegral lowest weights split by the $L_0$ residue decomposition. If $h_T<0$ is integral, the spectral theorem gives either $k=a/n$ or $k=a/(n+4)$ with $0<|k|<1$. The zero label would give $h_T=-n/4$, which is not integral for odd $n$. Lemma~\ref{A:rat:genericorbit} applies.

In an extension $0\to T\to E\to V\to0$, the nilpotent part of $L_0$ factors through a homomorphism $V\to T$, and is therefore zero. Thus $E$ is ordinary, so no generalized-category-$\mathcal O$ assertion is being imported. The disjoint Virasoro blocks split the sequence as Virasoro modules. The section with the Virasoro constituents of $V$ is stable under $W$: fusion of the $C_1$-cofinite $W$ with the diagonal constituents of $V$ has $C_1$-cofinite image, and cannot reach any of the singleton generic blocks of $T$ \cite{A:CJORY}. Hence this section is a $V$-submodule and splits the extension. Combining this with the previous $h_T>1$, nonintegral and $h_T=0$ arguments proves the last statement.
\end{proof}

\subsection{The first weight-one candidate}
Let $x_0\in X_0$, $y_0\in Y_{-n/4}$ be nonzero. The lattice construction gives
\[
 x_0=\one+2\lambda\varepsilon,\qquad
 L_{-1}x_0\ne0,\qquad (x_0,x_0)\ne0,
 \qquad o(W)x_0=w_Xx_0,\quad w_X\ne0.
\]
The vector $\phi=L_{-1}x_0$ is a Virasoro highest vector of weight one. Use the nonzero intertwiner $\mathcal Y$ of type $\binom X{J\ Y}$ and the invariant form on $X$. For homogeneous $v\in J$ of weight $b$, write $T_s(v)$ for the coefficient of $z^{s+n/4-b}$ in $\mathcal Y(v,z)y_0$, and set
\[
 \Phi_0(v)=(x_0,T_0(v)),\qquad
 \Phi_1(v)=(\phi,T_1(v)).
\]
The first functional is the actual lowest-space Zhu map, and $\Phi_0(j)\ne0$. The $L_1$ bracket identity gives
\begin{equation}\label{A:rat:Phi}
 \Phi_1(v)=\Phi_0(L_1v)+(b+n/4)\Phi_0(v),
 \qquad \Phi_1(j)=n\Phi_0(j)\ne0.
\end{equation}
The second functional is a Virasoro highest-to-highest functional with ordered weights $(h,h')=(-n/4,1)$; it is not a lowest-$V$-space map.

Put $w=2n$ and $C=W\circ j=\sum_{i=0}^w\binom wi W_{i-2}j$. Since $\Phi_0(C)=0$, \eqref{A:rat:Phi} and
\[
 [L_1,W_{i-2}]=(2w-i)W_{i-1}
\]
give, by shifting one binomial sum,
\begin{align*}
 \Phi_1(C)
 &=\sum_{i=0}^w\binom wi
       \bigl((2w-i)\Phi_0(W_{i-1}j)-i\Phi_0(W_{i-2}j)\bigr)\\
 &=w\sum_{i=0}^w\binom wi\Phi_0(W_{i-1}j)
 =2n\,w_X\Phi_0(j).
\end{align*}
For $n\geq3$, all terms of $C$ lie below the next Virasoro primary of $J$. Evaluating their first-channel reduction on $\Phi_1$ therefore proves the exact identity
\begin{equation}\label{A:rat:lowerexclusion}
 q\,R_{n,1}(1,-n/4)=2w_X\ne0.
\end{equation}


\begin{corollary}\label{A:rat:onepair}
For odd $n\geq3$, the module $Y$ is the unique simple of lowest weight $-n/4$. There is at most one simple module $T$ of lowest weight one, and, if it exists, it satisfies
\[
 h_{J\bo T}=\frac{3n}{4}+4.
\]
\end{corollary}
\begin{proof}
At $h=-n/4$, the spectral theorem permits partners of weight $0$ or $1$. Identity \eqref{A:rat:lowerexclusion} excludes the second; ordered-pair uniqueness identifies the first with $Y$. At $h=1$ the first quadratic branch is impossible because $n(n+4)$ is not a square. The two second-branch choices have partners $-n/4$ and $3n/4+4$. The former is excluded by applying the involution $J\bo-$ to \eqref{A:rat:lowerexclusion}. Ordered-pair uniqueness gives the remaining at-most-one assertion.
\end{proof}

\subsection{The second weight-one candidate}
Specialize the fourth-order operator to
\[
 x_r=1,\qquad x_l=\frac{3n}{4}+4.
\]
For its scalar solution $F(z)=z^{-2n}\sum_{\ell\geq0}f_\ell z^\ell$, set
\[
 P(z)=z^{2n}(1+z)^{2n}F(z).
\]
After this gauge transformation and division by $1+z$, the equation is
\begin{equation}\label{A:rat:gauged}
 \mathcal G_n P=0,\qquad
 \mathcal G_n=\sum_{k=0}^3 z^k C_k(\theta),\qquad \theta=z\partial_z,
\end{equation}
where the following four polynomials specify the operator completely:
\begin{align*}
 C_0(e)={}&e(e-3n)\bigl(e^2-(3n+2)e-6n+1\bigr),\\
 C_1(e)={}&3e^4-(22n+14)e^3+(51n^2+54n+13)e^2\\
          &+(-36n^3-18n^2+9n-6)e-54n^3-90n^2+36n,\\
 C_2(e)={}&3e^4-(26n+22)e^3+(79n^2+126n+45)e^2\\
          &-(96n^3+206n^2+129n+14)e+36n^4+78n^3+36n^2-32n,\\
 C_3(e)={}&(e-2n)(e-3n-4)
                  \bigl(e^2-(5n+6)e+6n^2+14n+9\bigr).
\end{align*}
The recurrence \eqref{A:rat:Dr} gives these coefficients as polynomial identities in $n$.

\begin{lemma}\label{A:rat:twocircle}
For every odd integer $n\geq3$,
\[
 \bigl(R_{n,1}(3n/4+4,1),R_{n,2}(3n/4+4,1)\bigr)\ne(0,0).
\]
Consequently there is no simple $V$-module of lowest weight one.
\end{lemma}
\begin{proof}
Suppose both circle coefficients vanish. The recurrence determines $f_\ell$ uniquely through $\ell=2n+2$, because those levels precede $3n$. Put
\[
 p_\ell=\sum_{i=0}^{\min(2n,\ell)}\binom{2n}{i}f_{\ell-i},
 \qquad P(z)=\sum_{\ell=0}^{2n}p_\ell z^\ell.
\]
Then $P(0)=1$ and $\deg P\leq2n$. The assumed circle relations say $p_{2n+1}=p_{2n+2}=0$. The gauged recurrence therefore makes every coefficient of $\mathcal G_nP$ through degree $2n+2$ vanish. Its only possible further coefficient is at degree $2n+3$, and this is $C_3(2n)p_{2n}=0$. Thus $\mathcal G_nP=0$. This step uses no coefficient at or beyond the resonant level $3n$.

The same equation has an explicit polynomial solution
\begin{equation}\label{A:rat:highsol}
 S_n(z)=z^{3n}H_n(z),
\end{equation}
where
\begin{equation}\label{A:rat:Hpoly}
 \begin{split}
 H_n(z)={}&1+2z+\frac{2(4n-1)}{6n-1}z^2\\
 &+\frac{2(2n-1)(5n-4)}{3(3n-4)(6n-1)}z^3
   +\frac{n(2n-1)}{3(3n-4)(6n-1)}z^4.
 \end{split}
\end{equation}
Substitution in \eqref{A:rat:gauged} verifies this identity coefficient by coefficient. In particular,
\begin{equation}\label{A:rat:Hminusone}
 \deg S_n=3n+4,\qquad H_n(0)=1,\qquad
 H_n(-1)=-\frac{4(2n-1)}{3(3n-4)(6n-1)}\ne0.
\end{equation}

At $z=-1$, the lowest-order coefficient of $\mathcal G_n(z+1)^r$ is the indicial polynomial
\begin{equation}\label{A:rat:indminusone}
 I_{-1}(r)=r(r-2n-8)\bigl(r^2-(2n+6)r+2n+9\bigr).
\end{equation}
The two roots of its quadratic factor are $n+3\pm\sqrt{n(n+4)}$, and are irrational. It follows directly by examining the first nonzero Taylor coefficient that any nonzero polynomial solution vanishing at $-1$ has vanishing order exactly $2n+8$. No connection formula or absence-of-logarithms assumption is needed for this assertion.

Now form the polynomial Wronskian
\[
 \mathscr W=P S_n'-P'S_n.
\]
At zero, $P(0)=1$ and $S_n=z^{3n}(1+O(z))$, so
\[
 \operatorname{ord}_0\mathscr W=3n-1;
\]
in particular, $\mathscr W\ne0$. Since $S_n(-1)\ne0$, the nonzero polynomial solution
\[
 Q=P-\frac{P(-1)}{S_n(-1)}S_n
\]
vanishes at $-1$. It is nonzero because $Q(0)=1$. By \eqref{A:rat:indminusone}, $\operatorname{ord}_{-1}Q=2n+8$. Replacing $P$ by $Q$ leaves $\mathscr W$ unchanged, so
\[
 \operatorname{ord}_{-1}\mathscr W=2n+7.
\]
Consequently this nonzero polynomial must have degree at least
\[
 (3n-1)+(2n+7)=5n+6.
\]
But its definition gives
\[
 \deg\mathscr W\leq\deg P+\deg S_n-1\leq5n+3,
\]
a contradiction. The two circle coefficients cannot both vanish. Corollary~\ref{A:rat:onepair} and the Zhu relations \eqref{A:rat:circles} therefore exclude the last weight-one candidate.
\end{proof}

\subsection{Projectivity of the vacuum and rationality}

\begin{lemma}[Projective unit criterion]\label{A:rat:projectiveunit}
Let $\mathcal D$ be a finite $\C$-linear abelian braided $r$-category with simple unit and right-exact tensor product. If its unit is projective, then $\mathcal D$ is rigid and semisimple.
\end{lemma}
\begin{proof}
Here an $r$-category is a Grothendieck--Verdier category whose dualizing object is the tensor unit; write $D$ for its contravariant duality. For a nonzero indecomposable $M$, the evaluation $DM\otimes M\to\one$ corresponds to $\id_{DM}$ under the defining adjunction and is therefore nonzero. Since the unit is simple, this evaluation is an epimorphism. Projectivity of $\one$ splits it, so $M$ is non-negligible in the exact sense of Etingof--Penneys \cite[Theorem 1.1]{A:EP}.

Moderate growth is automatic here and does not require already knowing rigidity. Let $K$ be the maximum of $1$ and the lengths of $S_i\otimes S_j$ over the finitely many simple objects. Right exactness in each variable gives
\[
 \ell(A\otimes B)\leq K\ell(A)\ell(B),\qquad
 \ell(M^{\otimes r})\leq K^{r-1}\ell(M)^r.
\]
Also $\dim\operatorname{Hom}(A,B)\leq\ell(A)\ell(B)$ by induction on composition series. Thus $\dim\operatorname{End}(M^{\otimes r})<r!$ for sufficiently large $r$. Etingof--Penneys applies to each indecomposable, proving rigidity. Finally,
\[
 \operatorname{Hom}(M,-)\cong
 \operatorname{Hom}(\one,M^*\otimes-)
\]
is exact, because rigid tensoring is exact and the unit is projective. Every object is therefore projective, so the finite abelian category is semisimple.
\end{proof}

\begin{theorem}[Uniform rationality]\label{A:rat:criterion}
For every odd $n\geq3$,
\[
 V_{n+2}\text{ is strongly rational}.
\]
In particular this holds for every $n=N^2$ with $N\geq3$ odd.
\end{theorem}
\begin{proof}
Lemma~\ref{A:rat:twocircle} proves that there is no simple of weight one. Proposition~\ref{A:rat:negativeext}, together with the other vacuum-extension vanishings, gives $\Ext^1(V,T)=0$ for every simple $T$. Induction on length then gives $\Ext^1(V,M)=0$ for every object $M$, so $V$ is projective. The $C_2$-cofinite theory makes the category of grading-restricted generalized $V$-modules finite abelian and braided, with right-exact tensor product \cite[Section~2.3]{A:McR}. Closure under contragredients gives Grothendieck--Verdier duality with dualizing object $V^\vee$ \cite[Section~4.1]{A:CMSY}. Since $V^\vee\cong V$, this is an $r$-category. Lemma~\ref{A:rat:projectiveunit} therefore proves semisimplicity. In the $C_2$-cofinite setting this is rationality \cite{A:McR}, and Section~\ref{sec:diagonal} establishes the other strong-rationality properties.
\end{proof}

\subsection{Foundational properties}\label{A:cat:inputs}
For odd $n\geq3$, let $\mathcal C_n=\operatorname{Rep}(V_{n+2})$. We collect the conclusions of Sections~\ref{sec:inputs}--\ref{A:app:rat} in a form suited to the categorical reconstruction and module classification.
\begin{theorem}[Foundational properties of the diagonal family]\label{E:foundations}
For every odd $n\geq3$, the following statements hold.
\begin{enumerate}[label=(H\arabic*),leftmargin=11mm,itemsep=5pt]
\item $V=V_{n+2}$ is simple, self-contragredient, of CFT type, $C_2$-cofinite and rational, with central charge $1-6n$. Its ordinary module category $\mathcal C_n$ is modular.
\item There is a self-dual invertible $J$ with $J^2=\one$, associator $F(J,J,J)=-1$, self-braiding $c_{J,J}=i^n$, twist $\theta_J=i^{-n}$, and spherical dimension $d(J)=-1$. Its centralizer gives the ribbon factorization
\[
 \mathcal C_n\simeq\Scat_n\bo\D_n,\qquad \Scat_n=\langle J\rangle.
\]
\item The deformed lattice VOA $\mathcal U_n=(V_{\Z\beta},T_\lambda)$, $\beta^2=2n$, $\lambda\ne0$, is a conformal extension of $V$, and as a $V$-module
\[
 A=\mathcal U_n=\one\oplus X,\qquad X\text{ simple},\quad h_X=0,\quad X\not\cong\one.
\]
It is simple and self-contragredient. The projection pairing $A\otimes A\xrightarrow{m}A\to\one$ is nondegenerate. Also $A\in\D_n$ and $\theta_A=\id$.
\item The local $A$-modules in $\D_n$ are the pointed category $\Pcat_n$ in \eqref{A:cat:NGcenter}. If $L_x$ denotes the local object labeled $x\in\Z/n$ and $U$ is restriction to $\D_n$, then
\begin{equation}\label{A:cat:branch}
 U(L_0)=\one\oplus X,\qquad U(L_x)=Z_x\quad(x\ne0),
\end{equation}
where each $Z_x$ is simple, $Z_x\cong Z_{-x}$, and these are the only repetitions. The $Z_x$ differ from $\one,X$. The pointed category has trivial associator in the gauge $c_{x,y}=e^{2\pi ixy/n}$ and dimensions one.
\item The simple $Y=JX$ is the unique simple module of minimum lowest weight $-n/4$, with one-dimensional lowest space. With
\[
 \Lambda_{a,m}(\tau)=\frac{1}{\eta(\tau)}
 \sum_{r\in\Z}q^{(2mr+a)^2/(4m)},\qquad q=e^{2\pi i\tau},
\]
the characters satisfy
\begin{equation}\label{A:cat:characters}
 \chi_V=\tfrac12(\Lambda_{n,n}-\Lambda_{n+4,n+4}),\quad
 \chi_X=\tfrac12(\Lambda_{n,n}+\Lambda_{n+4,n+4}),\quad
 \chi_Y=\tfrac12(\Lambda_{0,n}+\Lambda_{0,n+4}).
\end{equation}
\end{enumerate}
\end{theorem}
\begin{proof}
Section~\ref{sec:diagonal} establishes simplicity, CFT type, self-contragredience and strong generation by $T,W$. Theorem~\ref{A:c2:thm:main} proves $C_2$-cofiniteness. The lattice and simple-current constructions in Theorem~\ref{A:lat:main} then provide the inputs to Theorem~\ref{A:rat:criterion}, which proves rationality. McRae's modularity theorem \cite[Main Theorem~3]{A:McR} gives modularity of $\mathcal C_n$. This proves \textup{(H1)}.

Theorem~\ref{A:lat:main} proves the simple-current data and semion factorization in \textup{(H2)}, as well as the conformal embedding and decomposition $\mathcal U_n=V\oplus X$ in \textup{(H3)}. The invariant form constructed in \eqref{A:lat:ell} is nondegenerate and restricts nondegenerately to $V$. The orthogonal projection onto $V$ therefore identifies its product pairing with the projection pairing in \textup{(H3)}. Section~\ref{A:lat:condensation} proves that $J$ centralizes $A$, so $A\in\D_n$; the integral lattice grading gives $\theta_A=\id$.

The lattice local-module calculation in Section~\ref{A:lat:condensation} identifies the local category with \eqref{A:lat:pointed}. Together with the simplicity and isomorphism statements in Theorem~\ref{A:lat:main}(iii), it gives the branching and pointed structure in \textup{(H4)}. Finally, the character calculation in Section~\ref{A:app:lat} gives the three formulas in \textup{(H5)}. Theorem~\ref{A:rat:main} gives the lower bound $-n/4$, and Corollary~\ref{A:rat:onepair} proves that $Y$ uniquely attains it. Its lowest space is one-dimensional by the lattice construction. This proves \textup{(H5)}.
\end{proof}

\section{Categorical reconstruction}\label{A:sec:categorical}
Sections~\ref{A:sec:categorical}--\ref{B:sec:characters} use the foundational properties established in Theorem~\ref{E:foundations}. This section proves the categorical and extension-uniqueness assertions of Theorem~\ref{A:cat:main}.

\subsection{The extension algebra}
\begin{lemma}\label{A:cat:etale}
The algebra $A$ is connected \'etale, $\Ncat_n=\D_{n,A}$ is a fusion category, and the spherical structures and the center equivalence \eqref{A:cat:NGcenter} are defined.
\end{lemma}
\begin{proof}
Connectedness follows from $A=\one\oplus X$ with $X\not\cong\one$. It is commutative and associative by the VOA extension. Its only possible proper nonzero $\D_n$-subobjects are $\one$ and $X$. The unit action rules out $\one$ as an ideal, and the nonzero vacuum channel in $XX$ rules out $X$. Hence $A$ is a simple commutative algebra.

Creutzig--McRae--Shimizu--Yadav \cite[Theorem 2.26]{A:CMSY} proves that modules over a simple commutative algebra in a braided fusion category form a fusion category. This theorem requires neither unitarity nor a prior nonzero categorical-dimension calculation. Thus $\D_{n,A}$ is semisimple. Kong--Zheng \cite[Theorem 6.10]{A:KZ} then gives separability in characteristic zero. This proves that $A$ is connected \'etale.

For completeness, the nonzero dimension needed for the spherical module structure can be checked without assuming it at the beginning. Work with common left and right duals from the spherical structure. Let $b=\varepsilon m:A\otimes A\to\one$ be the nondegenerate projection pairing, and let $f:A\to A^*$ be its adjoint, so $b=\operatorname{ev}_A(f\otimes\id_A)$. Associativity makes $f$ an isomorphism of the appropriate regular $A$-modules. Braided commutativity gives $b\,c_{A,A}=b$. Together with $\theta_A=\id_A$, this is the symmetry identity
\[
 f^*j_A=f,
\]
where $j_A:A\to A^{**}$ is the ribbon pivotal map. This is the pairing calculation in \cite[Lemma~1.13]{A:KO}; it uses the inverse pairing, but no division by $d(A)$.

The algebra dimension used in \cite[Definition~6.1]{A:KZ} is therefore
\begin{align*}
 \dim_{\rm alg}(A)
 &=\operatorname{ev}_A(f\otimes f^{-1})\operatorname{coev}_A\\
 &=\operatorname{ev}_{A^*}\bigl((f^{-1})^*f\otimes\id_{A^*}\bigr)\operatorname{coev}_A
 =\operatorname{ev}_{A^*}(j_A\otimes\id_{A^*})\operatorname{coev}_A
 =d_{\D_n}(A).
\end{align*}
The regular $A$-module is simple and $\Hom(\one,A)=\C$. Separability and \cite[Theorem~6.3]{A:KZ} give $\dim_{\rm alg}(A)\ne0$, hence $d(A)\ne0$. Now the spherical module structure and pivotal free-module functor follow from \cite[Theorems~1.17--1.18 and Remark~1.19]{A:KO}. The center theorem \cite[Corollary~3.30]{A:DMNO} gives \eqref{A:cat:NGcenter}.
\end{proof}

\subsection{The rank of the module category}
\begin{lemma}\label{A:cat:ranklemma}
Let $A$ be connected \'etale in a nondegenerate braided fusion category $\D$, put $\Ncat=\D_A$, and let $\Pcat=\D_A^{\mathrm{loc}}$. Then
\begin{equation}\label{A:cat:rankidentity}
 \rank\Ncat=\sum_{M\in\Irr(\Pcat)}\dim_{\C}\End_{\D}(U(M)).
\end{equation}
\end{lemma}
\begin{proof}
Let $F:Z(\Ncat)\to\Ncat$ be forgetful and $I$ its right adjoint. For the canonical algebra $K=I(\one)$, adjunction and the center-induction formula \cite[Proposition 9.2.2]{A:EGNO} give
\[
 \dim\End_{Z(\Ncat)}K=\dim\Hom_{\Ncat}(F I(\one),\one)
 =\sum_{T\in\Irr(\Ncat)}\dim\Hom_{\Ncat}(T\otimes T^*,\one)
 =\rank\Ncat.
\]
Under $Z(\Ncat)\simeq\D\bo\Pcat^{\mathrm{rev}}$, the underlying object of $E\bo M$ is $F_A(E)\otimes_A M$. Therefore
\[
 \Hom_{Z(\Ncat)}(E\bo M,K)
 \cong\Hom_{\Ncat}(F_A(E),M^*)
 \cong\Hom_{\D}(E,U(M^*)).
\]
Consequently, as an object,
\[
 K\cong\bigoplus_{M\in\Irr(\Pcat)}U(M^*)\bo M.
\]
Different $M$ are orthogonal simple labels in the second factor. Taking endomorphism dimensions and relabeling by duality proves \eqref{A:cat:rankidentity}.
\end{proof}

Apply the lemma to \eqref{A:cat:branch}. The unit local module contributes two, and each of the other $n-1$ local modules contributes one. $Z_x\cong Z_{-x}$ does not merge the distinct second-factor labels $L_x,L_{-x}$. Hence
\begin{equation}\label{A:cat:rankN}
 \rank\Ncat_n=2+(n-1)=n+1.
\end{equation}
The $n$ local modules already give a full pointed subcategory with group $\Z/n$. There is exactly one remaining simple, denoted $R$. The next calculation proves that it is noninvertible and fixes its fusion coefficient.

\subsection{Frobenius--Perron dimensions}
The following lemma justifies the character method in this nonunitary setting.
\begin{lemma}\label{A:cat:PFlemma}
Let $V$ be a strongly rational VOA of central charge $c$, with a unique minimum-weight simple $Y$ whose lowest space is one-dimensional. Suppose the vacuum character has a nonzero contribution at the corresponding exponential rate, with
\[
 \lim_{\epsilon\downarrow0}e^{-\pi(c-24h_Y)/(12\epsilon)}\chi_V(i\epsilon)>0.
\]
Then
\[
 \FP(M)=\frac{S_{MY}}{S_{VY}}
 =\lim_{\epsilon\downarrow0}\frac{\chi_M(i\epsilon)}{\chi_V(i\epsilon)}
\]
for every simple $M$.
\end{lemma}
\begin{proof}
Modular transformation and uniqueness of $Y$ show that the coefficients $v_M=S_{MY}$ of the fastest exponential are real and nonnegative, since the characters have nonnegative graded dimensions. By hypothesis $v_V>0$. Verlinde diagonalization says that $v$ is a simultaneous eigenvector of the nonnegative fusion matrices $N_M$, with eigenvalue $\lambda_M=v_M/v_V\geq0$.

Since $M\otimes M^*$ contains the unit, $N_MN_{M^*}\geq I$ entrywise. Applying this to the nonzero vector $v\geq0$ yields $\lambda_M\lambda_{M^*}\geq1$. Thus all $\lambda_M$ and all entries of $v$ are strictly positive. A nonnegative matrix with a strictly positive eigenvector has that eigenvalue as its spectral radius. Therefore $\lambda_M=\FP(M)$, and the asserted limit follows.
\end{proof}

Poisson summation and the eta transformation give
\[
 \Lambda_{a,m}(i\epsilon)\sim\frac{1}{\sqrt{2m}}
 e^{\pi/(12\epsilon)}.
\]
By \eqref{A:cat:characters}, set
\begin{equation}\label{A:cat:tdef}
 t=\FP(X)=\frac{(2n)^{-1/2}+(2n+8)^{-1/2}}
 {(2n)^{-1/2}-(2n+8)^{-1/2}}
 =\frac{n+2+\sqrt{n(n+4)}}2.
\end{equation}
The denominator is positive. Useful identities are
\begin{equation}\label{A:cat:tidentities}
 t+t^{-1}=n+2,\qquad (t-1)^2=nt,\qquad \FP(A)=1+t.
\end{equation}

\subsection{The near-group fusion rules}\label{A:cat:neargroup}
The dimension formulas for connected \'etale algebras \cite[Lemma 3.11 and Corollary 3.32]{A:DMNO} imply
\begin{equation}\label{A:cat:globalFP}
 \FP(\D_n)=n(1+t)^2,\qquad \FP(\Ncat_n)=n(1+t).
\end{equation}
Let $r=\FP(R)$. Subtracting the $n$ known invertible simples gives
\[
 r^2=n(1+t)-n=nt=(t-1)^2.
\]
Therefore $r=t-1>1$, so $R$ is noninvertible. The local invertibles permute the noninvertible simples on both sides; there is just one. Hence $gR=Rg=R$ and $R^*=R$. Frobenius reciprocity makes every invertible occur once in $R^2$. Write
\[
 R^2=\bigoplus_g g\oplus mR.
\]
Dimensions now give
\[
 m=\frac{r^2-n}{r}=\frac{n(t-1)}{t-1}=n.
\]
This proves \eqref{A:cat:NG}.

An exact induction decomposition will also be needed. Adjunction and \eqref{A:cat:branch} show that $F_A(X)$ contains the unit once and no other local invertible. Thus $F_A(X)=\one\oplus bR$. Since the tensor functor $F_A$ preserves Frobenius--Perron dimensions, $t=1+b(t-1)$, and
\begin{equation}\label{A:cat:indX}
 F_A(X)=\one\oplus R.
\end{equation}

\subsection{Spherical dimensions}\label{A:cat:spherical}
Write $x=d_{\D_n}(X)$. The free-module functor is pivotal. By \eqref{A:cat:indX}, $x=1+d(R)$, while \eqref{A:cat:NG} and the fact that odd-order invertibles have spherical dimension one give
\[
 d(R)^2=n+n\,d(R).
\]
Consequently
\[
 x\in\{t,t^{-1}\},\qquad x>0.
\]
The character asymptotics determine the root. By (H5) and \eqref{A:cat:characters},
\begin{equation}\label{A:cat:a0}
 a:=S_{VY}=\frac12\left(\frac1{\sqrt{2n}}-\frac1{\sqrt{2(n+4)}}\right)>0.
\end{equation}
Because $Y=JX$ and $d(J)=-1$, $d(Y)=-x$. The categorical dimension formula for local modules gives
\[
 \dim\D_n=n(1+x)^2,\qquad \dim\mathcal C_n=2n(1+x)^2.
\]
The normalized VOA $S$-matrix satisfies $S_{VY}/S_{VV}=d(Y)$ and $S_{VV}^2=1/\dim\mathcal C_n$. Squaring eliminates any choice of sign in its normalization. Since $x>0$, \eqref{A:cat:a0} gives
\[
 \frac{x}{1+x}=\sqrt{2n}\,a
 =\frac12\left(1-\sqrt{\frac n{n+4}}\right).
\]
Solving yields
\begin{equation}\label{A:cat:negdim}
 d(X)=t^{-1},\qquad
 d(R)=t^{-1}-1=\frac{n-\sqrt{n(n+4)}}2.
\end{equation}
In particular the spherical near-group category obtained from the VOA has negative noninvertible dimension. Frobenius--Perron dimensions remain the positive numbers computed above.

\subsection{Condensation at a square parameter}\label{A:cat:condensation}
Now put $n=N^2$ with $N$ odd, and let $H=N\Z/N^2$. In $\Pcat_{N^2}^{\mathrm{rev}}$ its twists and braiding are trivial. Also
\[
 H^\perp=H,
\]
because the associated bicharacter is $e^{-4\pi ixy/N^2}$ and $2$ is invertible modulo $N$. Hence the central subcategory on $H$ is Tannakian and Lagrangian in this pointed factor.

Let $B=\bigoplus_{h\in H}h$ be its regular central algebra in
$Z(\Ncat_{N^2})\simeq\D_{N^2}\bo\Pcat_{N^2}^{\mathrm{rev}}$, and define
\begin{equation}\label{A:cat:Fdef}
 \Fcat_N:=\operatorname{Rep}_{\Ncat_{N^2}}(B).
\end{equation}
The forgetful functor identifies this central pointed subcategory with the invertibles $H\subset\Ncat_{N^2}$, so this is ordinary de-equivariantization by $\widehat H$. Its unit is simple. The local modules of the second factor over $B$ are $\mathrm{Vec}$. The center theorem for a central algebra \cite[Theorem 3.20]{A:DMNO} therefore gives
\begin{equation}\label{A:cat:centerF}
 Z(\Fcat_N)\simeq Z(\Ncat_{N^2})_B^{\mathrm{loc}}
 \simeq\D_{N^2}.
\end{equation}

There is also a concrete description directly in VOA modules. Put $\alpha=\beta/N$, so $\alpha^2=2$, and use the same conformal vector in the overlattice extension
\[
 B_N=(V_{\Z\alpha},T_\lambda).
\]
Condensation in stages identifies \eqref{A:cat:Fdef} with
\begin{equation}\label{A:cat:Fconcrete}
 \Fcat_N\simeq(\D_{N^2})_{B_N},\qquad
 B_N\cong\one\oplus X\oplus
 2\bigoplus_{r=1}^{(N-1)/2} Z_{Nr}.
\end{equation}
The regular algebra on $H$ is connected \'etale in the local $A$-module category, so its algebra over $A$, namely $B_N$, is connected \'etale in $\D_{N^2}$ \cite[Proposition 3.16]{A:DMNO}. Its twist is trivial. All simple $Z_{Nr}$ displayed here are distinct. The grading of this overlattice VOA can contain negative weights; no CFT-type extension theorem is used for it.

Each invertible local $A$-module has spherical dimension one, so its restriction has dimension $d(A)=1+t^{-1}$. The regular algebra on $H$ has $N$ such summands. Consequently
\[
d_{\D_{N^2}}(B_N)=N\,d_{\D_{N^2}}(A)=N(1+t^{-1})\ne0.
\]
Together with $\theta_{B_N}=\id$, this verifies the dimension hypothesis for the spherical module structure of \cite{A:KO} on \eqref{A:cat:Fconcrete}. Its central free-module functor is pivotal. Thus the braided equivalence \eqref{A:cat:centerF} also preserves twists: it is a ribbon equivalence with the inherited spherical center. As an independent rank check, the Lagrangian algebra $B_N$ corresponds to $I_{\Fcat_N}(\one)$, and hence
\[
 \rank\Fcat_N=\dim\End_{\D_{N^2}}(B_N)
 =1+1+4\frac{N-1}{2}=2N.
\]

\subsection{The Haagerup--Izumi fusion rules}\label{A:cat:HIfusion}
The proof first determines the simple objects, then their left and right actions and their products.

\textbf{Simple objects.} The free modules $gB$ on invertibles are indexed by $G/H=\Z/N$ and give precisely $N$ invertibles $\alpha_a$. By adjunction,
\[
 \End_B(RB)\cong\bigoplus_{h\in H}\Hom(R,Rh).
\]
Every summand is one-dimensional. Compositions are nonzero scalar multiples of the group law, so this algebra is a twisted group algebra of $H$. A cyclic group has $H^2(H,\C^\times)=0$; thus it is $\C^N$. Consequently $RB$ is a multiplicity-one sum of $N$ distinct simple objects $Q_i$. None is an $\alpha_a$, by adjunction. Induction is dominant, so these $2N$ objects exhaust the simples.

Put $\delta=(N+\sqrt{N^2+4})/2$, so $\delta^2=1+N\delta$ and $\FP(R)=N\delta$. If $d_i=\FP(Q_i)$, tensor induction and de-equivariantization give
\[
 \sum_i d_i=N\delta,\qquad
 \sum_i d_i^2=\frac{N^2+(N\delta)^2}{N}-N=N\delta^2.
\]
Equality in Cauchy--Schwarz implies $d_i=\delta$ for every $i$. They are therefore noninvertible.

\textbf{Free left and right actions.} For a noninvertible simple $Q$, write the total multiplicities of invertible and noninvertible terms in $QQ^*$ as $a,b$. Dimensions give
\[
 1+N\delta=a+b\delta.
\]
As $\delta$ is irrational, $a=1$ and $b=N$. Frobenius reciprocity identifies $a$ with the order of the left stabilizer of $Q$. Thus the left action of $\Z/N$ is free and, since there are $N$ noninvertibles, transitive. The same argument using $Q^*Q$ proves the right action free and transitive. Choose $\rho=Q_0$ and write the others as $\alpha_a\rho$. There is an automorphism $\varphi$ of $\Z/N$ with
\[
 \rho\alpha_a=\alpha_{\varphi(a)}\rho.
\]

\textbf{The dual action supplies all translations.} De-equivariantization gives $\Ncat_{N^2}\simeq\Fcat_N^{\widehat H}$. The dual action fixes every invertible $gB$: twisting the $B$-action is undone by rescaling the $H$-graded free factor. Noninvertible orbits each produce at least one noninvertible in the equivariantization. Since $\Ncat_{N^2}$ has just one, the $N$ noninvertibles of $\Fcat_N$ form a single orbit. Thus $\widehat H$, of order $N$, acts freely and transitively on them. It fixes the $\alpha_a$ and commutes with their left action, so on simple labels it acts by every translation
\[
 \alpha_a\rho\longmapsto\alpha_{a+t}\rho.
\]

\textbf{Inversion and self-duality.} The dimension argument above applies also to $\rho^2$, so it has a unique invertible summand, say $\alpha_\ell$. Every translation fixes that summand, but
\[
 (\alpha_t\rho)^2=\alpha_{t+\varphi(t)}\rho^2.
\]
Comparing the unique invertibles shows $t+\varphi(t)=0$ for every $t$; hence $\varphi=-1$. Write $\rho^*=\alpha_j\rho$. Dualizing again gives $\rho=\alpha_{2j}\rho$. Freeness and oddness force $j=0$. Thus $\rho^*=\rho$, and the invertible summand of $\rho^2$ is $\one$.

\textbf{Multiplicity one.} Write
\[
 \rho^2=\one\oplus\bigoplus_a m_a\alpha_a\rho,\qquad\sum_a m_a=N.
\]
Since $(\alpha_t\rho)^2=\rho^2$, applying each translation automorphism gives $m_{a+t}=m_a$. Therefore $m_a=1$ for all $a$. This proves all the rules \eqref{A:cat:HI}.

The spherical de-equivariantization is pivotal, so $d(R)=\sum_i d(Q_i)$. All $\alpha_a$ have dimension one (they have odd order), and the $Q_i=\alpha_a\rho$ have a common dimension. Equation \eqref{A:cat:negdim} now yields
\[
 d(\rho)=\frac{d(R)}N=\frac{N-\sqrt{N^2+4}}2,
\]
as asserted.

\subsection{The pointed associator}\label{A:cat:associator}
Work first with $G=\Z/N^2$, $H=N\Z/N^2$, and central half-braiding
\[
 \sigma_h(x)=e^{2\pi i k hx/N^2}\quad(h\in H,x\in G).
\]
The pointed category before condensation has trivial associator. Choose coset representatives $0\leq a<N$, and write $M_a=aB$. Let $u_{a,h}$ be its basis in degree $a+Nh$, $h\in\Z/N$, with standard right $B$-action. Its left action is
\[
 b_{Nr}\cdot u_{a,h}=e^{2\pi i k ra/N}u_{a,h+r}.
\]
Write $\overline{a+b}$ for reduction modulo $N$ and $s(a,b)=\lfloor(a+b)/N\rfloor$. The following are well-defined balanced multiplication isomorphisms:
\[
 f_{a,b}(u_{a,h}\otimes_B u_{b,t})
 =e^{2\pi i k hb/N}
 u_{\overline{a+b},\,h+t+s(a,b)}.
\]
For a triple of basis vectors, the left-associated composition of the $f$'s divided by the right-associated composition has scalar
\[
 e^{2\pi i k c s(a,b)/N}.
\]
Terms involving $h,t$ cancel modulo integral exponents, using
$b+c-\overline{b+c}=Ns(b,c)$. The skeletal associator is the map from the left-associated product to the right-associated one, so it is
\begin{equation}\label{A:cat:cocyclecalc}
 \omega_{-k}(a,b,c)=e^{-2\pi i k c s(a,b)/N}.
\end{equation}
In this construction the central factor is $\Pcat_{N^2}^{\mathrm{rev}}$, so $k=-1$. Equation \eqref{A:cat:cocyclecalc} is exactly \eqref{A:cat:omega}.

The carry identity proves the normalized pentagon equation. Its class is the standard generator of $H^3(\Z/N,\C^\times)\cong\Z/N$: the cyclic cocycle invariant
\[
 \prod_{j=0}^{N-1}\omega_1(1,j,1)=e^{2\pi i/N}
\]
is primitive. In particular this family is twisted, with pointed associator of maximal order $N$, including composite odd $N$. The remaining associators are inherited from relative tensor products in $\Fcat_N$.

\subsection{The center criterion}\label{A:cat:centercriterion}
Equation \eqref{A:cat:NGcenter} gives the Witt equality $[\D_n]=[\Pcat_n]$. A pointed nondegenerate category is Witt-trivial exactly when its metric group has a Lagrangian subgroup; see \cite[Section 5.3]{A:DMNO}. Such a subgroup has order $\sqrt n$, so it cannot exist for nonsquare $n$. For $n=N^2$, the subgroup $N\Z/N^2$ is Lagrangian as already proved. This establishes the final assertion of Theorem~\ref{A:cat:main}.

For odd squares $n=N^2$, one has $n\equiv1\pmod8$. Thus the semion factor in \eqref{A:cat:realization} can be specified uniformly by
\[
 J^2=\one,\quad F(J,J,J)=-1,\quad c_{J,J}=i,\quad
 \theta_J=-i,\quad d(J)=-1.
\]
This specifies the inherited ribbon structure, including its negative dimension.


\subsection{Uniqueness of the VOA extension}\label{sec:extension}
The near-group calculation gives
\[
 \dim\Hom(X^2,\one)=1,\qquad
 \dim\Hom(X^2,A)=\dim\Hom_{\N}((\one\oplus R)^2,\one)=2.
\]
Consequently $\dim\Hom(X^2,X)=1$. Fix nonzero maps $e:X^2\to\one$ and $f:X^2\to X$. Every unit-fixed multiplication on $A$ has the form
\[
 m|_{X^2}=a e+b f.
\]
For a simple extension, $a\ne0$: if $a=0$, its $X$ summand is a proper ideal. 

In strict notation supplied by coherence, put
\[
 U=e\otimes\id_X-\id_X\otimes e\in\Hom(X^3,X).
\]
This is nonzero. If it vanished, choose $c:\one\to X^2$ with $(\id\otimes e)(c\otimes\id)=\id_X$. Tensor the equality $U=0$ on the left by $\id_X$ and precompose with $c\otimes\id_{X^2}$. The result is $\id_{X^2}=ce$, making $X^2$ a retract of the unit. Then $X$ would be invertible, contrary to $\FP(X)=t>1$ in \eqref{A:cat:tdef}.

Projection of associativity on $X^3$ to the $X$ output gives
\[
 aU+b^2V=0,\qquad
 V=f(f\otimes\id)-f(\id\otimes f).
\]
Thus $b\ne0$. If $(a_0,b_0)$ are the constants of the existing lattice multiplication, the same equation and $U\ne0$ give $a/b^2=a_0/b_0^2$. Set $s=b/b_0$. The identities
\[
 a=s^2a_0,\qquad sb=s^2b_0
\]
are exactly the two components of $\varphi m=m_0(\varphi\otimes\varphi)$ for $\varphi=\id_{\one}\oplus s\id_X$. Thus $\varphi$ is an algebra isomorphism to the lattice algebra.

The VOA-extension correspondence transports this isomorphism to an isomorphism fixing $V$ \cite[Theorem 3.2 and Remark 3.3]{B:HKL}. The extension has a two-dimensional weight-zero space; Remark~3.3 supplies the version of the correspondence that allows this.

\section{The simple modules}\label{sec:spectrum}
The next step determines the simple objects of $\D=\D_n$. Write $F:\D\to\Ncat_n$ for the free-module functor. The reconstruction gives
\begin{equation}\label{B:FX}
 F(X)=\one\oplus R,\qquad
 t:=\FP(X)=\frac{n+2+\sqrt{n(n+4)}}2,\qquad \FP(R)=t-1,
\end{equation}
and
\begin{equation}\label{B:dims}
 \begin{gathered}
 \delta:=d(R)=\frac{n-\sqrt{n(n+4)}}2,\qquad d(X)=1+\delta=t^{-1},\\
 \dim\D=n(2+\delta)^2,\qquad \FP(\D)=n(1+t)^2.
 \end{gathered}
\end{equation}
The known objects are $\one,X$ and $Z_x$, $1\leq x\leq(n-1)/2$. The unknown induction multiplicities are first expressed categorically, then bounded using the lattice structure of $X$.

\subsection{Induction multiplicities}
\begin{lemma}\label{B:multiplicities}
Write $Q_i$ for the as-yet-unclassified simples of $\D$, with $i$ in a finite set $I$. Then
\begin{align}
 F(Z_x)&=x\oplus(-x)\oplus R, & F(Q_i)&=b_iR,\\
 \FP(Q_i)&=b_i(t-1), & d(Q_i)&=b_i\delta,
\end{align}
where $b_i$ are positive integers, and
\begin{equation}\label{B:squares}
 \sum_{i\in I}b_i^2=\frac{n+3}{2}=:M.
\end{equation}
Also
\begin{equation}\label{B:XX}
 X^2=\one\oplus X\oplus\bigoplus_{x=1}^{(n-1)/2}Z_x
          \oplus\bigoplus_{i\in I} b_iQ_i,
 \qquad N_{X,Q_i}^{Q_i}=b_i^2-1.
\end{equation}
\end{lemma}
\begin{proof}
Adjunction with $U(g)$ determines every invertible summand of a free module. The known $Z_x$ have exactly the two invertible summands $x,-x$; an unlisted simple has none. Dominance gives $b_i>0$. The dimension formula for restriction of an $A$-module gives $\FP(Z_x)=\FP(A)=1+t$, since its local module has dimension one. Since $\FP(R)=t-1$, this proves the displayed decomposition of $F(Z_x)$. Pivotality gives the spherical dimensions.

Subtract the dimensions of the known simples from $\FP(\D)=n(1+t)^2$. Using $(t-1)^2=nt$ gives
\[
 \sum_i b_i^2=
 \frac{n(1+t)^2-1-t^2-\frac{n-1}{2}(1+t)^2}{(t-1)^2}
 =\frac{n+3}{2}.
\]
Adjunction also gives
$U(R)=X\oplus\bigoplus_x Z_x\oplus\bigoplus_i b_iQ_i$.
Now $UF(X)=X\otimes A$ and $F(X)=\one\oplus R$, so $X^2=\one\oplus U(R)$. Finally
\[
 \dim\Hom_{\D}(Q_iQ_i^*,A)
 =\dim\Hom_{\N}(b_i^2R^2,\one)=b_i^2.
\]
The unit contributes one; Frobenius reciprocity identifies the remaining multiplicity with $N_{X,Q_i}^{Q_i}$.
\end{proof}

\begin{proposition}\label{B:squarefree}
Even without the bimodule theorem below, the inputs imply $b_i^2\mid n+4$. Hence $b_i$ is odd, and
\[
 \rank\D=n+3-\sum_i(b_i^2-1)\equiv n+3\pmod8.
\]
In particular the asserted ranks hold whenever $n+4$ is squarefree.
\end{proposition}
\begin{proof}
For a modular category, the column character $a\mapsto S_{ai}/S_{0i}$ of its integral fusion ring takes algebraic-integer values. Its formal codegree is the sum of the products of its values on $a$ and $a^*$, namely $1/S_{0i}^2$; this is an algebraic integer. This elementary form of the formal-codegree argument is also recorded in \cite{B:Ostrik}.
Here
\[
 S_{0i}^2=\frac{b_i^2\delta^2}{n(2+\delta)^2}
 =\frac{b_i^2}{n+4},
\]
using $\delta^2=n(1+\delta)$ and $(2+\delta)^2=(n+4)(1+\delta)$. Thus $(n+4)/b_i^2$ is a rational algebraic integer. The other conclusions follow from \eqref{B:squares}.
\end{proof}
For general $n$, the remaining task is to exclude $b_i>1$ when $n+4$ has a square divisor. The bimodule $A(X)$ supplies the required bound.

\subsection{The Virasoro structure of the complement}
Use the Fock spaces and blocks of Section~\ref{A:app:lat}. Write
\[
 h_j=nj(j+1),\qquad g_j=h_j+(2j+1)^2=(n+4)j(j+1)+1.
\]
Thus $M_j=M(c,h_j)$, $L_j=L(c,h_j)$, and $K_j=L(c,g_j)$, with the Verma and Fock structures in \eqref{A:lat:Verma}--\eqref{A:lat:Fock}. The charge filtration and its exact block projections are those of Lemmas~\ref{A:blocks:union} and \ref{A:lat:blocklemma}. Let $v_j=e^{(j+1)\beta}$ generate the embedded $M_j\subset B_j$, and let $W_j$ span the bottom of the unique $L_j\subset B_j$.

The orthogonal decomposition $\mathcal U=V\oplus X$ gives
\begin{equation}\label{B:Xblocks}
 X=\bigoplus_{j\geq0}X_j,\qquad
 0\longrightarrow M_j\longrightarrow X_j\longrightarrow K_j\longrightarrow0.
\end{equation}
Indeed the projection of the embedded $M_j$ into $X_j$ is injective, since $M_j$ has no $L_j$ submodule. Lemma~\ref{A:lat:essential} identifies the socle and, by self-duality, the head of $X_j$ with $K_j$.

Choose a bottom vector $x_j$ of $X_j$, a highest vector $\phi_j$ of its socle, and a homogeneous lift $u_j$ of the highest vector of its head. Then
\[
 \wt x_j=h_j,\qquad \wt\phi_j=\wt u_j=g_j,\qquad (u_j,\phi_j)\ne0.
\]
Every vector in $X_j$ is a sum of negative Virasoro descendants of $x_j$ and $u_j$. Since $n$ is odd, every $h_j$ is even and every $g_j$ is odd. This separation will make the reduction triangular.

\subsection{Transition maps between Virasoro blocks}
\begin{lemma}\label{B:couplings}
The vectors can be chosen so that:
\begin{enumerate}[label=(\alph*),leftmargin=7mm]
\item $x_j$ has nonzero lowest-charge component proportional to $e^{-j\beta}$; the projection of $W_{-2nj-1}x_j$ to the bottom of $X_{j+1}$ is nonzero.
\item For $j\geq1$, projection of $Y(W,z)\phi_j$ to $X_{j-1}$ has a nonzero coefficient proportional to $\phi_{j-1}$.
\item The coefficient of $u_{j+1}$ in the head of the $X_{j+1}$ component of $W_{2n-1-(g_{j+1}-g_j)}u_j$ is nonzero.
\end{enumerate}
\end{lemma}
\begin{proof}
The charge of $L_1^\lambda$ never decreases. Invariance expresses the pairing of two lattice states through the weight-zero functional, which is supported in charges zero and one. Hence $(v_j,v_j)=0$, since both inputs have charge $j+1$ and all further charge shifts are nonnegative. The bottom form of $B_j$ is nondegenerate, while the $W_j$ line is nondegenerate. Its orthogonal projection of $v_j$ into $X_j$ therefore subtracts a \emph{nonzero} multiple of $W_j$: otherwise a nonzero vector in the one-dimensional nondegenerate orthogonal complement would be isotropic. This proves the assertion about $x_j$.

Normalize $W=W_1$ to have lowest component $e^{-\beta}$. The lowest-charge contribution to $Y(W,z)x_j$ is a nonzero scalar times
\[
 Y(e^{-\beta},z)e^{-j\beta}
 =\epsilon_j z^{2nj}\bigl(e^{-(j+1)\beta}+\text{higher oscillator terms}\bigr),
 \quad\epsilon_j\ne0.
\]
Block projection to $B_{j+1}$ preserves this coefficient: on the charge $-(j+1)$ associated graded it is the identity. The action of $V$ preserves $X$. At weight $h_{j+1}$ its projection is thus a nonzero multiple of $x_{j+1}$. This proves (a).

The lowest-charge component $\phi_j^0$ of $\phi_j$ is the nonzero singular vector in $\mathsf F_{j+1}=M_j$. The lowest-charge part of $Y(W,z)\phi_j$ is
\[
 Y(e^{-\beta},z)\phi_j^0\in\mathsf F_j\{z\}.
\]
It is nonzero: the lattice exponential is a charge shift composed with invertible creation and annihilation exponentials. Its projection to $L_{j-1}$ is zero. A nonzero projected Virasoro intertwiner would, by contragredient symmetry, give an intertwiner of type $\binom{K_j}{L_1\ L_{j-1}}$. Its image would be $C_1$-cofinite by \cite[Theorem 4.1.7]{A:CJORY}, contrary to irreducibility and non-$C_1$-cofiniteness of $K_j$. Thus the nonzero series takes values in the submodule $K_{j-1}$ of $\mathsf F_j$.

Its leading coefficient is a highest vector: the Virasoro commutator formula with the primary $e^{-\beta}$ shows that every positive $L_r$ annihilates the first nonzero coefficient on the highest input $\phi_j^0$. Since $K_{j-1}$ is an irreducible Verma module, that coefficient is at its bottom $g_{j-1}$, not at a higher singular weight. Exact block projection preserves its charge-$j$ component. The filtered part of $B_{j-1}$ of charges at least $j$ is precisely its embedded Verma module; the next filtered part is zero. Hence the projected coefficient is a nonzero multiple of the socle vector $\phi_{j-1}$, proving (b).

For an even-weight primary, invariance gives $(W_r a,b)=(a,W_{4n-2-r}b)$. Apply this to (b), pairing with $u_j$ and $\phi_{j+1}$. The adjoint mode has weight increase $g_{j+1}-g_j$, and the pairing is nonzero. Since pairing with the socle detects the head quotient, this is precisely (c).
\end{proof}

\subsection{The two-generator bound}\label{B:sec:twogen}
For a $V$-module $M$, let $A(M)=M/O_V(M)$ be the Frenkel--Zhu bimodule. Put $x_l=[T]*-$ and $x_r=-*[T]$. For homogeneous $v$ the standard relations are
\begin{equation}\label{B:ZhuVir}
 [L_{-r}v]=(-1)^r(rx_r-x_l+\wt v)[v]\quad(r\geq1).
\end{equation}
Also the higher circle relations reduce $[W_{-s}v]$, $s\geq2$, to classes of $W_rv$ with $r\geq-1$, and
\begin{equation}\label{B:ZhuW}
 [W]*[v]=[W_{-1}v]+\sum_{a=1}^{2n}\binom{2n}{a}[W_{a-1}v].
\end{equation}
These follow directly from the defining residue relations; they hold without assuming $v$ is Virasoro highest.

\begin{theorem}\label{B:twogen}
The bimodule $A(X)$ is generated over $A(V)$ on the left and right by the two classes $[x_0]$ and $[u_0]$, of weights zero and one. Consequently
\begin{equation}\label{B:fusionbound}
 N_{X,S}^{T}\leq2
\end{equation}
for any simple $V$-modules $S,T$.
\end{theorem}
\begin{proof}
By \eqref{B:Xblocks} and \eqref{B:ZhuVir}, it is enough to reduce $[x_j],[u_j]$. Induct on their conformal weights. The two bottom generators are $x_0,u_0$. Their weights are distinct from all subsequent generator weights.

For $x_{j+1}$ use Lemma~\ref{B:couplings}(a). The target weight increase from $x_j$ is $2n(j+1)$. For $j=0$, the mode is $W_{-1}$, which is handled by \eqref{B:ZhuW}. For $j\geq1$ the mode is at most $-2$, and a higher circle relation replaces it by modes of resulting weight at most $h_j+2n<h_{j+1}$. All other block components of the original coefficient are descendants of generators of smaller weight. There is no other generator of weight $h_{j+1}$: the $h_l$ are strictly increasing even integers and the $g_l$ are odd. The coefficient of $x_{j+1}$ is nonzero, so this solves for its class.

For $u_{j+1}$ use Lemma~\ref{B:couplings}(c). Its weight increase is $2(n+4)(j+1)>2n$, so a higher circle relation reduces the relevant mode to states of smaller weight. After removing the nonzero head coefficient, the remaining part of its block lies in the Verma module generated by $x_{j+1}$, whose weight is strictly smaller than $g_{j+1}$. Other blocks have smaller generator weights as well, because no $h_l$ equals $g_{j+1}$ and the $g_l$ are strictly increasing. Thus induction also solves for $[u_{j+1}]$. This proves the generation assertion.

Strong generation by $T,W$ implies that $A(V)$ is generated by $[T],[W]$. Since $[T]$ is central it is commutative. Its simple modules, hence all simple $V$-module top spaces, are one-dimensional. The Frenkel--Zhu--Li injection \cite[Theorem 2.12]{A:AA} bounds an intertwining-operator space by the contraction of $A(X)$ with the two top-space characters. That contraction is generated by the images of the two displayed classes and has dimension at most two. This proves \eqref{B:fusionbound}.
\end{proof}

\begin{corollary}\label{B:rankthm}
Every $b_i$ in Lemma~\ref{B:multiplicities} is one. There are exactly $M=(n+3)/2$ new simples $Q_i$, and the two ranks in Theorem~\ref{A:cat:main} hold.
\end{corollary}
\begin{proof}
Apply \eqref{B:fusionbound} with $S=T=Q_i$. Equation \eqref{B:XX} gives $b_i^2-1\leq2$, so the positive integer $b_i$ is one. Now use \eqref{B:squares} and the semion factorization. Equation \eqref{B:XX} becomes $X^2=\bigoplus_{S\in\Irr(\D)}S$.
\end{proof}
In particular the missing objects have an intrinsic construction: split the complementary summand
\[
 Q=X^2\ominus\left(\one\oplus X\oplus\bigoplus_x Z_x\right).
\]
Its endomorphism algebra is $\C^M$, and its primitive idempotents give exactly the $Q_i$, each once. Their semion translates give the other $M$ missing $V$-modules.

Every simple is self-contragredient. On the contragredient top, $o(T)$ has the same eigenvalue, and so does $o(W)$ because $W$ is primary of even weight. The two Zhu characters coincide; Zhu's bijection gives the isomorphism. Rationality makes $A(V)$ semisimple, so its commutativity and the rank theorem give $A(V)\cong\C^{2n+6}$ as abstract algebras.

\section{Modular data}\label{sec:modular}
Order the simples of $\D_n$ as $\one,X,Z_x,Q_y$, with $1\leq x\leq(n-1)/2$ and $1\leq y\leq(m-1)/2$. The next formulas state that, for $u=n^{-1/2}$ and $v=m^{-1/2}$,
\begin{equation}\label{B:Sfull}
 S=\begin{pmatrix}
 (u+v)/2&(u-v)/2&u\mathbf1^{\!t}&-v\mathbf1^{\!t}\\
 (u-v)/2&(u+v)/2&u\mathbf1^{\!t}&v\mathbf1^{\!t}\\
 u\mathbf1&u\mathbf1&C_n&0\\
 -v\mathbf1&v\mathbf1&0&C_m
 \end{pmatrix},\qquad
 u=n^{-1/2},\quad v=m^{-1/2},
\end{equation}
\begin{equation}\label{B:Tfull}
 (C_r)_{ab}=\frac{2}{\sqrt r}\cos\frac{4\pi ab}{r},\qquad
 \Theta=\operatorname{diag}\left(1,1;
       e^{2\pi i x^2/n};e^{2\pi i y^2/m}\right).
\end{equation}

The matrix $\Theta$ records the ribbon twists. Use the positive square root of $\dim\D_n$ in the normalization of $S$. The end of this section recovers the character normalization and compares the data with the formulas of~\cite{A:EG,B:EG11}.

\subsection{The known blocks}
Normalize categorical $S$ by the positive root $\sqrt{\dim\D}=\sqrt n(2+\delta)$. Let $B$ be the restriction matrix from the local $A$-modules. Its only nonzero entries are
\[
 B_{\one,0}=B_{X,0}=1,\qquad B_{Z_x,x}=B_{Z_x,-x}=1.
\]
Kirillov--Ostrik's local-module modular-invariance theorem, which does not assume unitarity, gives
\begin{equation}\label{B:branchS}
 S_{\D}B=B S_{\Pcat_n},\qquad
 (S_{\Pcat_n})_{xy}=n^{-1/2}e^{4\pi ixy/n}
\end{equation}
\cite[Theorems 4.1 and 4.5]{A:KO}. Twists on restriction agree. Since $d(Z_x)=2+\delta$ and $d(Q_i)=\delta$, the vacuum row is
\[
 S_{00}=(u+v)/2,\quad S_{0X}=(u-v)/2,\quad
 S_{0Z_x}=u,\quad S_{0Q_i}=-v.
\]
Taking each column of \eqref{B:branchS}, and using symmetry, now gives exactly \eqref{B:Sfull} except for its last block, temporarily denoted $K=(K_{ij})$. The known twists are $\theta_X=1$ and $\theta_{Z_x}=e^{2\pi ix^2/n}$.

All fusion matrices are real symmetric because all simples are self-dual. Their eigenvalues $S_{ai}/S_{0i}$ are real; hence $S$ is real. Thus $S^2=1$, and its remaining blocks satisfy
\begin{equation}\label{B:Krelations}
 K\mathbf1=-\frac1{\sqrt m}\mathbf1,\qquad
 K^2=I_M-\frac2m\mathbf1\mathbf1^t.
\end{equation}
The normalized Gauss sum of $\D$ equals that of its local pointed category \cite[Theorem 4.5]{A:KO}:
\[
 \gamma=\frac1{\sqrt n}\sum_{x\bmod n}e^{2\pi ix^2/n}
 =\begin{cases}1&n\equiv1\pmod4,\\ i&n\equiv3\pmod4.\end{cases}
\]
If $w_i=\theta_{Q_i}$, the vacuum entry of the Gauss-sum identity gives
\begin{equation}\label{B:gaussmissing}
 1+2\sum_{i=1}^{M}w_i=\gamma\sqrt m.
\end{equation}
For example, expand $\sum_a S_{0a}^2\theta_a=\gamma S_{00}$ and use the ordinary quadratic Gauss sum for the $Z$ sector. This proves \eqref{B:gaussmissing} without choosing a proposed twist for any $Q_i$.

\subsection{The integral fusion algebra}\label{B:sec:metric}
We recover the unknown block by adapting the characteristic-two method of \cite[Section~5.2]{B:GSY}; this also treats square $m$.

Write $N_{ij}^k=N_{Q_i,Q_j}^{Q_k}$. Adjunction gives
\begin{equation}\label{B:QQ}
 Q_iQ_j=\delta_{ij}\one+(1-\delta_{ij})X+\sum_x Z_x+\sum_k N_{ij}^kQ_k,
 \qquad \sum_kN_{ij}^k=M-2+\delta_{ij}.
\end{equation}
Indeed the coefficient of $Z_x=U(x)$ is the multiplicity of $x$ in $R^2$, the sum of the unit and $X$ coefficients is one, and the last equality follows by applying $F$ and comparing $R$ coefficients. Similarly
\[
 XQ_i=X+\sum_x Z_x+\sum_{j\ne i}Q_j.
\]
In the complexified fusion ring define $c_i=\one+X-Q_i$. Direct substitution using $X^2=\sum_a a$ gives
\begin{equation}\label{B:cc}
 c_ic_j=2\delta_{ij}\one+\sum_k a_{ij}^kc_k,\qquad
 a_{ij}^k=1-N_{ij}^k\leq1,\qquad \sum_k a_{ij}^k=2-\delta_{ij}.
\end{equation}
The elements $\one,c_1,\ldots,c_M$ are linearly independent in the fusion ring.

Evaluate at the fusion-ring character belonging to $Q_j$. The evaluations are
\[
 x_i(j)=\frac{S_{0j}+S_{Xj}-S_{Q_i j}}{S_{0j}}=\sqrt m K_{ij}.
\]
They form a basis of functions on the $M$ labels, since $K$ is invertible. Equations \eqref{B:Krelations} imply
\begin{equation}\label{B:tracepair}
 \sum_i x_i=-1,\qquad
 \operatorname{Tr}(x_ix_j)=m\delta_{ij}-2.
\end{equation}
Here the trace is the sum over all $M$ evaluations. The lattice $\mathcal O=\sum_i\Z x_i$ is a unital integral algebra by \eqref{B:cc} and $1=-\sum_i x_i$. All its evaluations are algebraic integers, and its trace-pairing determinant is $m^{M-1}$, odd.

Modular Galois symmetry \cite[Theorem 10.1]{B:ENO} permutes the $Q_i$ among themselves. Their vacuum entry has Galois conjugates $\pm1/\sqrt m$. The other vacuum entries have different absolute values, since
\[
 0<\frac{u-v}{2}<v<\frac{u+v}{2}<u
 \qquad(n\geq3).
\] The Galois signs on all $Q_i$ are the same, namely $\sqrt m/\sigma(\sqrt m)$. Therefore
\[
 \sigma(x_i(j))=x_{\widehat\sigma(i)}(j).
\]
Choose a prime above two in a cyclotomic field containing the entries, and an automorphism in its decomposition group inducing residue-field Frobenius. Such a lift exists even if the cyclotomic field ramifies at two. Let its permutation on $I$ be $\pi$. Each evaluation of $x_i^2-x_{\pi(i)}$ vanishes modulo that prime. Hence its trace after multiplication by any $x_j$ is an even rational integer. The invertibility modulo two of the trace matrix \eqref{B:tracepair} proves
\begin{equation}\label{B:frobenius}
 x_i^2\equiv x_{\pi(i)}\pmod{2\mathcal O}.
\end{equation}

Write $x_i^2=\sum_k b_kx_k$. By \eqref{B:cc}, $b_k=a_{ii}^k-2\leq-1$ and $\sum_kb_k=1-2M$. Exactly the coefficient at $\pi(i)$ is odd by \eqref{B:frobenius}. Thus that coefficient is at most $-1$ and every other coefficient is at most $-2$. Their sum forces equality in every place. Lifting back to the independent elements $\one,c_k$ yields
\begin{equation}\label{B:twovalued}
 c_i^2=2\one+c_{\pi(i)}.
\end{equation}
For $i\ne j$, compare the coefficient of $\one$ in $(c_ic_j)^2=c_i^2c_j^2$. The permutation property of $\pi$ and \eqref{B:cc} give
\[
 \sum_k(a_{ij}^k)^2=2=\sum_k a_{ij}^k.
\]
Since $a(a-1)\geq0$ for every integer $a$, exactly two coefficients are one and all the others are zero. Thus $c_ic_j=c_p+c_q$ with $p\ne q$.

\subsection{The finite abelian group}
Put $c_0=2\one$. The preceding products define a commutative two-valued group on $\{0\}\cup I$: multiplication lists the two labels in a product $c_ac_b$. Associativity follows from associativity of the algebra and independence of the basis. Its identity is strong; every label is its own unique inverse. There is no nonidentity label whose square is $\{0,0\}$, by \eqref{B:twovalued}.

The following precise consequence is needed from the classification of involutive commutative two-valued groups \cite[Theorem 1.10]{B:BVG}: a finite such group with no nonidentity strong involution is $H/\{\pm1\}$ for a finite abelian group $H$ of odd order. Indeed the unipotent and special families, and any nontrivial Boolean factor in that classification, contain a nonidentity element squaring to the identity twice. In the principal family the same exclusion removes $H[2]$. Therefore
\begin{equation}\label{B:Hgroup}
 |H|=2M+1=m,\qquad
 c_{[h]}\longmapsto e_h+e_{-h}
\end{equation}
identifies the algebra with $\C[H]^{\{\pm1\}}$.

The column evaluations are consequently the characters of this algebra: for each $[k]\in H/\{\pm1\}$ there is a unique class $[\chi_k]\in\widehat H/\{\pm1\}$ such that
\[
 c_{[h]}([k])=\chi_k(h)+\chi_k(h)^{-1}.
\]
The identity evaluation sends every $c_{[h]}$ to two. Symmetry of $K$ and the product formula in \eqref{B:Hgroup} show that $[k]\mapsto[\chi_k]$ preserves the two-valued product.

Use the lifting fact proved in \cite[Lemma 5.4]{B:GSY}; the following argument records it in the present notation. A two-valued isomorphism $f:H/\{\pm1\}\to H'/\{\pm1\}$ for odd finite abelian groups lifts to a group isomorphism, uniquely up to a global sign. Choose $g\ne0$ and a representative $b$ of $f([g])$. For each $x$ choose the unique representative $F(x)$ of $f([x])$ for which
\[
 f([x+g])=[F(x)+b].
\]
Existence follows from the two-valued product, and uniqueness from oddness: $[u+b]=[-u+b]$ forces $u=0$ or $b=0$. Induction using the product with $[g]$ gives $f([x+rg])=[F(x)+rb]$. Apply the product rule to $x+rg$ and $y+sg$. The variables $r+s,r-s$ vary independently modulo the odd order of $b$. If the first outputs could not be matched, the sign classes $[F(x)-F(y)+(r-s)b]$ would all be constant, impossible for a cyclic orbit of odd length at least three. Thus $[F(x)+F(y)+tb]=[F(x+y)+tb]$ for every $t$. Taking $t=0,1$ and using the same uniqueness proves additivity. Bijectivity follows from that of $f$.

This gives a nondegenerate bicharacter $\beta:H\times H\to\C^\times$ with
\begin{equation}\label{B:Kbeta}
 K_{[h],[k]}=\frac{\beta(h,k)+\beta(h,k)^{-1}}{\sqrt m}.
\end{equation}
Symmetry says $\beta(h,-)+\beta(h,-)^{-1}=\beta(-,h)+\beta(-,h)^{-1}$. Independence of characters and the fact that a group cannot be the union of two proper subgroups show that $\beta$ is globally symmetric or globally skew. The twist argument next proves symmetry, including when $m$ is a square.

\subsection{The quadratic form}
Extend $w_i$ to an even function $w:H\to\C^\times$ by $w(0)=1$, and put $v(h)=w(h)^{-1}$. Vafa's theorem gives $|w(h)|=1$. The identity
\[
 S\Theta S=\gamma\Theta^{-1}S\Theta^{-1}
\]
in its vacuum--$Q$ and $Q$--$Q$ blocks, using \eqref{B:Kbeta}, gives respectively
\begin{align}
 \sum_{h\in H}w(h)\beta(h,k)&=\gamma\sqrt m\,v(k),\label{B:Fourierw}\\
 v(h+k)+v(h-k)&=v(h)v(k)\bigl(\beta(h,k)+\beta(h,k)^{-1}\bigr).
 \label{B:parsum}
\end{align}
Here is the block calculation for the second equation. The two distinguished columns contribute $2/m$. Extending the sum in $K\operatorname{diag}(w)K$ from sign classes to all of $H$ introduces the $h=0$ term $-2/m$, cancelling it. Its four character sums are the two sums in \eqref{B:Fourierw} at $h+k,h-k$, each twice. This works for a symmetric or a skew bicharacter, since the arguments occur in an even pair. Cancelling $\gamma/\sqrt m$ yields \eqref{B:parsum}.

The real number $\beta(h,k)+\beta(h,k)^{-1}$ is nonzero, because $\beta(h,k)$ has odd order. For two unit complex numbers with nonzero sum $z$, their product is $z/\overline z$. Applying this to \eqref{B:parsum} gives
\begin{equation}\label{B:parallelogram}
 v(h+k)v(h-k)=v(h)^2v(k)^2.
\end{equation}
It follows inductively that $v(rh)=v(h)^{r^2}$. If $h$ has odd order $d$, compare $r=(d+1)/2$ and $(d-1)/2$ using evenness; this gives $v(h)^d=1$. Write $v(h)=e^{2\pi i Q(h)/e}$, with $e$ the odd exponent of $H$. Equation \eqref{B:parallelogram} is the parallelogram identity in $\Z/e$. Division by two is defined in $\Z/e$. Set
\[
 P(h,k)=\frac{Q(h+k)-Q(h-k)}2
       =Q(h+k)-Q(h)-Q(k).
\]
Applying the parallelogram identity at $h+k$ and $h-k$ gives
$P(h+l,k)+P(h-l,k)=2P(h,k)$. Taking $l=h$ gives $P(2h,k)=2P(h,k)$, and then taking $h=(a+b)/2$, $l=(a-b)/2$ proves additivity in the first variable. Symmetry gives additivity in the second.

Thus $b_v(h,k)=v(h+k)/(v(h)v(k))$ is a symmetric bicharacter. Equation \eqref{B:parsum} says $b_v+b_v^{-1}=\beta+\beta^{-1}$. Character independence and the same two-subgroup argument give $b_v=\beta$ or $\beta^{-1}$ globally. In particular $\beta$ is symmetric and $b_v$ is nondegenerate. Hence $w$ is a nondegenerate homogeneous quadratic form on $H$ and
\begin{equation}\label{B:metricconclusion}
 K_{[h],[k]}=\frac{b_w(h,k)+b_w(h,k)^{-1}}{\sqrt m},\qquad
 \sum_{h\in H}w(h)=\gamma\sqrt m.
\end{equation}
Thus the missing block comes from a metric group of order $m$.

\subsection{Identification of the metric group}
\begin{lemma}\label{B:primitive}
Some simple $V$-module has twist $e^{2\pi i/m}$.
\end{lemma}
\begin{proof}
By (H5), $\ch X-\ch V=\Lambda_{m,m}$. Its modular transform is the exact Poisson-summation identity
\begin{equation}\label{B:thetaS}
 \Lambda_{m,m}(-1/\tau)=\frac1{\sqrt{2m}}
             \sum_{a\bmod2m}(-1)^a\Lambda_{a,m}(\tau).
\end{equation}
Let $D$ act on functions by $Df(\tau)=e^{2\pi ic/24}f(\tau+1)$, and put $\xi=e^{2\pi i/m}$. By (H1), the span of simple characters is modular invariant, and $D$ acts there with the ribbon twists as eigenvalues. Choose $L$ divisible by their orders and by $4m$. The polynomial operator
\[
 \Pi_\xi=\frac1L\sum_{r=0}^{L-1}\xi^{-r}D^r
\]
projects to the $\xi$-eigenspace both on the character span and on the theta functions in \eqref{B:thetaS}. The terms with $a=m-2$ and $a=m+2$ survive, since
\[
 \frac{(m-2)^2}{4m}-\frac n4=\frac1m.
\]
Any surviving index satisfies $a^2-mn\equiv4\pmod{4m}$, so $a^2\equiv mn\equiv1\pmod4$ and $a$ is odd. All surviving terms therefore have negative sign and nonnegative theta coefficients; they cannot cancel. Thus $\Pi_\xi\ne0$ on the character span, proving the assertion.
\end{proof}

\begin{proposition}\label{B:cyclicmetric}
The metric group in \eqref{B:metricconclusion} is isometric to
\[
 (\Z/m,\ y\longmapsto e^{2\pi iy^2/m}).
\]
\end{proposition}
\begin{proof}
All twists in $\D$ have odd order: this holds for its known objects and, by \eqref{B:metricconclusion}, for its other objects. Multiplication by $\theta_J=\pm i$ produces a twist with a factor four in its order. Therefore the simple in Lemma~\ref{B:primitive} belongs to $\D$, not its $J$-translate. It is neither $\one,X$ nor a $Z_x$, since $\gcd(n,m)=1$. It corresponds to some $h\in H$ with $w(h)=e^{2\pi i/m}$.

For a homogeneous quadratic form on an odd group, $w(h)$ has order dividing the order of $h$, as proved above. Thus $h$ has order $m=|H|$. It generates $H$, and homogeneity gives $w(rh)=e^{2\pi ir^2/m}$. Equation \eqref{B:metricconclusion} now gives precisely the last block in \eqref{B:Sfull} and the last twists in \eqref{B:Tfull}.
\end{proof}


\subsection{Comparison with earlier modular data}\label{B:sec:galois}
To specify the Galois conjugation, write the cyclic factor of~\cite{A:EG} using units $A\pmod n$, $B\pmod m$ with
\begin{equation}\label{B:EGgauss}
 \Gamma_n(A)\Gamma_m(B)=-1,\qquad
 \Gamma_r(a)=\frac1{\sqrt r}\sum_{x\bmod r}e^{2\pi iax^2/r}.
\end{equation}
In the positive-dimension normalization, its distinguished $2\times2$ block is
$\frac12\left(\begin{smallmatrix}u-v&u+v\\u+v&u-v\end{smallmatrix}\right)$;
its vacuum entries on the two further families are $u,v$; its $X$ entries there are $u,-v$; its remaining diagonal blocks are
\[
 \frac2{\sqrt n}\cos\frac{4\pi Axx'}n,\qquad
 -\frac2{\sqrt m}\cos\frac{4\pi Byy'}m.
\]
Its twists are $1,1;e^{-2\pi iAx^2/n};e^{2\pi iBy^2/m}$.

The comparison uses the full polarization of each quadratic form. In particular, for $q_B(y)=e^{2\pi iBy^2/m}$ its bicharacter is
\[
b_B(y,y')=\frac{q_B(y+y')}{q_B(y)q_B(y')}=e^{4\pi iByy'/m},
\]
and the primed block is $-m^{-1/2}(b_B+b_B^{-1})$. This normalization is compatible with the theta-character formulas of Section~\ref{B:sec:characters}. The comparison below refers to these explicitly displayed matrices; using the half-polarization $e^{2\pi iByy'/m}$ in that block would give a different matrix. The derivation of \eqref{B:Sfull}--\eqref{B:Tfull} above fixes the normalization without depending on that comparison.

These are the centralizing labels $a_0,b_0,c_{x,-x},d_{0,y}$ in \cite[Proposition 7(b), Section 4.4]{A:EG}. This statement concerns the explicitly defined modular data, not the separate conjecture identifying all near-group doubles with those data.

Choose an integer $\ell$ by the Chinese remainder theorem with
\begin{equation}\label{B:CRTgeneral}
 \ell\equiv1\pmod4,\qquad
 \ell\equiv-A^{-1}\pmod n,\qquad
 \ell\equiv B^{-1}\pmod m.
\end{equation}
It is a unit modulo $4nm$. Let $\sigma_\ell(\zeta_{4nm})=\zeta_{4nm}^{\ell}$. For odd $r$ and $\ell\equiv1\pmod4$, quadratic Gauss sums give
$\sigma_\ell(\sqrt r)=(\ell/r)\sqrt r$.
Put $s_n=(\ell/n)$ and $s_m=(\ell/m)$. Since $n\equiv m\pmod4$, \eqref{B:EGgauss} gives
\[
 s_ns_m=(-1/n)(A/n)(B/m)=-1.
\]
Thus $s_m=-s_n$. Substituting \eqref{B:CRTgeneral} into every entry gives
\begin{equation}\label{B:EGcomparison}
 \sigma_\ell(S^{EG})=s_n S,\qquad
 \sigma_\ell(\Theta^{EG})=\Theta.
\end{equation}
The single overall sign is exactly the permitted choice of square root of global dimension when normalizing $S$. This proves the prescribed Galois relationship for any admissible cyclic pair in the family $D_n$ of~\cite{A:EG}; it is an entrywise identity valid for composite odd parameters as well.

\subsection{The twisted Haagerup--Izumi comparison}\label{B:sec:HIcomparison}
Let $n=N^2$ and $\mu=N^2+4$. Compare with the explicit data $D^1\mathrm{Hg}_N$ of \cite[Section 3.2, equation (3.3)]{B:EG11}. Their two groups of nontrivial labels are the nonzero sign classes of a metric group of order $N^2$, and the nonzero sign classes of $\Z/\mu$.

Here is an explicit description of the first metric group, so that no classification of HI categories enters. For twist parameter $\omega$, use generators $g,h$ with
\[
 Nh=0,\qquad Ng=2\omega h,\qquad
 q(ag+bh)=\exp\!\left(2\pi i\frac{\omega a^2+Nab}{N^2}\right).
\]
Changing $a$ by $N$ and $b$ by $-2\omega$ preserves this form. For a unit $\omega\pmod N$, $g$ has order $N^2$, and the form is cyclic with coefficient $\omega$. The labels $a_i$ in \cite{B:EG11} are the nonzero $bh$ sign classes, and its $c_{a,b}$ are the classes with $a\ne0$. This gives every nonzero class exactly once. For $\omega=1$ the form is $q(x)=e^{2\pi ix^2/N^2}$. The other coefficient in \cite{B:EG11} is
$B=(\mu-1)/2\equiv-2^{-1}\pmod\mu$; its cosine block can equally be written $-2\mu^{-1/2}\cos(4\pi Byy'/\mu)$.

Choose
\begin{equation}\label{B:CRTHI}
 \ell\equiv1\pmod{4N^2},\qquad \ell\equiv-2\pmod\mu.
\end{equation}
Since $\mu\equiv5\pmod8$, $(-2/\mu)=-1$. Hence $\sigma_\ell$ fixes $N$ and sends $\sqrt\mu$ to $-\sqrt\mu$. It fixes the first quadratic form and sends the second coefficient $B$ to one. Direct substitution in the data displayed in~\cite{B:EG11} gives
\[
 \sigma_\ell\bigl(D^1\mathrm{Hg}_N\bigr)=(S,\Theta)
\]
with the positive square-root normalization of this paper. Combined with the ribbon equivalence \eqref{A:cat:realization}, this proves the modular-data assertion for the center of the constructed twisted HI category.


\subsection{The full VOA matrices}
In the order $\Irr(\D)$ followed by their $J$-translates, set
\[
 S_{\Scat}=\frac1{\sqrt2}\begin{pmatrix}1&-1\\-1&-1\end{pmatrix},
 \qquad \Theta_{\Scat}=\operatorname{diag}(1,i^{-n}).
\]
The categorical matrices are their tensor products with \eqref{B:Sfull}--\eqref{B:Tfull}, up to the global normalization of $S$. The genus-one/Verlinde comparison identifies the character and categorical matrices up to the common normalization of $S$ \cite{B:Huang}. The character normalization is
\begin{equation}\label{B:VOAmatrices}
 S^V=-S_{\Scat}\otimes S,\qquad
 T^V=e^{-2\pi i(1-6n)/24}(\Theta_{\Scat}\otimes\Theta).
\end{equation}
Indeed Poisson summation in (H5), using the unique minimum-weight simple, gives
\[
 S^V_{V,Y}=\tfrac12\big((2n)^{-1/2}-(2m)^{-1/2}\big)>0.
\]
Since $d(Y)=-d(X)<0$, this forces the negative normalization of $S^V_{V,V}$ and hence the minus sign in \eqref{B:VOAmatrices}. This sign distinction matters for this nonunitary ribbon structure.

\section{Characters and Zhu's algebra}\label{B:sec:characters}
The complete module list is now known. The character determination uses the space of holomorphic modular forms of weight $1/2$ for the two Weil representations appearing in the modular data. This also treats composite parameters, where a single character need not generate the full modular representation.

\subsection{Splitting the character vector into two Weil representations}
For $r=n,m$ put
\[
 \vartheta_{a,r}=\eta\Lambda_{a,r}
 =\sum_{t\in\Z}q^{(2rt+a)^2/(4r)},\qquad 0\leq a\leq r.
\]
Let $\rho_r^+$ be the even part of the Weil representation of
$(\Z/2r,a\mapsto a^2/(4r))$. Use component coordinates indexed by sign classes $0,1,\ldots,r$, so its theta transformation is
\begin{align}
 \vartheta_{a,r}(-1/\tau)&=\frac{\sqrt{-i\tau}}{\sqrt{2r}}
 \left(\vartheta_{0,r}+(-1)^a\vartheta_{r,r}
       +2\sum_{b=1}^{r-1}\cos(\pi ab/r)\vartheta_{b,r}\right),\label{B:WeilS}\\
 \vartheta_{a,r}(\tau+1)&=e^{2\pi ia^2/(4r)}\vartheta_{a,r}(\tau).
 \label{B:WeilT}
\end{align}
The factor $\sqrt{-i\tau}$ includes the usual $e^{-\pi i/4}$ in the metaplectic $S$ matrix.

Temporarily label the $Z_x,Q_y$ and their $J$-translates by $A_a,B_b$ using
\[
 A_{n-2x}=Z_x,\quad A_{2x}=JZ_x,\qquad
 B_{m-2y}=Q_y,\quad B_{2y}=JQ_y.
\]
All $1\leq a<n$ and $1\leq b<m$ occur once. These are categorical labels, not yet character assertions. Define vectors $f^{(n)},f^{(m)}$ by
\begin{align*}
 f^{(n)}_n&=\eta(\ch V+\ch X),& f^{(n)}_0&=\eta(\ch J+\ch Y),&
 f^{(n)}_a&=\eta\ch A_a,\\
 f^{(m)}_m&=\eta(\ch X-\ch V),& f^{(m)}_0&=\eta(\ch Y-\ch J),&
 f^{(m)}_b&=\eta\ch B_b.
\end{align*}
Substitution of \eqref{B:Sfull}, \eqref{B:Tfull}, and \eqref{B:VOAmatrices} gives precisely \eqref{B:WeilS}--\eqref{B:WeilT} for these two vectors. To check the signs, for example,
\[
 \cos\frac{\pi(r-2x)(r-2y)}r=-\cos\frac{4\pi xy}r
 \quad(r\text{ odd}),
\]
and the factors at $0,r$ account for the half-sums defining $V,X,J,Y$.
For the $T$ check, the exponent of $\eta\ch M$ is $h(M)+n/4$ modulo $\Z$.
For $B_{m-2y}$ this is $y^2/m+n/4\equiv(m-2y)^2/(4m)$, and the other three label types follow in the same way.

By (H5) every lowest weight is at least $-n/4$. Thus each component of these vectors has no negative Fourier exponent. Modular covariance then gives holomorphy at every cusp. Consequently
\begin{equation}\label{B:holomorphicWeil}
 f^{(r)}\in M_{1/2}(\rho_r^+),\qquad r=n,m.
\end{equation}

\subsection{The unary-theta basis}
The argument uses Bruinier--Schwagenscheidt's unary-theta basis theorem
\cite[Lemma 2.1]{B:BS}, a formulation of Skoruppa's critical-weight result:
for a positive integer $r$, $M_{1/2}(\rho_r)$ is spanned by the unary theta vectors at indices $r/d^2$, lifted from isotropic subgroups and acted on by discriminant-form automorphisms, where $d^2\mid r$. The involution $a\mapsto-a$ acts trivially at weight $1/2$, so this is also the space in \eqref{B:holomorphicWeil}. The required explicit consequences are verified next.

On the full coordinates $a\in\Z/2r$, for odd $r$ and $d^2\mid r$, define
\[
 (E_d)_{ab}=
 \begin{cases}1&d\mid a,\ d\mid b,\quad b\equiv a\pmod{2r/d},\\
 0&\text{otherwise}.
 \end{cases}
\]
Then $(E_d\vartheta_r)_a$ is zero unless $d\mid a$, and otherwise is
$\vartheta_{a/d,r/d^2}$. Indeed the $d$ residue classes $b=a+2rs/d$ partition the integers $d\ell$ with $\ell\equiv a/d\pmod{2r/d^2}$. This is exactly the isotropic lift in the basis theorem. An automorphism at index $r/d^2$ is multiplication by signs at its prime-power factors. Lift those signs to the corresponding prime powers of $r$, choosing either sign at any remaining factors. The lifted permutation $P_\sigma$ has square one, commutes with $E_d$, and gives the same lifted theta vector. Hence
\begin{equation}\label{B:thetaM}
 f^{(r)}=M_r\vartheta_r,\qquad
 M_r\in\operatorname{span}_{\C}\{E_dP_\sigma:d^2\mid r\}.
\end{equation}
All these matrices commute with the Weil matrices. For $T$, this follows from
$b^2\equiv a^2\pmod{4r}$ on the support of $E_d$. For $S$, summing the finite Fourier kernel over a coset of the isotropic subgroup proves the equality; $P_\sigma$ preserves the quadratic form. Moreover they are symmetric on the full coordinates. On sign-class coordinates this becomes
\begin{equation}\label{B:weightedtranspose}
 D_rM_r=M_r^tD_r,\qquad
 D_r=\operatorname{diag}(1,2,\ldots,2,1).
\end{equation}
These statements hold for complex linear combinations; no reality assumption on their coefficients is needed.

The functions $\vartheta_{a,r}$, $0\leq a\leq r$, are linearly independent. Their exponents are disjoint: equality of two squared integers implies equality of their sign classes modulo $2r$. Thus every scalar coefficient of $M_r$ can be recovered from Fourier coefficients.

It is already known that $f^{(r)}_r=\vartheta_{r,r}$ from (H5). Also
\[
 f^{(n)}_0+f^{(m)}_0=2\eta\ch Y=\vartheta_{0,n}+\vartheta_{0,m}.
\]
By the basis theorem, $f^{(r)}_0$ is a linear combination of
$\vartheta_{0,r/d^2}$, $d^2\mid r$: every automorphism fixes zero.
The positive exponents in the $n$ and $m$ combinations are disjoint. An intersection would make $n/m$ a rational square; since $\gcd(n,m)=1$, both $n,m$ would be odd squares, impossible when $m-n=4$. It follows that each difference $f^{(r)}_0-\vartheta_{0,r}$ is constant. Such a finite linear combination of these theta functions cannot be a nonzero constant: its smallest positive exponent comes from the largest $d$ with a nonzero coefficient, and cannot cancel. Therefore
\begin{equation}\label{B:endpoints}
 f^{(r)}_0=\vartheta_{0,r},\qquad f^{(r)}_r=\vartheta_{r,r}.
\end{equation}
Linear independence and \eqref{B:weightedtranspose} show that both endpoint rows and both endpoint columns of $M_r$ are the corresponding rows and columns of the identity.

\subsection{Positivity and the character formulas}
The next positivity fact is stronger than positivity of a VOA character alone.
\begin{lemma}\label{B:Virpositive}
Let $E$ be an ordinary, lower-bounded $V$-module. The coefficient of $q^{h+n/4}$ in $\eta\ch E$ is a nonnegative integer if either $h<0$, or $h\not\equiv0,-n/4\pmod\Z$.
\end{lemma}
\begin{proof}
At this central charge, a rational Virasoro highest weight can have a reducible Verma module only at a positive Kac label with $r=s$, hence at
$h=n(s^2-1)/4\geq0$. This follows by separating the irrational coefficient of $t$ in the Kac formula, as in Section~\ref{A:app:rat}. All these weights are congruent to $0$ or $-n/4$ modulo integers. The character of an irreducible Verma module becomes the single monomial $q^{h+n/4}$ after multiplication by $\eta$. A non-Verma irreducible can therefore contribute a subtraction only in those two congruence classes and at nonnegative weights.

Virasoro composition multiplicities are nonnegative integers. This character assertion also applies to a module not finitely generated over Virasoro: at any fixed weight cutoff, finitely many homogeneous generators suffice to reproduce the truncated character, and their category-$\mathcal O$ modules have finite length at this generic central charge. Composition-factor additivity then gives the stated coefficient. Equivalently one can use the finite weight truncations of the block decomposition. Under either condition on $h$, no negative subtraction contributes at that weight.
\end{proof}

For $1\leq b<r$ set $h_{b,r}=b^2/(4r)-n/4$. For $r=n$ these numbers are all negative. For $r=m$ they are negative when $b\leq m-3$; the only others are
\[
 h_{m-2,m}=1/m,\qquad h_{m-1,m}=\tfrac12+\tfrac1{4m}.
\]
Both lie outside the two exceptional congruence classes in Lemma~\ref{B:Virpositive}. In a nonendpoint row $a$, the matrix coefficient $(M_r)_{ab}$ equals the coefficient at $q^{b^2/(4r)}$ in $\eta\ch A_a$ or $\eta\ch B_a$, because the theta supports are disjoint and the leading coefficient of $\vartheta_{b,r}$ is one. Consequently
\[
 (M_r)_{ab}\in\Z_{\geq0}\qquad(1\leq a,b<r).
\]
The zero row of the Fourier matrix in \eqref{B:WeilS} is
$(1,2,\ldots,2,1)/\sqrt{2r}$. Commutation of $M_r$ with that matrix, together with \eqref{B:endpoints}, therefore gives
\[
 \sum_{a=1}^{r-1}(M_r)_{ab}=1\qquad(1\leq b<r).
\]
The interior block has $r-1$ rows and columns. Each column thus has a single entry one. No row is zero, since each represents the character of a nonzero simple, and the endpoint columns have already vanished on these rows. Thus each row also has one entry one: $M_r$ is a permutation fixing both endpoints.

Relabel by these permutations. Commutation with $T$ preserves parity of the indices, because an odd square and an even square cannot be congruent modulo $4r$. They commute with the modular matrices, so the formulas already proved for $S,T$ are unchanged. They fix $V,X,J,Y$; Verlinde's formula then shows that they also preserve fusion by $J$. This proves the following full classification.

\Needspace{15\baselineskip}
\begin{theorem}\label{B:characterthm}
For every odd $n\geq3$, the simple $V$-modules have the following lowest weights and characters. Every lowest space is one-dimensional; the index ranges are $1\leq a<n$ and $1\leq b<m$.
\[
\begin{array}{lll}\toprule
\text{module}&\text{lowest weight}&\text{character}\\\midrule
V&0&(\Lambda_{n,n}-\Lambda_{m,m})/2\\
X&0&(\Lambda_{n,n}+\Lambda_{m,m})/2\\
J&3n/4&(\Lambda_{0,n}-\Lambda_{0,m})/2\\
Y&-n/4&(\Lambda_{0,n}+\Lambda_{0,m})/2\\
A_a&a^2/(4n)-n/4&\Lambda_{a,n}\\
B_b&b^2/(4m)-n/4&\Lambda_{b,m}\\\bottomrule
\end{array}
\]
Moreover $JA_a=A_{n-a}$ and $JB_b=B_{m-b}$. The odd indices in the two families are the $Z$ and $Q$ objects in the centralizer. 
\end{theorem}
The weight of $J$ follows from the first surviving term of
$(\vartheta_{0,n}-\vartheta_{0,m})/2=q^n+\cdots$. All other weights follow directly from their displayed theta leading terms. 

\subsection{A presentation of Zhu's algebra}
The character theorem gives an exact uniform presentation after a harmless change of generators. Define
\begin{equation}\label{B:ZhuP}
\begin{split}
 P_n(t)={}&t(t-3n/4)(t+n/4)\\
 &\cdot\prod_{a=1}^{n-1}\left(t-\frac{a^2}{4n}+\frac n4\right)
 \prod_{b=1}^{n+3}\left(t-\frac{b^2}{4(n+4)}+\frac n4\right).
\end{split}
\end{equation}
Its $2n+5$ roots are distinct. Within either family this is immediate; between the two families equality would make $n/m$ a rational square. None equals $0$ or the two exceptional endpoint weights. Since $A(V)$ is semisimple and its simples have the displayed $L_0$ eigenvalues,
$P_n([T])=0$, and this is the minimal polynomial of $[T]$.

The only repeated $L_0$ eigenvalue is zero, on $V,X$. The $W$ eigenvalues on these two tops differ: otherwise the two characters of the algebra generated by $[T],[W]$ would coincide. The vacuum eigenvalue is zero. Rescale $W$ to make its $X$ eigenvalue one, and put
\[
 e=\frac{P_n([T])/[T]}{P_n'(0)}[W].
\]
Here $P_n(t)/t$ means polynomial division before evaluation, not division by a possibly noninvertible element. Then $e$ is the primitive idempotent of $X$. It follows that
\begin{equation}\label{B:ZhuPresentation}
 A(V)\cong\C[t,e]/\big(P_n(t),\ te,\ e^2-e\big),
 \qquad t\longmapsto[T].
\end{equation}
Indeed these elements separate all $2n+6$ simple characters, so the map is surjective; its source has dimension $(2n+5)+1=2n+6$, proving injectivity. This presentation follows from the module classification; the rank argument used only the independent bound on $A(X)$.

\subsection{Proof of Theorem~\ref{A:cat:main}}\label{A:cat:proof-main}
\begin{proof}
Theorem~\ref{E:foundations} establishes the five foundational properties of $V_{n+2}$ for every odd $n\geq3$, including strong rationality, the lattice extension and its branching rules. We now collect the deductions that prove the four parts of the theorem.

\emph{Part~\textup{(i)}.} Theorem~\ref{B:twogen} bounds the relevant fusion multiplicities by two. Applied to the induction multiplicities in \eqref{B:XX}, this gives $b_i^2-1\leq2$, hence $b_i=1$. Corollary~\ref{B:rankthm} therefore proves that the simple objects already constructed exhaust the category and that the ranks are $2n+6$ and $n+3$. Section~\ref{B:sec:twogen} also proves that the lowest spaces are one-dimensional and that every simple is self-contragredient: the contragredient has the same top-space eigenvalues for the two generators $[T],[W]$ of Zhu's algebra. Section~\ref{sec:modular} determines the modular matrices \eqref{B:Sfull}, \eqref{B:Tfull} and \eqref{B:VOAmatrices}, and Section~\ref{B:sec:galois} proves the stated Galois comparison. Finally, Theorem~\ref{B:characterthm} identifies the simple modules by their explicit characters and lowest weights. These results prove every assertion in part~\textup{(i)}.

\emph{Part~\textup{(ii)}.} The lattice extension exists by Theorem~\ref{A:lat:main}. Section~\ref{sec:extension} proves uniqueness by determining its multiplication on $X\bo_V X$: both the vacuum channel and the $X$ channel are one-dimensional, and associativity fixes the ratio $a/b^2$ of their coefficients. Rescaling $X$ then identifies every simple extension multiplication with the lattice multiplication, while acting identically on $V$.

\emph{Part~\textup{(iii)}.} Lemma~\ref{A:cat:etale} proves that $A=\one\oplus X$ is connected \'etale, equips $\Ncat_n=(\D_n)_A$ with its spherical structure and proves the center factorization \eqref{A:cat:NGcenter}. Lemma~\ref{A:cat:ranklemma} and the lattice branching rules give exactly $n+1$ simple $A$-modules. The $n$ local modules form the pointed subcategory on $\Z/n$; the remaining simple is $R$. Section~\ref{A:cat:neargroup} computes $\FP(R)=t-1$, proves $gR=Rg=R$ and $R^*=R$, and uses Frobenius reciprocity and dimensions to obtain $R^2=\bigoplus_g g\oplus nR$. This proves part~\textup{(iii)}.

\emph{Part~\textup{(iv)}.} At $n=N^2$, Section~\ref{A:cat:condensation} constructs $\Fcat_N$ by condensing the regular algebra of $N\Z/N^2$ in the central pointed factor. This subgroup is Lagrangian, so its condensation removes that factor and gives the ribbon equivalence \eqref{A:cat:centerF}. The inherited spherical structure is identified there using the equivalent description $\Fcat_N\simeq(\D_{N^2})_{B_N}$. Section~\ref{A:cat:HIfusion} determines all $2N$ simple objects, their left and right actions, self-duality and the multiplicity-one decomposition of $\rho^2$. Together with the spherical dimension calculation in Section~\ref{A:cat:spherical}, it gives $d(\rho)=(N-\sqrt{N^2+4})/2$. Section~\ref{A:cat:associator} computes the associator from balanced multiplication maps, obtaining \eqref{A:cat:omega}; its cyclic invariant is $e^{2\pi i/N}$, so its class has exact order $N$. The semion factorization in Theorem~\ref{E:foundations} now gives both equivalences in \eqref{A:cat:realization}. Section~\ref{B:sec:HIcomparison} proves the Galois comparison with $D^1\mathrm{Hg}_N$. Finally, Section~\ref{A:cat:centercriterion} proves that $\D_n$ is a center exactly when $n$ is a square, by the Lagrangian-subgroup criterion for its pointed Witt class. This completes part~\textup{(iv)} and the theorem.
\end{proof}

\section{The inverse ribbon semion}\label{E:auxiliary}
From this point onward $n=N^2$ with $N\geq3$ odd. In particular $n\equiv1\pmod8$. The results for the diagonal family above give
\begin{equation}\label{E:Jdata}
J^2\cong V,\qquad h_J=\frac{3n}{4},\qquad
c_{J,J}=i,\qquad\theta_J=-i,\qquad d(J)=-1.
\end{equation}
Write $\Scat^-:=\langle J\rangle$; the superscript records the negative dimension, not a reversal of braiding. Thus
\[
\Rep(V)\simeq\D_n\bo\Scat^-.
\]

\subsection{A conformal shift of the \texorpdfstring{$E_7$}{E7} lattice}
Let $K$ be the positive-definite $E_7$ root lattice. Its dual quotient has order two. A minuscule weight $u$ representing the nonzero class satisfies
\begin{equation}\label{E:u}
u\in K^*,\quad u\notin K,\quad 2u\in K,\quad u^2=\frac32.
\end{equation}
Define the shifted lattice VOA $\mathcal W=V_{K,u}$ by keeping its vacuum, state-field correspondence and translation, and replacing its conformal vector by
\begin{equation}\label{E:shift}
\omega_u=\omega_K+u(-2)\one.
\end{equation}
Its Virasoro operators, central charge and momentum weights are
\begin{align}
L_u(r)&=L_K(r)-(r+1)u(r),& c(\mathcal W)&=7-12u^2=-11,\label{E:shiftmodes}\\
\wt_u(e^p)&=\frac{p^2}{2}-(u,p)
=\frac{(p-u)^2}{2}-\frac34.\label{E:shiftweight}
\end{align}
Oscillators add their usual nonnegative integer degrees. Since $u\in K^*$, $\mathcal W$ is integer graded. Completing the square proves boundedness below and finite-dimensionality of every weight space.

There are exactly two simple weak modules, with underlying lattice modules $V_K$ and $V_{K+u}$. Both remain ordinary for $L_u(0)$. The fields and weak-module identities are unchanged by the conformal shift, so the regularity of the positive-definite lattice vertex algebra gives regularity here. The shifted-lattice theorems of Dong--Mason \cite[Theorems~3.1, 4.1 and 4.4]{E:DM} give these facts and self-contragredience, the latter because $2u\in K$.

Denote the second shifted module by $M=V_{K+u,u}$. At a momentum $p=u+a$, $a\in K$, formula \eqref{E:shiftweight} becomes
\begin{equation}\label{E:Mweight}
\wt_u(e^{u+a})=\frac{a^2}{2}-\frac34.
\end{equation}
It has a one-dimensional lowest space at weight $-3/4$, and hence $\theta_M=i$. Lattice fusion gives $M\bo_{\mathcal W} M\cong \mathcal W$.

\begin{lemma}\label{E:inverse}
The category of ordinary $\mathcal W$-modules, with its shifted ribbon structure, is ribbon equivalent to $(\Scat^-)^{\mathrm{rev}}$. In particular
\[
F(M,M,M)=-1,\qquad c_{M,M}=-i,\qquad
\theta_M=i,\qquad d(M)=-1.
\]
\end{lemma}
\begin{proof}
The shift leaves the lattice vertex operators and translation operator unchanged. To identify the tensor structures, apply \cite[Proposition~2.17]{A:McR} to $v=u(-1)\one$: here $v_0\omega_K=0$, and the original and shifted vacuum gradings are nonnegative, the latter by Lemma~\ref{E:Wzero}. The proposition identifies the module categories with the same braided tensor structure. The twist, however, must be calculated with the shifted $L(0)$. The nontrivial lattice class has self-braiding
\[
c_{M,M}=e^{\pi i u^2}=-i
\]
and the order-two semion associator. The twist is changed by the new $L(0)$; \eqref{E:Mweight} gives $\theta_M=i$. The rank-two spin-statistics relation gives $d(M)=\theta_M/c_{M,M}=-1$. These data are the inverse braiding and inverse twist of \eqref{E:Jdata}, with the same dimension. They determine the claimed ribbon equivalence.
\end{proof}

This shift addresses two distinct extension conditions. The following table records the nontrivial auxiliary object $s$ and the product current $J\otimes s$:
\begin{center}
\begin{tabular}{lcccc}
\toprule
Auxiliary lattice module&$c_{s,s}$&$\theta_s$&$c_{Js,Js}$&$\theta_{Js}$\\
\midrule
Unshifted $A_1$&$i$&$i$&$-1$&$1$\\
Unshifted $E_7$&$-i$&$-i$&$1$&$-1$\\
Shifted $E_7$, namely $M$&$-i$&$i$&$1$&$1$\\
\bottomrule
\end{tabular}
\end{center}
The last row supplies both braided commutativity and integral conformal grading for an ordinary VOA extension.

\subsection{Nonnegative grading and the zero-weight space}
\begin{lemma}\label{E:Wzero}
The shifted VOA $\mathcal W$ has no negative weights and $\dim \mathcal W_0=56$.
\end{lemma}
\begin{proof}
For a nonzero $p\in K$, put $p^2=2r$, where $r$ is a positive integer. Since $(u,p)$ is integral and $(u,p)\leq\sqrt{3r}$, it is at most $r$ for $r=1,2,3$. For $r\geq4$ one has $\sqrt{3r}<r$. Thus \eqref{E:shiftweight} is nonnegative. No oscillator contributes at weight zero, so $\dim \mathcal W_0$ is the number of momenta satisfying $p^2/2=(u,p)$.

For an explicit count, realize
\[
K=E_8\cap(e_1+e_2)^\perp,\qquad
u=(1/2,-1/2,1,0,0,0,0,0)
\]
in the usual orthonormal coordinates. The $E_8$ roots are the vectors $\pm e_i\pm e_j$ and the half-integral vectors $(\pm1/2)^8$ with an even number of minus signs. Among the roots perpendicular to $e_1+e_2$, exactly $27$ have $(u,a)=1$: eleven are integral and sixteen are half-integral. Every zero-weight momentum has $r\leq3$. For $r=1$ these are the $27$ roots just counted; the involution $p\mapsto2u-p$ pairs them with all $r=2$ solutions. At $r=3$, equality in Cauchy--Schwarz forces $p=2u$. Together with $p=0$, the count is
\[
1+27+27+1=56.
\]
Equivalently, these are the $56$ shortest vectors in the coset $K+u$, translated by $u$.
\end{proof}

The invariant bilinear form on $\mathcal W$ is nondegenerate even though its vacuum is isotropic. In the shifted contragredient pairing, the vacuum pairs with $e^{2u}$. Positivity of a Hermitian form, and the normalization $(\one,\one)=1$, are not hypotheses of the extension result used below.

\section{The extension and ribbon cancellation}\label{E:extension}
We first construct the VOA and establish the finiteness and duality properties needed for its tensor category. We then identify that entire category by local induction, keeping track of the tensor product, braiding and twist.

Put
\[
B=V\otimes \mathcal W,\qquad Q=J\otimes M.
\]
This use of $B$ is local to this section; it is distinct from the central algebra used to define $\Fcat_N$. The product category is
\begin{equation}\label{E:basecat}
\Rep(B)\simeq\D_n\bo\Scat^-\bo(\Scat^-)^{\mathrm{rev}}.
\end{equation}

\subsection{Existence, simplicity and regularity}
Equations \eqref{E:Jdata} and \eqref{E:Mweight} imply
\begin{equation}\label{E:Qdata}
Q^2\cong B,\quad c_{Q,Q}=1,\quad\theta_Q=1,\quad d(Q)=1,
\quad h_Q=\frac{3(n-1)}4.
\end{equation}
For odd $N\geq3$, this last number is a positive multiple of six.

\begin{proposition}\label{E:existence}
The direct sum $B\oplus Q$ admits a simple, regular, $C_2$-cofinite VOA structure extending $B$, unique up to an isomorphism fixing $B$. It is nonnegatively integer graded and has central charge $-6n-10$.
\end{proposition}
\begin{proof}
\emph{The base VOA.} Theorem~\ref{E:foundations} makes $V$ strongly rational. Thus the tensor product $V\otimes V_K$ is regular and $C_2$-cofinite. Its weak-module identities depend only on the fields, which the conformal shift leaves unchanged. Every weak $B$-module therefore decomposes into simple tensor-product modules. These remain ordinary: their $V$-weights are bounded below, and their lattice weights are a positive-definite quadratic function of the momentum, as in \eqref{E:shiftweight}, plus nonnegative oscillator degrees. Each total weight space is finite-dimensional. Hence $B$ is regular, and the unchanged $C_2$ space proves $C_2$-cofiniteness. The product of the invariant forms on $V$ and $\mathcal W$ is nondegenerate, so $B$ is self-contragredient as well as simple. Its vacuum grading is nonnegative by Lemma~\ref{E:Wzero}. Finally, the conformal-shift comparison used in Lemma~\ref{E:inverse}, together with the tensor-product construction, gives the braided factorization \eqref{E:basecat}.

\emph{The extension multiplication.} There is one nonzero intertwining-operator channel from $Q\times Q$ to $B$, since $Q\bo_B Q\cong B$. The two semion associators can be written in normalized gauges as $F(J,J,J)=F(M,M,M)=-1$. Their product on $Q$ is therefore
\begin{equation}\label{E:Qcoherence}
F(Q,Q,Q)=1,\qquad c_{Q,Q}=i(-i)=1.
\end{equation}
The associators involving the unit are normalized to one. Thus the order-two extension has trivial associativity and commutativity obstructions. Also $h_Q=3(n-1)/4\in6\Z_{>0}$, so its fields have the integral grading required for a VOA.

These checks make $Q$ an even simple-current unit. Carnahan's extension theorem \cite[Theorem~3.2.12]{E:Carnahan} now applies: its base is a simple $C_2$-cofinite M\"obius vertex algebra with finite weight spaces and a nondegenerate invariant form, and is rational. It gives a VOA extension on $B\oplus Q$ that is regular and $C_2$-cofinite. In particular, the theorem allows the present $B_0\neq\C\one$.

More explicitly, choose the nonzero intertwining operator $\mathcal I$ of type $\binom{B}{Q\ Q}$ supplied by this construction. The four components of the state-field map are
\begin{equation}\label{E:extensionfields}
\begin{aligned}
Y(b,z)b'&=Y_B(b,z)b',& Y(b,z)q&=Y_Q(b,z)q,\\
Y(q,z)b&=e^{zL(-1)}Y_Q(b,-z)q,& Y(q,z)q'&=\mathcal I(q,z)q',
\end{aligned}
\end{equation}
for $b,b'\in B$ and $q,q'\in Q$. The vacuum and conformal vector lie in $B$. Equation~\eqref{E:Qcoherence} is precisely the compatibility of the last component with associativity and skew-symmetry; the extension theorem supplies the Jacobi identity for these fields.

\emph{Simplicity and uniqueness.} As a $B$-module, $B\oplus Q$ is a multiplicity-free sum of two inequivalent simples. Any VOA ideal is a $B$-submodule. If it contains $B$, it contains the vacuum and hence the whole extension. If it contains $Q$, the nonzero $Q\times Q\to B$ product forces it to contain $B$. Thus the extension is simple.

For any other simple extension on the same $B$-module, the first three components of \eqref{E:extensionfields} are fixed by the $B$-action and skew-symmetry. The last is $a\mathcal I$ for a scalar $a$, and $a\neq0$ since otherwise $Q$ would be an ideal. Choose $s\in\C^\times$ with $s^2=a$. The map $\id_B\oplus s\id_Q$ is an isomorphism from that extension to the one with last component $\mathcal I$: on $Q\times Q$ the required equality is exactly $a=s^2$. This proves uniqueness with $B$ fixed.

\emph{Grading and central charge.} The $B$ summand has nonnegative integral weights, while $Q$ starts at the positive integer $h_Q\geq6$ and has integral weight increments. The conformal vector is that of $B$, giving
\[
c(\mathscr{V}_N)=c(V)+c(\mathcal W)=(1-6n)-11=-6n-10.
\]
\end{proof}

\begin{proposition}\label{E:selfdual}
The extension $\mathscr{V}_N$ is self-contragredient.
\end{proposition}
\begin{proof}
Since $h_Q\geq6$, its degree-zero and degree-one spaces coincide with those of $B$, with the same $L(1)$ action. Consequently
\[
(\mathscr{V}_N)_0/L(1)(\mathscr{V}_N)_1=B_0/L(1)B_1.
\]
The space of invariant bilinear forms on $B$ is $\Hom_B(B,B')$, where $B'$ is its restricted contragredient. Since $B$ is simple and self-contragredient, this space has dimension one. Li's invariant-form theorem \cite{A:Li} identifies it with the dual of $B_0/L(1)B_1$. Applying the same theorem to $\mathscr{V}_N$ therefore gives a one-dimensional space of invariant forms on the extension. Choose a nonzero form in this space. Its radical is an ideal; simplicity from Proposition~\ref{E:existence} forces that radical to be zero. Invariance makes distinct $L(0)$-eigenspaces orthogonal, and each eigenspace is finite-dimensional. The nondegenerate form thus gives an isomorphism $\mathscr{V}_N\cong\mathscr{V}_N'$ of modules. No normalization of the vacuum's self-pairing is needed.
\end{proof}

\subsection{Modularity without CFT type}\label{E:nonCFTsection}
\begin{proposition}\label{E:nonCFTmodular}
The category of grading-restricted $\mathscr{V}_N$-modules is a semisimple modular
tensor category, despite $\dim (\mathscr{V}_N)_0=56$.
\end{proposition}
\begin{proof}
Proposition~\ref{E:existence} gives a simple, regular, $C_2$-cofinite,
nonnegatively $\Z$-graded VOA $\mathscr{V}_N$, and
Proposition~\ref{E:selfdual} proves it is self-contragredient. In particular,
regularity implies rationality and hence semisimplicity of Zhu's algebra
$A(\mathscr{V}_N)$. McRae's Main Theorem~3 \cite{A:McR} applies to precisely an
$\mathbb N$-graded, simple, self-contragredient, $C_2$-cofinite VOA with semisimple
Zhu algebra. It therefore proves that the full category of
grading-restricted generalized $\mathscr{V}_N$-modules is a semisimple modular
tensor category. No hypothesis that $(\mathscr{V}_N)_0=\C\one$ is used. This argument
establishes modularity independently of the subsequent identification of
the category with $\D_n$; in particular, that identification is not used
to deduce the rigidity required in McRae's modularity results.
\end{proof}

\subsection{The diagonal algebra and its local modules}\label{E:localmodules}
Under \eqref{E:basecat}, the extension algebra is
\begin{equation}\label{E:diagonal}
\mathscr A=\one\oplus Q
=\one_{\D_n}\bo\bigl((\one\bo\one)\oplus(J\bo M)\bigr).
\end{equation}
The factor in parentheses is the diagonal algebra in $\Scat^-\bo(\Scat^-)^{\mathrm{rev}}$. Equations~\eqref{E:Qdata}--\eqref{E:Qcoherence} identify the symmetric subcategory generated by $Q$ with $\Rep(\Z/2)$. The multiplication on $\one\oplus Q$ is its regular commutative algebra, so it is separable over $\C$. It is connected because
\[
\Hom_{\Rep(B)}(B,\mathscr A)=\Hom_B(B,B)\cong\C;
\]
connectedness here concerns the multiplicity of the tensor unit, not the dimension of the weight-zero space. Moreover $\theta_{\mathscr A}=\id$ and $d(\mathscr A)=2$. Thus $\mathscr A$ is a connected \'etale algebra with trivial twist, and its module category is semisimple. We next calculate its local modules and their ribbon structure explicitly.

\begin{proposition}\label{E:category}
There is a ribbon equivalence
\[
\Rep(\mathscr{V}_N)\simeq\D_n\simeq Z(\Fcat_N).
\]
For each simple $T\in\D_n$, the corresponding irreducible ordinary $\mathscr{V}_N$-module is
\begin{equation}\label{E:ET}
\mathsf M_T=(T\otimes \mathcal W)\oplus((J\bo_V T)\otimes M).
\end{equation}
These modules are pairwise nonisomorphic and exhaust the irreducible ordinary modules.
\end{proposition}
\begin{proof}
\emph{The extension correspondence.} Put $\mathcal C=\Rep(B)$. Since $B$ is $\mathbb N$-graded and $C_2$-cofinite, its full category of grading-restricted generalized modules has the Huang--Lepowsky--Zhang vertex tensor structure \cite[Section~2.3]{A:McR}. Regularity makes this category semisimple. Proposition~2.17 of the same reference identifies the braided structure before and after the lattice conformal shift, as checked in Lemma~\ref{E:inverse}; thus it is the product structure in \eqref{E:basecat}. The extension $\mathscr{V}_N=B\oplus Q$ is an object of this category with the same conformal vector as $B$. The extension theorem of Creutzig--Kanade--McRae \cite[Theorem~3.65]{E:CKMext} consequently gives a braided equivalence
\begin{equation}\label{E:localcorrespondence}
\mathcal C_{\mathscr A}^{\mathrm{loc}}\simeq\Rep(\mathscr{V}_N),
\end{equation}
initially for those modules whose restrictions lie in $\mathcal C$.

This accounts for the full grading-restricted category. Indeed, an ordinary $\mathscr{V}_N$-module restricts, by regularity of $B$, to a direct sum of simple ordinary $B$-modules. There are finitely many such types. Infinite multiplicity of any one type would give an infinite-dimensional $L(0)$-eigenspace, because the conformal vector has not changed. Its restriction is therefore a finite direct sum and belongs to $\mathcal C$. Regularity of $\mathscr{V}_N$ also makes every grading-restricted generalized module a finite direct sum of ordinary simples. Conversely a local $\mathscr A$-module in $\mathcal C$ acquires the extension action with the same grading, so is grading-restricted. This proves \eqref{E:localcorrespondence} for precisely the category in the proposition.

\emph{Locality and the simple objects.} Write the simple objects of $\mathcal C$ as
\[
S_{T,a,b}=(J^{\bo_V a}\bo_V T)\otimes M^{\bo_{\mathcal W}b},
\qquad T\in\Irr(\D_n),\quad a,b\in\{0,1\},
\]
where the zeroth powers mean $V$ and $\mathcal W$. In either semion factor the monodromy of the nontrivial object with itself is $-1$. Since $T$ centralizes $J$, the monodromy with $Q$ is
\begin{equation}\label{E:localitytest}
c_{S_{T,a,b},Q}\,c_{Q,S_{T,a,b}}
=(-1)^{a+b}\id_{Q\bo_B S_{T,a,b}}.
\end{equation}
Let $\operatorname{Ind}(S)=\mathscr A\bo_B S$ denote free induction. By \cite[Proposition~2.65]{E:CKMext}, it is local exactly when the monodromy of $S$ with $\mathscr A$ is trivial. Formula~\eqref{E:localitytest} says that this is exactly $a=b$. Tensoring with $Q$ sends $(T,a,b)$ to $(T,a+1,b+1)$, with the last two labels read modulo two. Hence the local labels occur in the free pairs
\[
\{S_{T,0,0},S_{T,1,1}\},\qquad T\in\Irr(\D_n).
\]
There are no fixed labels, since the second lattice factor changes from $\mathcal W$ to $M$ or conversely.

Induction--restriction adjunction gives, for every simple $S$,
\begin{equation}\label{E:inductionend}
\End_{\mathscr A}(\operatorname{Ind}(S))
\cong\Hom_{\mathcal C}(S,S\oplus(Q\bo_B S))\cong\C.
\end{equation}
The last equality uses the absence of fixed labels. Since the $\mathscr A$-module category is semisimple, this proves that each induced module is simple. Induction identifies the two labels in each pair: the isomorphism $\mathscr A\bo_B Q\cong\mathscr A$ follows from its regular action. For the representative $S_{T,0,0}=T\otimes\mathcal W$, the underlying $B$-module is exactly
\[
\operatorname{Ind}(T\otimes\mathcal W)
=(T\otimes\mathcal W)\oplus((J\bo_V T)\otimes M),
\]
as in \eqref{E:ET}. The same adjunction gives
\begin{equation}\label{E:inductionhom}
\Hom_{\mathscr A}(\mathsf M_T,\mathsf M_{T'})
\cong\Hom_{\D_n}(T,T'),
\end{equation}
because the cross term has lattice factors $\mathcal W$ and $M$, which are nonisomorphic. This proves pairwise nonisomorphism.

To prove exhaustion, let $L$ be a simple local $\mathscr A$-module. Its action $Q\bo_B L\to L$ is invertible: acting twice is the nonzero multiplication $Q\bo_B Q\cong B$ followed by the unit action. Locality then forces the monodromy with $Q$ on its underlying object to be the identity. Any simple $B$-constituent $S$ of $L$ therefore has $a=b$. Adjunction gives a nonzero $\mathscr A$-module map $\operatorname{Ind}(S)\to L$. Both modules are simple by \eqref{E:inductionend}, so this map is an isomorphism. Thus every local simple is one of the displayed $\mathsf M_T$.

\emph{The tensor product and braiding.} Define
\begin{equation}\label{E:ribbonfunctor}
\Phi:\D_n\longrightarrow\Rep(\mathscr{V}_N),\qquad
\Phi(T)=\operatorname{Ind}(T\otimes\mathcal W).
\end{equation}
Every object $T\otimes\mathcal W$ centralizes $\mathscr A$. Induction on this centralizer is a braided tensor functor \cite[Theorems~2.67 and 3.68]{E:CKMext}. In particular its tensor and unit isomorphisms are
\begin{equation}\label{E:inductiontensor}
\Phi(T)\bo_{\mathscr{V}_N}\Phi(T')\cong\Phi(T\bo_V T'),
\qquad \Phi(V)\cong\mathscr{V}_N.
\end{equation}
These isomorphisms arise from multiplication of the two $\mathscr A$ factors in the relative tensor product. The centralizer condition makes them compatible with the braiding, and associativity of $\mathscr A$ gives their tensor coherence. Equation~\eqref{E:inductionhom}, extended over direct sums, proves full faithfulness; the preceding exhaustion argument proves essential surjectivity. Thus $\Phi$ is a braided tensor equivalence, rather than merely a matching of fusion rules or modular matrices.

\emph{The ribbon structure and the center.} The twist on a VOA module is $e^{2\pi i L(0)}$. On the first summand of $\mathsf M_T$ it is $\theta_T$, since $\theta_{\mathcal W}=1$. On the second, centralization of $J$ by $T$ gives
\begin{equation}\label{E:twistpreserved}
\theta_{J\bo_V T}\theta_M
=\theta_J\theta_T\theta_M=(-i)i\theta_T=\theta_T.
\end{equation}
Thus $\Phi$ preserves the twist on each simple and hence on every object. The two categories are rigid, by the results for $V$ and Proposition~\ref{E:nonCFTmodular}, so this proves that $\Phi$ is a ribbon equivalence.

Finally Section~\ref{A:cat:condensation} constructs the spherical category $\Fcat_N$ and proves the ribbon equivalence $\D_n\simeq Z(\Fcat_N)$ in \eqref{A:cat:centerF}. Composing it with $\Phi^{-1}$ proves the asserted realization of the center, with its inherited spherical structure and associators.
\end{proof}

Together with Proposition~\ref{E:nonCFTmodular}, this proves the existence and categorical parts of Theorem~\ref{E:main} without applying a CFT-type modularity theorem to $\mathscr{V}_N$. The fusion coefficients of the $\mathsf M_T$ are exactly those of $\D_n$, and the $S$ matrix is the categorical matrix of Section~\ref{sec:modular}, with its stated normalization. The categorical twists should be distinguished from the trace $T$ matrix, which also includes the factor $e^{-2\pi i c(\mathscr{V}_N)/24}$.

\section{A rank-eight lattice parent}\label{E:lattice}
This section uses the concrete lattice embedding and the half-lattice decomposition of Section~\ref{A:app:lat}. It gives actual vertex operators for \eqref{E:def}, in addition to the abstract simple-current construction.

\subsection{The lattice and its discriminant form}
Choose a root $\alpha\in E_8$, with $\alpha^2=2$, and identify $K=E_8\cap\alpha^\perp$ with $E_7$. Choose the sign of $u$ so that
\[
E_8=(\Z\alpha\oplus K)+\Z(\alpha/2+u).
\]
Set $\beta=N\alpha$ and
\begin{equation}\label{E:LN}
L_N=(\Z\beta\oplus K)+\Z(\beta/2+u),\qquad h_N=\beta/2+u.
\end{equation}
The new coset has order two over $\Z\beta\oplus K$, and
\[
h_N^2=\frac{N^2+3}{2}\in2\Z.
\]
Thus $L_N$ is even and positive definite of rank eight. Its determinant is $(2N^2\cdot2)/2^2=N^2$.

\begin{lemma}\label{E:discriminant}
There are equalities and an isometry of finite quadratic modules
\begin{align}
L_N&=\{x\in E_8:(x,\alpha)\equiv0\pmod N\},\label{E:kernel}\\
(L_N^*/L_N,q)&\cong(\Z/N^2,\ x\longmapsto x^2/N^2\pmod\Z).\label{E:discriminantform}
\end{align}
In particular $[E_8:L_N]=N$.
\end{lemma}
\begin{proof}
Every generator in \eqref{E:LN} has pairing with $\alpha$ divisible by $N$, and $h_N\in E_8$ because $N$ is odd. The homomorphism $E_8\to\Z/N$ given by $x\mapsto(x,\alpha)$ is surjective: roots of $E_8$ have pairing one with $\alpha$. Its kernel has index $N$, which agrees with the determinant of $L_N$, proving \eqref{E:kernel}.

The vector $\alpha/N$ lies in $L_N^*$ by \eqref{E:kernel}. Moreover $L_N\cap\R\alpha=N\Z\alpha$: a vector in the odd glue coset has a nontrivial $K^*/K$ component, and cannot lie on the $\alpha$ line. Therefore $\alpha/N$ has exact order $N^2$ modulo $L_N$. Since the discriminant group has order $N^2$, it is generated by this element. Its quadratic value is $(\alpha/N)^2/2=1/N^2$, proving the formula.
\end{proof}

With the standard conformal vector, this is the concrete inner cyclic orbifold
\begin{equation}\label{E:orbifold}
V_{L_N}=V_{E_8}^{\langle\exp(2\pi i\alpha(0)/N)\rangle}.
\end{equation}
Indeed the automorphism fixes oscillators and acts on $e^x$ by $\exp(2\pi i(x,\alpha)/N)$. Its fixed momenta are exactly \eqref{E:kernel}. The subsequent conformal deformation is essential for the embedding of the diagonal family.

\subsection{The conformal vector and the subalgebra}
Fix $\lambda\ne0$ and put
\begin{equation}\label{E:ambientomega}
\omega_N^{\mathrm{amb}}
=\omega_{L_N}+h_N(-2)\one+\lambda D^2e^\beta.
\end{equation}
For the first two terms, $e^\beta$ is primary of weight zero, because $\beta^2/2=(h_N,\beta)=n$. Its self-OPE is regular. If $L^{(h_N)}_r$ denotes the shifted lattice Virasoro operator, the additional term gives
\[
L^{\mathrm{amb}}_r=L^{(h_N)}_r+\lambda r(r+1)(e^\beta)_{r-1}.
\]
For a weight-zero primary, $[L_r,(e^\beta)_s]=(-r-s-1)(e^\beta)_{r+s}$; the modes of $e^\beta$ commute with one another. Substitution proves the Virasoro relations, since the cross-term coefficient is
\[
\lambda(r+s)\bigl(r(r+1)-s(s+1)\bigr)
=\lambda(r-s)(r+s)(r+s+1).
\]
There is no added central term. Both $L(0)$ and $L(-1)$ are unchanged by $\lambda D^2e^\beta$. The resulting regular VOA $\mathcal P_N$ has
\[
c(\mathcal P_N)=8-12h_N^2=-6N^2-10.
\]
Regularity follows again from the unchanged lattice fields and the ordinary shifted grading of every lattice module.

The two lattice cosets in \eqref{E:LN} give
\begin{equation}\label{E:Pbranch}
\mathcal P_N=(\mathcal U_n\otimes \mathcal W)
\oplus(\mathcal M_n\otimes M),
\end{equation}
where $\mathcal U_n=V\oplus X$ and $\mathcal M_n=J\oplus Y$ are the deformed rank-one lattice VOA and its half-lattice module. The $V\otimes\mathcal W$ action preserves $J\otimes M$. For the remaining product, $J\bo_V J=V$ and $M\bo_{\mathcal W}M=\mathcal W$ force the image of $(J\otimes M)\times(J\otimes M)$ into $V\otimes\mathcal W$. This product is nonzero: its first factor is the nonzero lattice intertwiner on $J\times J$ proved in Section~\ref{A:app:lat}, and its second factor is the lattice intertwiner $M\times M\to\mathcal W$. Thus all products in
\[
(V\otimes \mathcal W)\oplus(J\otimes M)
\]
are closed under the actual lattice fields, and the subspace contains the vacuum and conformal vector \eqref{E:ambientomega}. It is therefore a conformal subalgebra of $\mathcal P_N$. The two-summand ideal argument in Proposition~\ref{E:existence} proves that it is simple, and the uniqueness assertion of that proposition identifies its vertex operators with those of $\mathscr{V}_N$, fixing $V\otimes\mathcal W$. This supplies a concrete lattice realization of the extension fields \eqref{E:extensionfields}. The complementary terms of \eqref{E:Pbranch} give the precise branching
\begin{equation}\label{E:parentbranch}
\mathcal P_N\downarrow \mathscr{V}_N=\mathscr{V}_N\oplus\mathsf M_X,
\qquad
\mathsf M_X=(X\otimes \mathcal W)\oplus(Y\otimes M).
\end{equation}

The inclusion $L_N\subset E_8$ gives a further conformal extension with the same vector \eqref{E:ambientomega}. Its underlying vertex algebra is $V_{E_8}$; with this conformal vector it is regular and has only one irreducible ordinary module. This use of holomorphicity permits negative conformal weights. Both ambient lattice algebras have minimum weight
\[
-\frac{h_N^2}{2}=-\frac{N^2+3}{4},
\]
attained at $e^{h_N}$. In contrast $\mathscr{V}_N$ itself is nonnegatively graded. The discriminant group in \eqref{E:discriminantform} is cyclic of order $N^2$; it is different from the split group $(\Z/N)^2$ occurring in an untwisted pointed parent.

\section{Characters and examples}\label{E:characters}
For a lattice coset write
\[
\Theta_{K+a}(\tau)=\sum_{p\in K+a}q^{p^2/2},\qquad
\theta_{r,s}(\tau)=\sum_{\ell\in\Z}q^{(2s\ell+r)^2/(4s)},
\qquad q=e^{2\pi i\tau}.
\]
The character convention is $\ch F=\operatorname{tr}_F q^{L(0)-c/24}$. Completing the square in \eqref{E:shiftweight} gives
\begin{equation}\label{E:auxchars}
\ch \mathcal W=\frac{\Theta_{K+u}}{\eta^7},\qquad
\ch M=\frac{\Theta_K}{\eta^7}.
\end{equation}
Thus the conformal shift exchanges the two standard $E_7$ characters.

\subsection{All irreducible characters}\label{E:allcharacters}
By \eqref{E:ET}, the character attached to $T\in\Irr(\D_n)$ is
\begin{equation}\label{E:allchars}
\ch\mathsf M_T
=\frac{\chi_T\Theta_{K+u}+\chi_{JT}\Theta_K}{\eta^7}.
\end{equation}
These are traces on constructed modules. Theorem~\ref{B:characterthm} makes \eqref{E:allchars} completely explicit. With $m=n+4$, the vacuum and the distinguished non-vacuum character are
\begin{align}
\ch \mathscr{V}_N&=
\frac{(\theta_{n,n}-\theta_{m,m})\Theta_{K+u}
+(\theta_{0,n}-\theta_{0,m})\Theta_K}{2\eta^8},\label{E:vacchar}\\
\ch\mathsf M_X&=
\frac{(\theta_{n,n}+\theta_{m,m})\Theta_{K+u}
+(\theta_{0,n}+\theta_{0,m})\Theta_K}{2\eta^8}.\label{E:Xchar}
\end{align}
For odd $a$ with $1\leq a<n$ and odd $b$ with $1\leq b<m$, the remaining characters are
\begin{align}
\ch\mathsf M_{A_a}&=
\frac{\theta_{a,n}\Theta_{K+u}+\theta_{n-a,n}\Theta_K}{\eta^8},\label{E:Achar}\\
\ch\mathsf M_{B_b}&=
\frac{\theta_{b,m}\Theta_{K+u}+\theta_{m-b,m}\Theta_K}{\eta^8}.\label{E:Bchar}
\end{align}
The odd labels are exactly the centralizer labels. The number of modules in this list is
\[
2+\frac{n-1}{2}+\frac{m-1}{2}=n+3=N^2+3.
\]
The lowest weight of \eqref{E:ET} is
\begin{equation}\label{E:lowest}
h(\mathsf M_T)=\min\{h_T,\ h_{JT}-3/4\}.
\end{equation}
If the two terms coincide, their bottom-space dimensions add. Their characters have nonnegative coefficients since both terms in \eqref{E:ET} are actual module characters.

\subsection{Rank-eight theta differences}
For a positive integer $r\equiv1\pmod4$, let $\beta_r^2=2r$, perpendicular to $K$, and define
\[
K(r)=(\Z\beta_r\oplus K)+\Z(\beta_r/2+u).
\]
This is an even positive-definite lattice of rank eight and determinant $r$. Its theta series is
\begin{equation}\label{E:thetaKr}
\Theta_{K(r)}=\theta_{0,r}\Theta_K+\theta_{r,r}\Theta_{K+u}.
\end{equation}
Since $n=N^2\equiv1\pmod8$, both $n$ and $n+4$ are allowed, and $K(n)=L_N$. The two distinguished characters simplify to
\begin{equation}\label{E:thetadiff}
\ch \mathscr{V}_N=\frac{\Theta_{K(n)}-\Theta_{K(n+4)}}{2\eta^8},
\qquad
\ch\mathsf M_X=\frac{\Theta_{K(n)}+\Theta_{K(n+4)}}{2\eta^8}.
\end{equation}
Their sum is $\Theta_{L_N}/\eta^8$, as required by \eqref{E:parentbranch}. In the ambient character, the shift $h_N\in L_N$ merely reindexes the lattice theta sum; the $D^2e^\beta$ term does not alter the grading. This explains the sum-and-difference identities at the level of the lattice extension.

\subsection{Numerical invariants}\label{E:invariants}
Since the second extension summand begins at $3(n-1)/4\geq6$,
\[
\dim(\mathscr{V}_N)_0=\dim V_0\dim \mathcal W_0=56.
\]
All simple $V$-modules have lowest weight at least $-n/4$, and all simple $\mathcal W$-modules have lowest weight at least $-3/4$. Formula \eqref{E:ET} therefore gives a lower bound $-(n+3)/4$ for every $\mathscr{V}_N$-module. It is attained in the $Y\otimes M$ summand of $\mathsf M_X$. Consequently
\begin{equation}\label{E:ceff}
h_{\min}(\mathscr{V}_N)=-\frac{n+3}{4},\qquad
c_{\mathrm{eff}}(\mathscr{V}_N)=(-6n-10)+6(n+3)=8.
\end{equation}
\begin{center}
\begin{tabular}{rrrrr}
\toprule
$N$&$c(\mathscr{V}_N)$&$h_Q$&$|\Irr(\mathscr{V}_N)|$&$h_{\min}(\mathscr{V}_N)$\\
\midrule
3&$-64$&6&12&$-3$\\
5&$-160$&18&28&$-7$\\
7&$-304$&36&52&$-13$\\
9&$-496$&60&84&$-21$\\
11&$-736$&90&124&$-31$\\
\bottomrule
\end{tabular}
\end{center}
Every entry follows from the uniform formulas; each row has $\dim(\mathscr{V}_N)_0=56$ and effective central charge eight. The row $N=9$ also illustrates that the construction includes composite odd parameters: the pointed class has order nine.

\subsection{Proof of Theorem~\ref{E:main}}\label{E:proof-main}
\begin{proof}
Fix an odd $N\geq3$ and put $n=N^2$. We verify the assertions in the order of construction, categorical realization and numerical invariants.

\emph{Construction and finiteness properties.} Theorem~\ref{E:foundations} supplies the strongly rational VOA $V_{n+2}$ and its order-two simple current $J$. Lemmas~\ref{E:inverse} and \ref{E:Wzero} construct the nonnegatively graded shifted lattice VOA $\mathcal W$ with two simple modules and identify its inverse ribbon semion. The product current $Q=J\otimes M$ has integral lowest weight $3(n-1)/4$, trivial self-braiding and trivial twist. Proposition~\ref{E:existence} consequently constructs
\[
\mathscr{V}_N=(V_{n+2}\otimes\mathcal W)\oplus(J\otimes M)
\]
as a simple, regular, $C_2$-cofinite VOA, proves uniqueness fixing the first summand and establishes its nonnegative integral grading. Equation~\eqref{E:extensionfields} specifies its state-field map, and Section~\ref{E:lattice} realizes these fields inside the deformed rank-eight lattice VOA. Proposition~\ref{E:selfdual} proves self-contragredience by an invariant form on the extension. Proposition~\ref{E:nonCFTmodular} then establishes modularity for its full grading-restricted module category.

\emph{The center and its spherical structure.} Section~\ref{A:cat:condensation} constructs $\Fcat_N$ by central condensation and proves the ribbon equivalence $\D_n\simeq Z(\Fcat_N)$. Sections~\ref{A:cat:HIfusion} and \ref{A:cat:associator} prove its Haagerup--Izumi fusion rules and calculate its pointed associator. In particular,
\[
d(\rho)=\frac{d(R)}N=\frac{N-\sqrt{N^2+4}}2,
\qquad \Omega_N=\omega_1.
\]
The invariant $\prod_{j=0}^{N-1}\omega_1(1,j,1)=e^{2\pi i/N}$ has order $N$, proving the exact order of $[\Omega_N]$ without a primality restriction. Proposition~\ref{E:category} constructs the braided tensor equivalence $\Phi:\D_n\to\Rep(\mathscr{V}_N)$ by local induction. Its tensor compatibility is \eqref{E:inductiontensor}, and \eqref{E:twistpreserved} proves twist preservation. Composing $\Phi^{-1}$ with the preceding center equivalence gives
\[
\Rep(\mathscr{V}_N)\simeq\D_n\simeq Z(\Fcat_N)
\]
as ribbon categories. This identifies the constructed center, with its associators and spherical structure.

\emph{Modules and numerical invariants.} The simplicity and exhaustion argument in Proposition~\ref{E:category} gives one module $\mathsf M_T$ for each simple $T\in\D_n$. Corollary~\ref{B:rankthm} therefore gives exactly $n+3=N^2+3$ irreducible modules; their explicit characters are \eqref{E:vacchar}--\eqref{E:Bchar}. The unchanged extension conformal vector has charge $(1-6n)-11=-6n-10$. Since $Q$ begins in weight at least six, the weight-zero space is $V_0\otimes\mathcal W_0$, of dimension $1\cdot56$. Finally every summand of an induced module has lowest weight at least $-n/4-3/4$, and the $Y\otimes M$ summand of $\mathsf M_X$ attains this bound. Thus
\[
h_{\min}(\mathscr{V}_N)=-\frac{n+3}{4},\qquad
c_{\mathrm{eff}}(\mathscr{V}_N)=(-6n-10)-24\left(-\frac{n+3}{4}\right)=8.
\]
These are all the claims in \eqref{E:maininvariants} and \eqref{E:HIinvariants}, completing the theorem.
\end{proof}

\section{CFT type, twisting and unitarity}\label{E:scope}
\subsection{The grading limitation}
The dimension $\dim(\mathscr{V}_N)_0=56$ shows that $\mathscr{V}_N$ is not of CFT type. Thus the term \emph{strongly rational}, as defined in Section~\ref{sec:inputs}, does not apply, even though $\mathscr{V}_N$ is regular, $C_2$-cofinite and self-contragredient. Nevertheless its representation category is modular by Proposition~\ref{E:nonCFTmodular}, using McRae's theorem rather than a theorem requiring CFT type. The obstruction here is specific and visible in the auxiliary VOA. Any nonzero self-contragredient lattice shift with $2u\in K$ places both $\one$ and $e^{2u}$ in degree zero. Adding a term $D^2e^{2u}$ leaves $L(0)$ unchanged. Tensoring this construction with a CFT-type holomorphic VOA also leaves its degree-zero multiplicity unchanged.

A sufficient replacement would be a fixed strongly rational CFT-type VOA $P$ with exactly two simple modules $P,R$ and ribbon data
\begin{equation}\label{E:CFTproblem}
R^2\cong P,\qquad d(R)=-1,\qquad\theta_R=i,\qquad c_{R,R}=-i.
\end{equation}
Then $(V_{N^2+2}\otimes P)\oplus(J\otimes R)$ would have the same center as its ribbon category and would be of CFT type for all sufficiently large odd $N$, since $3N^2/4+h_R>0$ eventually. No existence assertion for such a $P$ is made here. Nor does this sufficient auxiliary problem exclude other constructions of CFT-type realizations.

\subsection{The twisted target}
The invariant
\[
\prod_{b=0}^{N-1}\Omega_N(1,b,1)=e^{2\pi i/N}
\]
shows that the pointed associator is nontrivial for every parameter under consideration. The cancellation in Section~\ref{E:extension} preserves $\Fcat_N$ and its center; it does not remove this cohomology class or change $d(\rho)$ to its positive Frobenius--Perron value.

There is also a distinction directly in the center. Section~\ref{sec:modular} gives an object of twist $e^{2\pi i/N^2}$, of exact order $N^2$. For the canonical untwisted Haagerup--Izumi data the twist orders divide $N$ or $N^2+4$ \cite[Section~3.2]{B:EG11}; none has order $N^2$, because $\gcd(N,N^2+4)=1$ for odd $N$. A cyclotomic Galois automorphism preserves the exact order of a root of unity, so Galois conjugation cannot erase this distinction.

\subsection{The nonunitary ribbon structure}
The category $\D_n$ has negative-dimensional objects $Q_y$, with
\[
d(Q_y)=\frac{n-\sqrt{n(n+4)}}2<0.
\]
Hence its specified ribbon structure cannot be the representation category of a unitary VOA. This is separate from the elementary fact that the constructed $\mathscr{V}_N$ has negative central charge.

The signs cannot be changed by a different pivotal structure on the same underlying fusion category. Corollary~\ref{B:rankthm} gives
\[
X^2=\bigoplus_{T\in\Irr(\D_n)}T,
\]
including both $\one$ and $X$. If a tensor automorphism of the identity has scalar $\lambda_T$ on $T$, then every displayed summand satisfies $\lambda_T=\lambda_X^2$, while the $X$ summand gives $\lambda_X=\lambda_X^2$. Thus $\lambda_X=1$ and every $\lambda_T=1$. Pivotal structures, when they exist, differ by such automorphisms; the existing one is unique.

Finally the Haagerup--Izumi fusion rules themselves allow no nontrivial group grading of $\Fcat_N$. The summand $\rho$ of $\rho^2$ forces $\deg\rho=e$; the other summands force $\deg\alpha_a=e$. Consequently a cocycle twist of a grading of the entire category does not provide an automatic inverse to \eqref{E:HIinvariants}. A realization of the canonical untwisted unitary series remains a different reconstruction problem.

\appendix
\section{The level-nine singular-vector calculation}\label{app:null}
This appendix gives the coefficient data used in Section~\ref{A:app:c2}. The singular vector is normalized by the coefficient of $L_{-1}^9W$; the recurrence follows by applying the Ward operators in PBW order.

\subsection{The singular vector}
Use PBW order $L_{-\lambda}=L_{-\lambda_1}\cdots L_{-\lambda_\ell}$ with $\lambda_1\geq\cdots\geq\lambda_\ell$. In
\[
s_{3,3}=\sum_{\lambda\vdash9}c_\lambda(n)L_{-\lambda}W,\qquad
c_\lambda(n)=\sum_{a=0}^6c_{\lambda,a}n^a,
\]
all coefficients are as follows. A repeated part is written as a power; for example $(2,1^7)$ has eight parts.
\Needspace{8\baselineskip}
{\small\setlength{\tabcolsep}{6pt}
\begin{longtable}{@{}lrrrrrrr@{}}
\toprule
$\lambda$&$c_0$&$c_1$&$c_2$&$c_3$&$c_4$&$c_5$&$c_6$\\\midrule
\endfirsthead
\toprule
$\lambda$&$c_0$&$c_1$&$c_2$&$c_3$&$c_4$&$c_5$&$c_6$\\\midrule
\endhead
$(9)$ & 0 & -2016 & 12544 & 0 & -1792 & -2048 & 0 \\
$(8,1)$ & 432 & -7936 & 2464 & -3136 & 2432 & -256 & 0 \\
$(7,2)$ & 0 & 1664 & -8192 & 4544 & 1280 & -256 & 0 \\
$(7,1^{2})$ & -384 & 1008 & 2832 & -3552 & 384 & 0 & 0 \\
$(6,3)$ & 0 & -3392 & 18432 & -2816 & -5632 & -1024 & 0 \\
$(6,2,1)$ & 64 & -2048 & -6272 & 3072 & 1536 & 0 & 0 \\
$(6,1^{3})$ & 480 & -1280 & 1920 & -640 & 0 & 0 & 0 \\
$(5,4)$ & 0 & 416 & -5312 & 2432 & 128 & -256 & 0 \\
$(5,3,1)$ & 808 & -7664 & 2824 & 3696 & 96 & -192 & 0 \\
$(5,2^{2})$ & 0 & 1536 & -1536 & 960 & 384 & 0 & 0 \\
$(5,2,1^{2})$ & -1328 & 4256 & -3600 & -960 & 192 & 0 & 0 \\
$(5,1^{4})$ & 288 & -440 & 472 & -48 & 0 & 0 & 0 \\
$(4^{2},1)$ & -624 & 7968 & -3648 & -192 & 384 & 0 & 0 \\
$(4,3,2)$ & 0 & -2560 & 11392 & -288 & -2304 & -384 & 0 \\
$(4,3,1^{2})$ & 312 & -624 & -3096 & 96 & 288 & 0 & 0 \\
$(4,2^{2},1)$ & 256 & -9088 & -1152 & 1728 & 384 & 0 & 0 \\
$(4,2,1^{3})$ & 864 & 3520 & -224 & -384 & 0 & 0 & 0 \\
$(4,1^{5})$ & -216 & -180 & 72 & 0 & 0 & 0 & 0 \\
$(3^{3})$ & 0 & 1152 & -5088 & -424 & 1904 & 672 & 64 \\
$(3^{2},2,1)$ & -256 & 3328 & 2208 & -2832 & -1536 & -192 & 0 \\
$(3^{2},1^{3})$ & -324 & 176 & 260 & 192 & 48 & 0 & 0 \\
$(3,2^{3})$ & 0 & 0 & -1024 & -256 & 0 & 0 & 0 \\
$(3,2^{2},1^{2})$ & 384 & -1344 & 1728 & 1248 & 192 & 0 & 0 \\
$(3,2,1^{4})$ & -144 & -240 & -336 & -96 & 0 & 0 & 0 \\
$(3,1^{6})$ & 36 & 6 & 12 & 0 & 0 & 0 & 0 \\
$(2^{4},1)$ & 0 & 1024 & 256 & 0 & 0 & 0 & 0 \\
$(2^{3},1^{3})$ & -256 & -640 & -384 & -64 & 0 & 0 & 0 \\
$(2^{2},1^{5})$ & 144 & 192 & 48 & 0 & 0 & 0 & 0 \\
$(2,1^{7})$ & -24 & -12 & 0 & 0 & 0 & 0 & 0 \\
$(1^{9})$ & 1 & 0 & 0 & 0 & 0 & 0 & 0 \\
\bottomrule
\end{longtable}}
The coefficient of $L_{-1}^9W$ is one. The verification uses only
\[
[L_a,L_b]=(a-b)L_{a+b}+\frac{1-6n}{12}(a^3-a)\delta_{a+b,0},\qquad L_0W=2nW,
\]
and $L_mW=0$ for $m>0$. Negative-mode reordering uses
$L_{-a}L_{-b}=L_{-b}L_{-a}+(b-a)L_{-(a+b)}$ for $a<b$.
Thus every coefficient of each $L_ms_{3,3}$ reduces to the zero polynomial, and $s_{3,3}$ vanishes in $L(1-6n,2n)$.

\subsection{The recurrence polynomials}
For complete coefficient data, write
\[
R_j(k,n)=2^{9-j}\sum_{m=0}^{2j+1}a_{j,m}(n)k^m,\qquad j=1,2,3.
\]
\Needspace{8\baselineskip}
\noindent\textbf{Polynomial $R_1$ (scale $256$).}
\begin{center}\small\begin{tabular}{@{}r l@{}}\toprule
$m$ & $a_{1,m}(n)$\\\midrule
3 & $2 n^{3} + 12 n^{2} + 20 n + 8$ \\
2 & $3 n^{3} + 2 n^{2} - 34 n + 12$ \\
1 & $- 13 n^{3} - 58 n^{2} - 22 n + 4$ \\
0 & $4 n^{4} + 16 n^{3}$ \\
\bottomrule\end{tabular}\end{center}
\Needspace{10\baselineskip}
\noindent\textbf{Polynomial $R_2$ (scale $128$).}
\begin{center}\small\begin{tabular}{@{}r l@{}}\toprule
$m$ & $a_{2,m}(n)$\\\midrule
5 & $- 12 n^{2} - 48 n - 36$ \\
4 & $24 n^{3} + 84 n^{2} - 84$ \\
3 & $12 n^{4} + 138 n^{3} + 363 n^{2} + 96 n - 51$ \\
2 & $- 24 n^{5} - 150 n^{4} - 273 n^{3} - 141 n^{2} + 98 n - 5$ \\
1 & $- 48 n^{5} - 270 n^{4} - 255 n^{3} + 390 n^{2} - 20 n + 2$ \\
0 & $24 n^{5} + 156 n^{4} + 192 n^{3} - 96 n^{2}$ \\
\bottomrule\end{tabular}\end{center}
\Needspace{12\baselineskip}
\noindent\textbf{Polynomial $R_3$ (scale $64$).}
\begin{center}\small\begin{tabular}{@{}r l@{}}\toprule
$m$ & $a_{3,m}(n)$\\\midrule
7 & $24 n + 48$ \\
6 & $- 96 n^{2} - 132 n + 216$ \\
5 & $48 n^{3} - 360 n^{2} - 1038 n + 396$ \\
4 & $192 n^{4} + 1248 n^{3} + 1144 n^{2} - 1835 n + 354$ \\
3 & $- 168 n^{5} - 312 n^{4} + 1614 n^{3} + 2280 n^{2} - 1670 n + 152$ \\
2 & $- 96 n^{6} - 1260 n^{5} - 3228 n^{4} - 45 n^{3} + 2029 n^{2} - 545 n + 23$ \\
1 & $96 n^{7} + 624 n^{6} + 268 n^{5} - 1904 n^{4} + 1923 n^{3} - 73 n^{2} - 1$ \\
0 & $144 n^{7} + 984 n^{6} + 924 n^{5} - 2694 n^{4} + 552 n^{3} - 24 n^{2}$ \\
\bottomrule\end{tabular}\end{center}
The end polynomials are $R_0(k,n)=-512n(n+4)k$ and the factorized $R_4$ in Lemma~\ref{A:c2:lem:div}. These formulas can also be recovered directly from the singular-vector table: apply each Ward operator from the rightmost PBW factor to the leftmost, acting on a monomial $z^{e+s}$ by
\[
D_1z^{e+s}=(e+s)z^{e+s-1},\qquad
D_rz^{e+s}=\delta_{r,2}z^{e+s}+(-1)^{r-2}\bigl((r-1)2n-e-s\bigr)z^{e+s-r}.
\]
Sum the thirty contributions and use $R_j(k,n)=P_{4-j}(-2k,n)$.

\section{An alternative uniform finiteness and Zhu argument}\label{H:uniform}
This appendix gives an alternative proof of uniform rationality. It derives a sharper $C_2$ bound and proves Zhu semisimplicity by reducing the singular-vector equation to a tensor product of hypergeometric equations. Proposition~\ref{H:gaugeclaim} gives the required differential-gauge identities. The coefficient recursions are specified below, and Appendix~\ref{H:gauge} gives the complete scalar row.

Throughout put $h=2n$, $c=1-3h$, $t+t^{-1}=n+2$, $t>1$, and $m=n+4$. Let $W$ be the primary generator of weight $h$, and let $s_\lambda(h)=c_\lambda(h/2)$ be the thirty singular-vector coefficients in Appendix~\ref{app:null}. Write $\nu=W_{2h-1}W\in\C^\times$ and let $C\ne0$ be the coefficient of $W$ in $W_{h-1}W$. Their nonvanishing is supplied by the lattice construction. The notation $x=[T]$, $y=[W]$ will be used in either the $C_2$ algebra or Zhu's algebra, as specified.

\subsection{Complete recursive coefficient definitions}
For a partition $\lambda=(r,\mu)$ of a number at most nine and an integer $s\geq0$, define polynomials in an auxiliary variable $z$ by
\begin{align*}
F_{\varnothing,s}(a,h;z)&=1,\\
F_{(r,\mu),s}(a,h;z)&=\delta_{r,2}zF_{\mu,s}(a,h;z)
+\bigl(2h+8-2a-s+(h-1)(r-1)\bigr)F_{\mu,s+r}(a,h;z),\\
Q_p(a,h)&=[z^p]\sum_{\lambda\vdash9}s_\lambda(h)F_{\lambda,0}(a,h;z).
\end{align*}
Only $0\leq p\leq4$ occur. For Zhu's algebra define
\begin{align*}
G_{\varnothing,s}(d,h,x;z)&=z^{9-s},\\
G_{(r,\mu),s}(d,h,x;z)&=(-1)^r\bigl((r-1)x+d-9+s+|\mu|\bigr)G_{\mu,s}(d,h,x;z)\\
&\quad+\bigl(2h+8-d-s+(h-1)(r-1)\bigr)G_{\mu,s+r}(d,h,x;z),\\
K_p(d,h,x)&=[z^p]\sum_{\lambda\vdash9}s_\lambda(h)G_{\lambda,0}(d,h,x;z).
\end{align*}
Only $0\leq p\leq9$ occur, $K_0(d,h,x)=Q_0(d/2,h)$, and $\deg_xK_p\leq\lfloor p/2\rfloor$. These fixed recursions involve the thirty partitions of nine.

Commuting $W_a$ through the singular vector before passing to the relevant quotient gives
\begin{align}
\sum_{p=0}^4Q_p(a,h)x^pR_{a-p}&=0,
&R_a&=[W_{2h-1-2a}W]\in V/C_2(V),\label{H:C2rec}\\
\sum_{p=0}^9K_p(d,h,x)Z_{d-p}&=0,
&Z_d&=[W_{2h-1-d}W]\in A(V).\label{H:Zhurec}
\end{align}
Negative indices mean zero. The identities follow from
\[
[W_a,L_{-r}]=\bigl(a+(h-1)(r-1)\bigr)W_{a-r},\qquad
[L_{-r}v]=(-1)^r\bigl((r-1)x+\wt v\bigr)[v]
\]
in Zhu's algebra. In the $C_2$ quotient a leftmost $L_{-r}$ survives only for $r=2$. Thus the calculation does not assume that positive modes preserve $C_2(V)$.

\subsection{The sharper uniform \texorpdfstring{$C_2$}{C2} bound}
Here $R_a=0$ for $a<0$ or $a>h$, $R_0=\nu$, and $R_h=y^2$. The recursive coefficients give
\[
Q_4(a,h)=-128h(h+8)(a-h-4).
\]
Define $b_0=1$, $b_j=0$ for $j<0$, and
\begin{equation}\label{H:backward}
b_j=-\frac{\sum_{p=0}^3Q_p(h+4-j,h)b_{j+p-4}}{128h(h+8)j}
\quad(j\geq1).
\end{equation}
Put $B_j=2^j j!b_j/h^2$ for $j\geq1$. The first values are
\begin{align*}
B_1&=1,\\
B_2&=h^2-h+3,\\
B_3&=h^4-3h^3+3h^2+32h+57,\\
B_4&=h^6-6h^5-3h^4+200h^3-279h^2+6750h+3417.
\end{align*}
For $j\geq5$,
\[
B_j=\sum_{q=1}^4A_q(h,j)B_{j-q},\qquad
A_q(h,j)=-\frac{2^q(j-1)\cdots(j-q+1)Q_{4-q}(h+4-j,h)}{128h(h+8)}.
\]
For $h\equiv2\pmod4$, the following are lower bounds on $v_2(A_q)$; starred bounds are equalities:
\[
\begin{array}{c|rrrr}
j\pmod4&q=1&q=2&q=3&q=4\\\hline
0&1&1&2&2^*\\
1&2&4&4&5\\
2&2&2&4&4\\
3&1&2&2^*&5
\end{array}
\]
They are finite polynomial congruences: substitute $h=4v+2$, $j=4r+s$ in the numerator of $A_q$. Its denominator is $2^9(2v+1)(2v+5)$. After division by $2^{9+e_{s,q}}$, all coefficients are integral; for $(s,q)=(0,4),(3,3)$ the numerator divided by $2^{11}$ is identically one modulo two. These statements are completely specified by the preceding recursion and the table of $s_\lambda$.

Induction in blocks of four now gives
\begin{equation}\label{H:valuation}
v_2(B_j)\geq2\left\lfloor\frac{j-1}{4}\right\rfloor,
\end{equation}
with equality for $j\equiv0,3\pmod4$. The first four values are odd. In a block beginning at $4r+1$, the first two new values have valuation at least $2r$. For $j=4r+3$, the $q=3$ term has valuation exactly $2r$, and the other terms have larger valuation; for $j=4r+4$, the $q=4$ term is the unique term of valuation $2r$. This proves \eqref{H:valuation}, and in particular
\[
v_2(B_{h+1})=n-1,\qquad b_{h+1}\ne0.
\]

\begin{theorem}[Sharper uniform finiteness bound]\label{H:C2bound}
For every odd $n\geq3$,
\[
x^{2n+3}=x^{n+3}y=0,\qquad \dim V_{n+2}/C_2(V_{n+2})\leq3n+6.
\]
\end{theorem}
\begin{proof}
Put $T_j=x^jR_{h-j}$ for $0\leq j\leq h+1$, and $T_j=0$ at negative indices. Multiply \eqref{H:C2rec} at $a=h+4-j$ by $x^{j-1}$. The resulting relation, compared with \eqref{H:backward}, proves inductively
\[
x^3T_j=b_jx^3R_h.
\]
Since $T_{h+1}=0$ and $b_{h+1}\ne0$, one gets $x^3R_h=0$. At $j=h$ and $j=n$ this gives $\nu x^{h+3}=0$ and $x^{n+3}R_n=0$. The relation $R_n=\alpha x^n+Cy$ with $C\ne0$ implies both asserted vanishings. Also $y^2\in\C x^{2n}+\C x^ny$. The monomials
\[
1,x,\ldots,x^{2n+2},\qquad y,xy,\ldots,x^{n+2}y
\]
span the quotient and number $3n+6$.
\end{proof}

\subsection{A universal Zhu polynomial}
Write $\mathcal B=A(V)$, a finite-dimensional commutative algebra generated by $x,y$. Set $Z_d$ as in \eqref{H:Zhurec}; then $Z_0=\nu\ne0$ and $Z_h=Cy+b(x)$ with $C\ne0$. There is a polynomial $P(w)\in\mathcal B[w]$ with
\begin{equation}\label{H:P}
\sum_{d\geq0}Z_dz^d=P\!\left(\frac{z^2}{1+z}\right),\qquad
\deg P\leq h,\qquad P(0)=\nu\one.
\end{equation}
Indeed higher Zhu circle relations make $(1+z)^h\sum_dZ_dz^d$ a polynomial of degree at most $2h$. Skew symmetry and $[D^ra]=(-1)^r(\wt a)_r[a]$ make the series invariant under $z\mapsto-z/(1+z)$. The invariant coordinate is $w=z^2/(1+z)$, and the pole bounds at $-1,\infty$ give a polynomial in $w$ of degree at most $h$.

Let $D_z=z\partial_z$, $D_w=w\partial_w$. Equation \eqref{H:Zhurec} specifies the scalar operator
\begin{equation}\label{H:Lz}
\mathcal L_z=\sum_{p=0}^9z^pK_p(D_z+p,h,x).
\end{equation}
Its symmetric version $\mathcal L_w$ is specified uniquely by
\begin{equation}\label{H:coordinate}
\mathcal L_z(P(w(z)))=(z+2)(z+1)^4(\mathcal L_wP)(w(z)).
\end{equation}
Thus $\mathcal L_wP=0$. Equations \eqref{H:Zhurec} and \eqref{H:coordinate} give an explicit finite definition of this order-nine operator.

\subsection{The rank-three systems and the rational gauge}
For $H=2t-1$ or $2/t-1$, consider
\begin{equation}\label{H:rankthree}
\left[D_w(D_w-H/2)(D_w-(3H+1)/2)
+\frac w4(D_w-H)\bigl(D_w^2-(H-1)D_w-2(H+1)x\bigr)\right]F=0.
\end{equation}
For the first system, choose $\lambda^2=n+4x$ and $q=\sqrt t\lambda$. In the coordinate $-w/4$, its hypergeometric parameters are
\begin{equation}\label{H:parameters}
\alpha=(1-2t,1-t+q,1-t-q),\qquad
\beta=(1,3/2-t,2-3t).
\end{equation}
The second system replaces $t$ by $t^{-1}$. Write $\mathcal H_t(x)$ and $\mathcal H_{1/t}(x)$ for their rank-three monodromy representations. At $w=-4$ the exceptional local exponent is always $1/2$, besides $0,1$.

\begin{proposition}[Exact gauge identities]\label{H:gaugeclaim}
The scalar row in Appendix~\ref{H:gauge} defines a rational gauge $\Gamma(w,x)$ from the tensor product of \eqref{H:rankthree} to $\mathcal L_w$, in derivative coordinates. Set
\[
A=4x-3t+2,\quad B=4tx+2t-3,\quad C_t=4tx-3t^2+6t-3,
\]
and
\[
\mathfrak d=(2t-1)^2(2/t-1)^2x(4x-3t+2)(4x-3/t+2)(8x-3h).
\]
Then $\Gamma_{\rm num}=\mathfrak d\Gamma$ is polynomial in $x$, and
\begin{equation}\label{H:det}
\det\Gamma=\frac{K(t)}{w^9(w+4)^9x^4A^2B^2C_t},
\end{equation}
where
\[
K(t)=\frac{5038848(t-2)^3(t-1)^{21}(t+1)^5(t+2)^2
(2t-3)^9(2t-1)^3(2t+1)^2(3t-2)^9}{t^{24}}.
\]
The residual ranks of $\Gamma_{\rm num}$ at $x=0,3n/4$ are four and one. For origin Frobenius solutions with exponents
\[
\rho_0(t)=0,\qquad\rho_1(t)=t-1/2,\qquad\rho_2(t)=3t-1,
\]
the scalar row shifts the leading exponent of the product indexed by $(a,b)$ by exactly $-ab$. Its normalized leading coefficient for $(0,0)$ is one; for $(2,2)$ it is a nonzero function of $t$ divided by $\mathfrak d$, so has a simple pole at $x=3n/4$.
\end{proposition}

\begin{proof}
The scalar row and the rank-three equations reduce the assertion to polynomial identities. Apply the scalar row to the nine tensors $(D_w^pF_1)(D_w^qF_2)$, $0\leq p,q\leq2$, and eliminate all higher derivatives with \eqref{H:rankthree}. Successive Euler derivatives through order eight give $\Gamma$. Substituting these rows into the operator specified by \eqref{H:Lz} and \eqref{H:coordinate}, and clearing denominators, gives zero in each of the nine tensor components. Thus the differential identity holds over $\Q(t,x,w)$. At the origin, the rank-three recurrences give the Frobenius coefficients. Substitution into the scalar row cancels every power preceding the stated shift and gives the nine nonzero leading coefficients, with the stated normalizations. The trace difference of the Euler first-order systems is $-18(w+2)/(w+4)$, giving the factor $w^{-9}(w+4)^{-9}$ in the determinant. The product of Frobenius leading coefficients and the ratio of the two Vandermonde determinants give its remaining factor. For the exceptional ranks, a nonzero $4\times4$ minor at $x=0,w=1$, using derivative rows $0,1,2,3$ and tensor columns $(1,1),(1,2),(2,1),(2,2)$, is
\[
-\frac{820125(t-2)^6(t-1)^{11}(t+1)(2t-3)^8(2t-1)^6(3t-2)^8}{t^{20}}.
\]
At $x=3n/4,w=1$, the row-zero, column-$(0,0)$ entry is
\[
\frac{441(t-2)^2(t-1)^4(2t-3)(2t-1)^2(3t-4)(3t-2)(4t-3)}{4t^6}.
\]
At zero the five columns with $p=0$ or $q=0$ vanish; at $3n/4$ the $2\times2$ minors vanish identically. The displayed nonzero certificates give the asserted ranks. All $t$-factors displayed above are nonzero, since $t>4$ at the parameters considered here.
\end{proof}

\subsection{The regular inverse and rational exceptional weights}
We apply Proposition~\ref{H:gaugeclaim}. Over the discrete valuation ring in $x$ at zero, $\det\Gamma_{\rm num}$ has valuation five and residual rank four, so its elementary-divisor valuations are $0,0,0,0,1,1,1,1,1$. Since $\mathfrak d$ has a simple zero, $\Gamma^{-1}$ is regular and has residual rank five. At $x=3n/4$, the corresponding determinant valuation is eight and the residual rank is one; the inverse is again regular, now of rank eight.

These statements are valid on the solution spaces at the origin as well. The scalar exponents are independent of $x$, with the only integral resonance class consisting of $0,n,3n$. For generic $x$ the gauge gives a log-free basis; the finitely many resonance obstructions are polynomial identities in $x$ with denominators independent of $x$, so vanish identically. Choosing the free resonant coefficients to be zero gives a basis holomorphic in $x$. In this basis
\[
(\Gamma^{-1}P)_{00}=P(0)
\]
at both exceptional parameters and after base change to every local Artin algebra. At $x=3n/4$, the inverse kernel is the scalar Frobenius line beginning with $w^{3n}$: the $(2,2)$ coefficient has a simple pole there, and the inverse has a one-dimensional kernel.

For rational $\xi$, either $\mathcal H_t(\xi)$ or $\mathcal H_{1/t}(\xi)$ can be reducible only at
\begin{equation}\label{H:exceptional}
\xi\in\mathscr E=\{0,1,3n/4,3n/4+4\}.
\end{equation}
To check this, apply the hypergeometric criterion $\alpha_i-\beta_j\notin\Z$ for irreducibility \cite[Propositions~3.2--3.3]{H:BH}. The differences involving $1-2t$ are irrational. An integral remaining difference forces $q=at+b$ with $a\in\{0,\pm1,\pm2\}$ and rational $b$. Using $q^2=t(n+4\xi)$ and $t^2=(n+2)t-1$, comparison in the basis $1,t$ gives $b^2=a^2$. The case $a=0$ comes from a nonzero half-integer $b$ and is impossible. The other cases give $n+4\xi\in\{n,m,4n,4m\}$, exactly \eqref{H:exceptional}.

\subsection{Reduced local Zhu factors away from the exceptions}
Let $\mathcal B_\xi$ be a local Artin factor of $A(V)$, with residue $x=\xi\in\Q\setminus\mathscr E$. Its polynomial \eqref{H:P}, under the inverse gauge, is a single-valued horizontal tensor, equivalently an intertwiner
\[
T:\mathcal H_t(x)^*\longrightarrow\mathcal H_{1/t}(x)
\]
over $\mathcal B_\xi$. Monodromy matrices depend holomorphically on $x$ at a nonsingular base point, so substitution of their finite Taylor series makes this statement meaningful over the Artin algebra. The determinant of the gauge is a unit there. The reduction of $T$ is nonzero since $P(0)=\nu\one$; irreducibility and Schur's lemma make it an isomorphism. It remains an isomorphism over $\mathcal B_\xi$.

Comparing the traces of monodromy at infinity gives the identity
\begin{align}
G(x)&=\cos(2\pi\sqrt{t(n+4x)})
-\cos(2\pi\sqrt{(n+4x)/t})\notag\\
&=-2\sin(\pi\sqrt m\sqrt{n+4x})
\sin(\pi\sqrt n\sqrt{n+4x})=0\quad\text{in }\mathcal B_\xi.\label{H:G}
\end{align}
This is entire in $x$, independent of the square-root choices. Its zeros are
\[
\xi=\frac{a^2}{4n}-\frac n4\quad\text{or}\quad
\xi=\frac{b^2}{4m}-\frac n4\qquad(a,b\in\Z).
\]
Every zero is simple as a zero in $x$. At $n+4x=0$, the product is a nonzero constant times $n+4x$; elsewhere both sine factors cannot vanish because $\sqrt{nm}$ is irrational. Thus $G(x)=(x-\xi)u(x)$ with $u(\xi)\ne0$, and \eqref{H:G} gives $x=\xi\one$ in the local algebra.

With $x$ scalar, Schur's lemma makes the space of single-valued tensor solutions one-dimensional. The condition $P(0)=\nu\one$ fixes its multiple, so every coefficient $Z_d$ is scalar. Since $Z_h=Cy+b(x)$ with $C\ne0$, $y$ is scalar as well, and $\mathcal B_\xi=\C$. The same normalization determines the same $y$ at every factor over $\xi$; hence at most one maximal ideal lies over any such $\xi$.

\subsection{The two upper exceptional parameters}
At $x=1$, the numerator parameters are $(1-2t,-2t,2)$. Hypergeometric contiguity \cite[Propositions~2.5 and~2.7]{H:BH} gives
\[
0\longrightarrow R_t\longrightarrow\mathcal H_t(1)
\longrightarrow\one\longrightarrow0,
\]
with $R_t$ irreducible of rank two. There is no trivial subrepresentation: its scalar solution would be a rational function regular at $w=0,-4$ and with only integral infinity exponent $-2$, hence a polynomial of negative degree. Consequently any map from $\mathcal H_t(1)^*$ to $\mathcal H_{1/t}(1)$ kills the trivial source submodule and has image in $R_{1/t}$. Its $(0,0)$ origin coefficient is zero. The only integral-exponent products are $(0,0),(1,1),(2,2)$, and the other two become series starting at $w^n,w^{3n}$ under the gauge. Thus $P(0)=0$, contradicting \eqref{H:P}.

At $x=3n/4+4$, the numerator parameters are $(1-2t,t+3,-3t-1)$, giving
\[
0\longrightarrow\chi_t\longrightarrow\mathcal H_t(x)
\longrightarrow R_t\longrightarrow0.
\]
Here $\chi_t$ has origin exponent $3t-1$. This sequence is nonsplit: a quotient $\chi_t$ would give an adjoint solution with character $\chi_t^{-1}$, equal to $w^{-3t+1}$ times a polynomial of degree $-3$. Indeed its infinity exponent is $-3t-2$, and its exponents at $-4$ are $0,1,1/2$. An intertwiner from the dual of the first system to the second must therefore factor through the rank-one source quotient into the rank-one target submodule. Its only origin component is $(2,2)$, which the scalar row sends to a series starting at $w^{3n}$. It cannot give a nonzero polynomial of degree at most $h=2n$. Neither upper exceptional value is a Zhu eigenvalue.

\subsection{The two lower exceptional parameters}
Let $\mathcal B_\xi$ now be a local factor at $\xi=0$ or $3n/4$. The regular inverse gauge produces an intertwiner $T$ whose reduction is nonzero, because its $(0,0)$ coefficient is $P(0)=\nu\one$.

At zero the numerator parameters are $(1-2t,2-2t,0)$, giving the nonsplit sequence
\[
0\longrightarrow\one\longrightarrow\mathcal H_t(0)
\longrightarrow R_t\longrightarrow0.
\]
A trivial quotient would yield a rational adjoint solution regular at the two finite singularities and with infinity exponent $-1$, which is impossible. A nonzero intertwiner from the dual of the first system to the second has rank one and factors through their trivial constituents. At infinity the eigenvalue one is isolated from the other two eigenvalues, both $e^{-4\pi it}$ in the first representation. Its spectral projector varies holomorphically in $x$. The component of $T$ between the matching one-dimensional eigenlines is a unit over the Artin algebra. Equality of the eigenvalues therefore forces
\[
e^{2\pi i s(x)}-1=0,\qquad
s(x)=-n+\sqrt n\sqrt{n+4x}.
\]
Since $s(0)=0$ and $s'(0)=2$, this gives $x=0$ in $\mathcal B_0$.

At $x=3n/4$, the parameters are $(1-2t,t-1,3-3t)$, and
\[
0\longrightarrow R_t\longrightarrow\mathcal H_t(3n/4)
\longrightarrow\chi_t\longrightarrow0
\]
is nonsplit. A $\chi_t$ submodule would give $w^{3t-1}$ times a polynomial with infinity exponent $3t-3$, hence degree $-2$. A nonzero intertwiner between the dual and the other system factors as an isomorphism of the irreducible rank-two constituents. All infinity eigenvalues are distinct, and one of the matched moving eigenline components is a unit. The exponent difference is again $s(x)$; now $s(3n/4)=n$ and $s'(3n/4)=1$. Thus $x=(3n/4)\one$ in its local Artin factor.

To finish, at $x=0$ the Zhu quadratic is $y(y-w_0)=0$ with $w_0\ne0$, and both roots are realized by $V$ and $X$. The whole summand is therefore $\C\oplus\C$. At $x=3n/4$, the inverse-gauge kernel begins at $w^{3n}$ and meets no polynomial of degree at most $2n$. Such polynomials inject into the one-dimensional homomorphism space of the rank-two constituents. Normalization fixes the polynomial uniquely, and $Z_h=Cy+b(x)$ makes $y$ scalar. The current $J$ realizes this factor, which is exactly $\C$.

\begin{theorem}[Uniform Zhu semisimplicity]\label{H:Zhuconclusion}
For every odd $n\geq3$, the Zhu algebra $A(V_{n+2})$ is semisimple. The zero-weight summand is $\C\oplus\C$, each nonzero weight fibre has dimension at most one, and weight one is absent. Consequently $V_{n+2}$ is strongly rational.
\end{theorem}
\begin{proof}
Theorem~\ref{H:C2bound} gives uniform $C_2$-cofiniteness. Rationality of conformal weights under $C_2$-cofiniteness \cite[Corollary~5.10]{H:Miyamoto} makes every residue of $x$ rational. At a rational value outside \eqref{H:exceptional}, the other factors in \eqref{H:det} do not vanish: $A,B$ could vanish only at irrational $x$, and $C_t$ vanishes at $3n/4$. The cases just treated therefore exhaust all local Artin factors, proving semisimplicity. The simple, self-contragredient CFT-type VOA then satisfies McRae's rationality criterion \cite[Main Theorem~3]{A:McR}. This proves strong rationality.
\end{proof}

The argument uses the uniform identities above and is independent of the character classification and categorical reconstruction in the main text.

\section{Complete coefficients of the uniform Zhu gauge}\label{H:gauge}
This appendix gives all thirty-nine nonzero scalar-row coefficients of the rational gauge. Together with the thirty singular-vector coefficients in Appendix~\ref{app:null} and the fixed recursions in Appendix~\ref{H:uniform}, it specifies every rational function in Proposition~\ref{H:gaugeclaim}. The symbols $E_j$ in this appendix denote auxiliary polynomials.

For solutions $F_1,F_2$ of the rank-three equations \eqref{H:rankthree}, the scalar row is
\begin{equation}\label{H:scalarrow}
\frac1{\mathfrak d}\sum_{r=-4}^{0}\sum_{p,q=0}^{2}
R_{r,p,q}(t,x)w^r(D_w^pF_1)(D_w^qF_2),
\end{equation}
where $\mathfrak d$ is as in Proposition~\ref{H:gaugeclaim}. Put
\[
\begin{gathered}
\ell_1=t-2,\quad\ell_2=t-1,\quad\ell_3=2t-3,\quad\ell_4=2t-1,\\
\ell_5=3t-4,\quad\ell_6=3t-2,\quad\ell_7=4t-3.
\end{gathered}
\]
All omitted coefficients vanish. The table is grouped by the Laurent power $r$.
{\small\renewcommand{\arraystretch}{1.7}
\begin{longtable}{@{}rrr l@{}}
\toprule
$r$&$p$&$q$&$R_{r,p,q}(t,x)$\\\midrule
\endfirsthead
\toprule
$r$&$p$&$q$&$R_{r,p,q}(t,x)$\\\midrule
\endhead
$-4$&1&1&$-36\ell_1^2\ell_3^2\ell_4^2\ell_5\ell_6^2\ell_7/t^5$\\
$-4$&1&2&$-72\ell_1\ell_3^2\ell_4^2\ell_5\ell_6^2\ell_7/t^4$\\
$-4$&2&1&$72\ell_1^2\ell_3^2\ell_4\ell_5\ell_6^2\ell_7/t^5$\\
$-4$&2&2&$144\ell_1\ell_3^2\ell_4\ell_5\ell_6^2\ell_7/t^4$\\\addlinespace
$-3$&0&1&$-72\ell_1^2\ell_3^2\ell_4^2\ell_5\ell_6\ell_7x/t^4$\\
$-3$&0&2&$-144\ell_1\ell_3^2\ell_4^2\ell_5\ell_6\ell_7x/t^3$\\
$-3$&1&0&$72\ell_1^2\ell_3\ell_4^2\ell_5\ell_6^2\ell_7x/t^5$\\
$-3$&1&1&$36E_1\ell_1\ell_2^2\ell_3\ell_4\ell_6/t^5$\\
$-3$&1&2&$-18E_2\ell_1\ell_3\ell_4\ell_6/t^4$\\
$-3$&2&0&$-144\ell_1^2\ell_3\ell_4\ell_5\ell_6^2\ell_7x/t^5$\\
$-3$&2&1&$-18E_3\ell_1\ell_3\ell_4\ell_6/t^5$\\
$-3$&2&2&$72E_4\ell_1\ell_3\ell_4\ell_6/t^4$\\\addlinespace
$-2$&0&0&$144\ell_1^2\ell_3\ell_4^2\ell_5\ell_6\ell_7x^2/t^4$\\
$-2$&0&1&$12E_5\ell_1\ell_3\ell_4^2\ell_6x/t^4$\\
$-2$&0&2&$-12E_6\ell_1\ell_3\ell_4^2\ell_6x/t^3$\\
$-2$&1&0&$12E_7\ell_1^2\ell_3\ell_4\ell_6x/t^5$\\
$-2$&1&1&$6E_8\ell_1\ell_3\ell_4\ell_6/t^5$\\
$-2$&1&2&$3E_9\ell_1\ell_3\ell_4\ell_6/t^4$\\
$-2$&2&0&$-12E_{10}\ell_1^2\ell_3\ell_4\ell_6x/t^5$\\
$-2$&2&1&$-3E_{11}\ell_1\ell_3\ell_4\ell_6/t^5$\\
$-2$&2&2&$3E_{12}\ell_1\ell_3\ell_4\ell_6/t^4$\\\addlinespace
$-1$&0&0&$48E_{13}\ell_1^2\ell_3\ell_4^2\ell_6x^2/t^4$\\
$-1$&0&1&$6E_{14}\ell_1\ell_3\ell_4^2\ell_6x/t^4$\\
$-1$&0&2&$-6E_{15}\ell_1\ell_3\ell_4^2\ell_6x/t^3$\\
$-1$&1&0&$-6E_{16}\ell_1^2\ell_3\ell_4\ell_6x/t^5$\\
$-1$&1&1&$-6E_{17}\ell_1\ell_3\ell_4\ell_6/t^5$\\
$-1$&1&2&$3E_{18}\ell_1\ell_3\ell_4\ell_6/(2t^4)$\\
$-1$&2&0&$-6E_{19}\ell_1^2\ell_3\ell_4\ell_6x/t^5$\\
$-1$&2&1&$-3E_{20}\ell_1\ell_3\ell_4\ell_6/(2t^5)$\\
$-1$&2&2&$3E_{21}\ell_1\ell_3\ell_4\ell_6/(2t^4)$\\\addlinespace
0&0&0&$6E_{22}\ell_1^2\ell_3\ell_4^2\ell_6x/t^4$\\
0&0&1&$3E_{23}\ell_1\ell_3\ell_4^2\ell_6x/t^3$\\
0&0&2&$-3\ell_1\ell_2\ell_3\ell_4^2\ell_6x(16x-1)/t^2$\\
0&1&0&$-3E_{24}\ell_1^2\ell_3\ell_4\ell_6x/t^4$\\
0&1&1&$-3E_{25}\ell_1\ell_3\ell_4\ell_6x/t^4$\\
0&1&2&$3E_{26}\ell_1\ell_3\ell_4\ell_6x/t^3$\\
0&2&0&$-3\ell_1^2\ell_2\ell_3\ell_4\ell_6x(16x-1)/t^4$\\
0&2&1&$-3E_{27}\ell_1\ell_3\ell_4\ell_6x/t^4$\\
0&2&2&$3E_{13}\ell_1\ell_3\ell_4\ell_6x/t^3$\\
\bottomrule
\end{longtable}}

The twenty-seven auxiliary polynomials are as follows. Each display is an exact expression over $\Z[t,x]$.
\begin{align*}
E_1={}&80t(t^2-t+1)x-13(t-2)(2t-3)(2t-1)(3t-2),\\
E_2={}&16t(t-1)(4t^2+2t-7)x\\
&+(2t-3)(2t-1)(3t-2)(40t^2-87t+46),\\
E_3={}&16t(t-1)(7t^2-2t-4)x\\
&-(t-2)(2t-3)(3t-2)(46t^2-87t+40),\\
E_4={}&40t(t-1)^2x+(2t-3)(3t-2)(17t^2-35t+17),\\
E_5={}&4t^2(16t-17)x-(t-2)(2t-3)(54t^2-97t+42),\\
E_6={}&8t(8t-9)x+(2t-3)(72t^2-143t+66),\\
E_7={}&4t(17t-16)x+(2t-1)(3t-2)(42t^2-97t+54).
\end{align*}
\begin{align*}
E_8={}&8t^2(8t^2-11t+8)x^2\\
&+2t(108t^4-308t^3+395t^2-308t+108)x\\
&-27(t-2)(t-1)^2(2t-3)(2t-1)(3t-2),\\
E_9={}&32t^2(4t-7)x^2-8t(60t^3-110t^2+29t+18)x\\
&-3(t-1)(2t-3)(2t-1)(3t-2)(22t-27),\\
E_{10}={}&8t(9t-8)x+(3t-2)(66t^2-143t+72),\\
E_{11}={}&32t^2(7t-4)x^2+8t(18t^3+29t^2-110t+60)x\\
&-3(t-2)(t-1)(2t-3)(3t-2)(27t-22),\\
E_{12}={}&64t^2x^2+16t(39t^2-79t+39)x\\
&+3(2t-3)(3t-2)(32t^2-65t+32),\\
E_{13}={}&4tx+9t^2-19t+9,\\
E_{14}={}&32t^2x^2+8t(t^2+5t-7)x\\
&-(t-2)(t-1)(36t^2-79t+36).
\end{align*}
\begin{align*}
E_{15}={}&4t(16t-17)x+36t^3-124t^2+134t-45,\\
E_{16}={}&32t^2x^2-8t(7t^2-5t-1)x\\
&-(t-1)(2t-1)(36t^2-79t+36),\\
E_{17}={}&32t^3x^3-8t^2(11t^2-19t+11)x^2\\
&+t(12t^4-95t^3+170t^2-95t+12)x\\
&+3(t-2)(t-1)^2(2t-3)(2t-1)(3t-2),\\
E_{18}={}&16t^2(8t-13)x^2-4t(84t^3-196t^2+128t-21)x\\
&-3(t-1)(2t-3)(2t-1)(3t-2)(4t-5),\\
E_{19}={}&4t(17t-16)x+45t^3-134t^2+124t-36,\\
E_{20}={}&16t^2(13t-8)x^2-4t(21t^3-128t^2+196t-84)x\\
&-3(t-2)(t-1)(2t-3)(3t-2)(5t-4),\\
E_{21}={}&64t^2x^2+8t(33t^2-68t+33)x
+15(t-1)^2(2t-3)(3t-2).
\end{align*}
\begin{align*}
E_{22}={}&8tx^2+2(t^2-3t+1)x+(t-1)^2,\\
E_{23}={}&16tx^2-4(3t^2-9t+7)x+(t-1)(3t-4),\\
E_{24}={}&16tx^2-4(7t^2-9t+3)x+(t-1)(4t-3),\\
E_{25}={}&16t^2x^2-4t(9t^2-16t+9)x\\
&+18t^4-78t^3+121t^2-78t+18,\\
E_{26}={}&4t(2t-3)x-18t^3+46t^2-36t+9,\\
E_{27}={}&4t(3t-2)x-9t^3+36t^2-46t+18.
\end{align*}

For completeness, the change of differential coordinates in \eqref{H:coordinate} can be recovered without solving for an unspecified operator. Set $N_{0,0}=1$, use zero for out-of-range indices, and define
\[
N_{r,j}(z)=z\bigl((z+1)N'_{r-1,j}(z)-(r-1)N_{r-1,j}(z)\bigr)
+(z+2)N_{r-1,j-1}(z).
\]
Then
\[
D_z^r=\sum_{j=0}^rN_{r,j}(z)(z+1)^{-r}D_w^j.
\]
Substitute this expression into \eqref{H:Lz}, divide by $(z+2)(z+1)^4$, and replace $z^2/(z+1)$ by $w$. The result is the symmetric operator $\mathcal L_w$. Applying successive Euler derivatives to \eqref{H:scalarrow}, reducing each rank-three factor by \eqref{H:rankthree}, then gives the full gauge matrix. Thus the data needed for the differential identity, Frobenius leading terms, determinant and exceptional-rank assertions are all contained in the manuscript.

\begin{thebibliography}{99}\footnotesize\raggedright\setlength{\itemsep}{0pt}
\bibitem{A:Evans} D. E. Evans, \emph{A vertex operator algebra for the S-fold realisation of a non-unitary Haagerup--Izumi modular fusion category}, arXiv:2610.26341 (2026). \url{https://arxiv.org/abs/2610.26341}.
\bibitem{I:Riesen3} A. Riesen, \emph{The Haagerup VOA at $c=8$}, arXiv:2610.13933 (2026). \url{https://arxiv.org/abs/2610.13933}.
\bibitem{I:Huang3} T.-C. Huang, \emph{Haagerup conformal field theory at central charge eight}, arXiv:2610.14441 (2026). \url{https://arxiv.org/abs/2610.14441}.
\bibitem{I:Huang5} T.-C. Huang, \emph{An order-five Haagerup--Izumi vertex operator algebra at central charge eight}, Hexagon 2610.00213v2 (2026). \url{https://hexagonmath.org/pdf/2610.00213v2}.
\bibitem{I:GK} D. Gang and D. Kim, \emph{Generalized non-unitary Haagerup--Izumi modular data from 3D S-fold SCFTs}, J. High Energy Phys. 03 (2023), 185. \url{https://arxiv.org/abs/2211.13561}.
\bibitem{I:GKL} D. Gang, D. Kim and S. Lee, \emph{A non-unitary bulk-boundary correspondence: Non-unitary Haagerup RCFTs from S-fold SCFTs}, SciPost Phys. 17 (2024), 064. \url{https://arxiv.org/abs/2310.14877}.
\bibitem{I:Flohr} M. Flohr, \emph{$\mathcal W$-algebras, new rational models and completeness of the $c=1$ classification}, Commun. Math. Phys. 157 (1993), 179--212. \url{https://arxiv.org/abs/hep-th/9207019}.
\bibitem{A:Li} H.-S. Li, \emph{Symmetric invariant bilinear forms on vertex operator algebras}, J. Pure Appl. Algebra 96 (1994), 279--297. \url{https://doi.org/10.1016/0022-4049(94)90104-X}.
\bibitem{A:CMSY} T. Creutzig, R. McRae, K. Shimizu, and H. Yadav, \emph{Commutative algebras in Grothendieck--Verdier categories, rigidity, and vertex operator algebras}, arXiv:2409.14618v2, Theorem~2.26. \url{https://arxiv.org/abs/2409.14618}.
\bibitem{A:KZ} L. Kong and H. Zheng, \emph{Semisimple and separable algebras in multi-fusion categories}, arXiv:1706.06904, Theorems~6.3 and~6.10. \url{https://arxiv.org/abs/1706.06904}.
\bibitem{A:KO} A. Kirillov Jr. and V. Ostrik, \emph{On a $q$-analogue of the McKay correspondence and the ADE classification of $\widehat{\mathfrak{sl}}_2$ conformal field theories}, Adv. Math. 171 (2002), 183--227; Theorems 4.1 and 4.5. \url{https://arxiv.org/abs/math/0101219}.
\bibitem{A:DMNO} A. Davydov, M. M\"uger, D. Nikshych, and V. Ostrik, \emph{The Witt group of non-degenerate braided fusion categories}, arXiv:1009.2117, Sections~3--5. \url{https://arxiv.org/abs/1009.2117}.
\bibitem{A:EGNO} P. Etingof, S. Gelaki, D. Nikshych, and V. Ostrik, \emph{Tensor Categories}, AMS, 2015, Proposition~9.2.2. \url{https://math.mit.edu/~etingof/egnobookfinal.pdf}.
\bibitem{A:FF} B. Feigin and D. Fuchs, \emph{Representations of the Virasoro algebra}, in \emph{Representations of Lie Groups and Related Topics}, Adv. Stud. Contemp. Math. 7 (1990), 465--554.
\bibitem{A:Fr} E. Frenkel, \emph{Determinant formulas for the free field representations of the Virasoro and Kac--Moody algebras}, Phys. Lett. B 286 (1992), 71--77. \url{https://doi.org/10.1016/0370-2693(92)90160-6}.
\bibitem{A:BNW} B. D. Boe, D. K. Nakano and E. Wiesner, \emph{Category $\mathcal O$ for the Virasoro algebra: cohomology and Koszulity}, Pacific J. Math. 234 (2008), 1--21, Sections 2--3. \url{https://arxiv.org/abs/math/0701644}.
\bibitem{A:CJORY} T. Creutzig, C. Jiang, F. Orosz Hunziker, D. Ridout and J. Yang, \emph{Tensor categories arising from the Virasoro algebra}, Adv. Math. 380 (2021), 107516; Theorem 4.1.7 (Miyamoto), Theorem 5.1.2. \url{https://arxiv.org/abs/2002.03180}.
\bibitem{A:CKM} T. Creutzig, S. Kanade, and R. McRae, \emph{Gluing vertex algebras}, arXiv:1906.00119. \url{https://arxiv.org/abs/1906.00119}.
\bibitem{A:Orb} R. McRae, \emph{On the tensor structure of modules for compact orbifold vertex operator algebras}, arXiv:1810.00747, Corollary~4.8. \url{https://arxiv.org/abs/1810.00747}.
\bibitem{A:AA} T. Abe and Y. Arike, \emph{Intertwining operators and fusion rules for vertex operator algebras arising from symplectic fermions}, J. Algebra 373 (2013), 39--64; Section 2.3 and Theorem 2.12 (Frenkel--Zhu--Li injection). \url{https://arxiv.org/abs/1108.1823}.
\bibitem{A:CKL} T. Creutzig, S. Kanade, and A. R. Linshaw, \emph{Simple current extensions beyond semi-simplicity}, arXiv:1511.08754, Lemma~2.13. \url{https://arxiv.org/abs/1511.08754}.
\bibitem{A:McR} R. McRae, \emph{On rationality for $C_2$-cofinite vertex operator algebras}, Camb. J. Math. 14 (2026), no.~1, 1--115; Main Theorem~3 of arXiv:2108.01898v3. \url{https://arxiv.org/abs/2108.01898}.
\bibitem{A:EP} P. Etingof and D. Penneys, \emph{Rigidity of non-negligible objects of moderate growth in braided categories}, arXiv:2412.17681, Theorem~1.1. \url{https://arxiv.org/abs/2412.17681}.
\bibitem{A:Ind} T. Creutzig, R. McRae, and J. Yang, \emph{Direct limit completions of vertex tensor categories}, Commun. Contemp. Math.~24 (2022), 2150033, Theorems~7.1 and~7.5. \url{https://arxiv.org/abs/2006.09711}.
\bibitem{A:EG} D. E. Evans and T. Gannon, \emph{Near-group fusion categories and their doubles}, Adv. Math. 255 (2014), 586--640; Proposition 7(b), Section 4.4. \url{https://arxiv.org/abs/1208.1500}.
\bibitem{B:Ostrik} V. Ostrik, \emph{On formal codegrees of fusion categories}, Math. Res. Lett. 16 (2009), 895--901; Section 2.2. \url{https://arxiv.org/abs/0810.3242}.
\bibitem{B:ENO} P. Etingof, D. Nikshych and V. Ostrik, \emph{On fusion categories}, Ann. Math. 162 (2005), 581--642; Theorem 10.1. \url{https://arxiv.org/abs/math/0203060}.
\bibitem{B:BVG} V. M. Buchstaber, A. A. Gaifullin and A. P. Veselov, \emph{Classification of involutive commutative two-valued groups}, Russian Math. Surveys 77 (2022), 651--727; Theorem 1.10. \url{https://arxiv.org/abs/2302.05909}.
\bibitem{B:GSY} T. Gannon, A. Schopieray and H. Yadav, \emph{On Haagerup--Izumi fusion categories}, 2026; Section 5.2. \url{https://arxiv.org/abs/2609.25185}.
\bibitem{B:HKL} Y.-Z. Huang, A. Kirillov Jr. and J. Lepowsky, \emph{Braided tensor categories and extensions of vertex operator algebras}, Commun. Math. Phys. 337 (2015), 1143--1159; Theorem 3.2 and Remark 3.3. \url{https://arxiv.org/abs/1406.3420}.
\bibitem{B:Huang} Y.-Z. Huang, \emph{Vertex operator algebras and the Verlinde conjecture}, Commun. Contemp. Math. 10 (2008), 103--154. \url{https://arxiv.org/abs/math/0406291}.
\bibitem{B:BS} J. H. Bruinier and M. Schwagenscheidt, \emph{Algebraic formulas for the coefficients of mock theta functions and Weyl vectors of Borcherds products}, J. Algebra 478 (2017), 38--57; Section 2.1 and Lemma 2.1. \url{https://arxiv.org/abs/1607.02701}.
\bibitem{B:EG11} D. E. Evans and T. Gannon, \emph{The exoticness and realisability of twisted Haagerup--Izumi modular data}, Commun. Math. Phys. 307 (2011), 463--512; Section 3.2, equation (3.3). \url{https://arxiv.org/abs/1006.1326}.

\bibitem{E:GR} T. Gannon and A. Riesen, \emph{Orbifolds of pointed vertex operator algebras I}, Adv. Math. 482 (2025), 110546; Theorem~2.3. \url{https://doi.org/10.1016/j.aim.2025.110546}.
\bibitem{E:DM} C. Dong and G. Mason, \emph{Shifted vertex operator algebras}, Math. Proc. Cambridge Philos. Soc. 141 (2006), 67--80; Theorems~3.1, 4.1 and 4.4. \url{https://arxiv.org/abs/math/0411526}.
\bibitem{E:Carnahan} S. Carnahan, \emph{Building vertex algebras from parts}, Comm. Math. Phys. 373 (2020), no.~1, 1--43; Theorem~3.2.12 in arXiv:1408.5215v4. \url{https://arxiv.org/abs/1408.5215}.
\bibitem{E:CKMext} T. Creutzig, S. Kanade and R. McRae, \emph{Tensor categories for vertex operator superalgebra extensions}, Mem. Amer. Math. Soc. 295 (2024), no.~1472; theorem numbering follows arXiv version~2. \url{https://arxiv.org/abs/1705.05017v2}.
\bibitem{H:BH} F. Beukers and G. Heckman, \emph{Monodromy for the hypergeometric function ${}_nF_{n-1}$}, Invent. Math. 95 (1989), 325--354; Proposition~2.7 and Propositions~3.2--3.3.
\bibitem{H:Miyamoto} M. Miyamoto, \emph{Modular invariance of vertex operator algebras satisfying $C_2$-cofiniteness}, Duke Math. J. 122 (2004), 51--91; Corollary~5.10 of arXiv:math/0209101v2. \url{https://arxiv.org/abs/math/0209101}.
\end{thebibliography}
\end{document}
